\chapter{Aplicaciones de la integral}\label{aplicaciones}

En este capítulo desarrollamos las siguientes aplicaciones: área entre curvas, volumen de sólidos con función de área transversal conocida, volumen de sólidos de revolución con secciones transversales discos y arandelas, volumen de sólidos de revolución por el método de capas cilíndricas, longitud de arco de una curva plana, área de una superficie de revolución, trabajo de problemas de resortes, problemas de bombeo de un líquido en un tanque, presión y fuerza de un fluido, fuerza ejercida por un líquido sobre una lámina vertical, centros de masa y momentos de inercia en barras delgadas y placas planas.

\section{Áreas entre curvas}

Sea $R$ la región acotada por las curvas $y=f(x)$, $y=g(x)$, $x=a$ \ y \
$x=b$, \ donde $g(x)\leq f(x)$ \ en \ $[a, b]$, con $f$ y $g$ continuas en $[a,b]$; entonces, el área de la
región $R$ está dada por
\[
  A=\int_{a}^{b}\left( f(x)-g(x)\right) dx
\]
La representación figura para aplicar la anterior fórmula es la siguiente:

\begin{figure}[H]
  \centering
  \caption{Área entre curvas, funciones de $x$:$f(x)\geq g(x)$}
  \includegraphics[width=0.30\textwidth]{Gráficas/GrafAreaEntreCurvFx.pdf}
  \propia
\end{figure}

Si la región $R$\ está acotada por las curvas $x=f(y)$, $x=g(y)$, $y=c$
\ y \ $y=d$, \ donde $f(y)\leq g(y)$ \ en \ $c\leq y\leq d.$ \ \ El área de
la región $R$ está dada por
\[
  A=\int_{c}^{d}\left( g(y)-f(y)\right) dy
\]
La representación figura para aplicar la anterior fórmula es la siguiente:

\begin{figure}[H]
  \centering
  \caption{Área entre curvas, funciones de $y$:$g(y)\geq f(y)$}
  \includegraphics[width=0.38\textwidth]{Gráficas/GrafAreaEntreCurvFy.pdf}
  \propia
\end{figure}

Si la región $R$ comprendida entre las curvas $y=f(x)$, $y=g(x)$, $x=a$
\ y \ $x=b$, \ donde $g(x)\leq f(x)$ \ para algunos valores de $x$, pero $%
f(x)\leq g(x)$ para otros valores de $x,$ entonces el área está dada por:
\begin{flalign*}
  A & =\int_{a}^{b}\left\vert f(x)-g(x)\right\vert dx, & & \\ % Alinea la integral a la izquierda
  & \text{donde } |f(x)-g(x)| =
  \begin{cases}
    f(x)-g(x), & \text{si } g(x)\le f(x)\\[6pt]
    g(x)-f(x), & \text{si } f(x)\le g(x)
  \end{cases} &&
\end{flalign*}

\begin{figure}[H]
  \centering
  \caption{Área entre curvas: cruce de curvas}
  \includegraphics[width=0.35\textwidth]{Gráficas/GrafAreaEntreCurvFG.pdf}
  \propia
\end{figure}

%\begin{center}---------
%\centerline{ \graf{Área entre curvas: cruce de curvas}}
%\includegraphics[scale=0.3]{GrafAreaEntreCurvFG.pdf}
%\textbf{Fuente:} elaboración propia
%\end{center}
%
\begin{ejemplo}
  Dibujar la región limitada por las curvas dadas por $y=x^{2}+5x$ y $y=3x+8$ y hallar el valor de su área.
\end{ejemplo}

\emph{Solución.} La figura \ref{fig:area_curvas} muestra la región acotada por las curvas dadas:
%\clearpage

\begin{figure}[H]
  \centering
  \caption{Área entre parábola y recta}
  \label{fig:area_curvas}
  \includegraphics[width=0.30\textwidth]{Gráficas/GrafAreaEntreCurv1.pdf}
  \propia
\end{figure}

%\begin{center}---------
%\centerline{ \graf{Área entre parábola y recta}}
%\includegraphics[scale=0.35]{GrafAreaEntreCurv1.pdf}
%\textbf{Fuente:} elaboración propia
%\end{center}

Para plantear los límites de integración encontremos sus puntos de corte: Igualamos sus ecuaciones:
{\small
  \begin{align*}
    x^{2}+5x&=3x+8 \\
    x^{2}+2x-8&=0\\
    (x+4)(x-2)&=0,\\
  \end{align*}
}

es decir, $x=-4$ y $x=2$. Con esos valores de $x$, los valores de $y$ respectivos son $y=-4$ y $y=14$. Esta información puede verificarse en la figura anterior.

El área de la región entonces, está dada por
{\footnotesize
  \begin{align*}
    a(R)&=\int_{-4}^{2}(3x+8-(x^{2}+5x))dx=\int_{-4}^{2}(8-2x-x^{2})dx \\
    & =\left(8x-x^{2}-\frac{1}{3}x^{3}\right)\Bigg\vert_{-4}^{2}
    =16-4-\frac{8}{3}-\left(-32-16+\frac{64}{3}\right)=60-24=36
  \end{align*}
}

Si nuestras unidades de medida, por ejemplo, son los centímetros, entonces el área de la región es $36$ cm$^2$.

\begin{ejemplo}
  Halle el área de la región acotada por las curvas $y=9-x^{2}$ \ y $y=x^{2}-9.$
\end{ejemplo}

\emph{Solución.} En la figura se muestra la región a medir su área:

%\clearpage

\begin{figure}[H]
  \centering
  \caption{Área entre parábolas}
  \includegraphics[width=0.30\textwidth]{Gráficas/GrafAreaEntreCurv2.pdf}
  \propia
\end{figure}

%\begin{center}---------
%\centerline{ \graf{Área entre parábolas}}
%\includegraphics[scale=0.35]{GrafAreaEntreCurv2.pdf}
%\textbf{Fuente:} elaboración propia
%\end{center}

Hallamos inicialmente los puntos de intersección de las curvas:
\begin{align*}
  x^{2}-9&=9-x^{2} \\
  2x^{2}&=18 \\
  x&=\pm 3
\end{align*}
Con ello, la integral es
{\small
  \begin{align*}
    A&=\int_{-3}^{3}\left( 9-x^{2}-\left( x^{2}-9\right) \right)
    dx\\
    &=\int_{-3}^{3}\left( 18-2x^{2}\right) dx \\
    &=\left( 18x-\frac{2}{3}x^{3}\right)\Bigg\vert_{-3}^{3}\\
    &=18\left( 3\right) -\left( \frac{2}{3}\right) 3^{3}-\left( 18\left(
    -3\right) -\left( \frac{2}{3}\right) \left( -3\right) ^{3}\right)\\
    &=54-18-(-54+18)\\
    &=54-18+54-18
    =72
  \end{align*}
}
Así, el área de la región es $72$ unidades cuadradas.

\begin{ejemplo}
  Halle el área de la región acotada por $x-y-5=0$\ y\ $x-y^{2}+1=0$.
\end{ejemplo}
\emph{Solución.} En este caso, una de las curvas está en términos de $y$, es decir, $y$ no está de grado uno, pero $x$ sí. Encontremos los puntos de corte, pero en términos de $y$:
\begin{align*}
  y+5&=-1+y^{2}\\
  y^{2}-y-6&=0\\
  \left(y+2\right) \left( y-3\right) &=0
\end{align*}
De allí tenemos que $ y=-2$ y $y=3$. La figura \ref{fig:area_recta_parabola} muestra la forma de la región:

\begin{figure}[H]
  \centering
  \caption{Área entre parábola y recta}
  \label{fig:area_recta_parabola}
  \includegraphics[width=0.30\textwidth]{Gráficas/GrafAreaEntreCurv3.pdf}
  \propia
\end{figure}

%\begin{center}---------
%\centerline{ \graf{Área entre parábola y recta}}
%\includegraphics[scale=0.35]{GrafAreaEntreCurv3.pdf}
%\textbf{Fuente:} elaboración propia
%\end{center}

Dado que la variable $y$ nos da la información de la figura, es conveniente integrar con respecto a $y$:

\resizebox{\textwidth}{!}{$
  \begin{aligned}
    A&=\int_{-2}^{3}\left( y+5-\left( y^{2}-1\right) \right)
    dy=\int_{-2}^{3}\left( 6+y-y^{2}\right) dy\\
    &=\left( 6y+\frac{1}{2}y^{2}-\frac{1}{3}y^{3}\right)\Bigg\vert_{-2}^{3} \\
    &=6y+\left( \frac{1}{2}\right) y^{2}-\frac{1}{2}y^{3}=6\left( 3\right)
    +\left( \frac{1}{2}\right) \left( 3\right) ^{2}-\left( \frac{1}{3}\right)
    \left( 3\right) ^{3}-\left( 6\left(-2\right) +\left( \frac{1}{2}\right)
    \left( -2\right) ^{2}-\left( \frac{1}{3}\right) \left( -2\right) ^{3}\right) \\
    &=18+\left( \frac{9}{2}\right) -9-\left( -12+2+\left( \frac{8}{3}\right)
    \right) =9+\frac{9}{2}+10-\frac{8}{3} \\
    &=\frac{114+27-16}{6}\\
    &=\frac{125}{6} \text{  unidades cuadradas}.
  \end{aligned}
$}

\begin{ejemplo}
  Halle el área de la región limitada por $x=y^{2}+y+6$ y $x=9-y^{2}.$
\end{ejemplo}
Precisemos un poco la figura con los puntos de corte:
\[
  y^{2}+y+6=9-y^{2}\rightarrow 2y^{2}+y-3=0
\]
En esta parte podemos factorizar la expresión o usar la fórmula cuadrática:
\[
  y=\frac{-1\pm \sqrt{1-4\left( 2\right) \left( -3\right)
  }}{4}=\frac{-1\pm 5}{4}\Rightarrow y=-\frac{3}{2},\text{ y } y=1
\]
%\clearpage

\begin{figure}[H]
  \centering
  \caption{Región cuya área se va a medir}
  \includegraphics[width=0.30\textwidth]{Gráficas/GrafAreaEntreCurv4.pdf}
  \propia
\end{figure}

%\begin{center}---------
%\centerline{ \graf{Región cuya área se va a medir}}
%\includegraphics[scale=0.35]{GrafAreaEntreCurv4.pdf}
%\textbf{Fuente:} elaboración propia
%\end{center}

Con ello calculamos el área:
{\small
  \begin{align*}
    A&=\int_{-\frac{3}{2}}^{1}\left( 9-y^{2}-\left( y^{2}+y+6\right)
    \right) dy =\int_{-\frac{3}{2}}^{1}\left( 3-y-2y^{2}\right) dy \\
    &=\left( 3y-\frac{1}{2}y^{2}-\frac{2}{3}y^{3}\right)\Bigg\vert _{-\frac{3}{2}}^{1}\\
    &=\left( 3-\frac{1}{2}-\frac{2}{3}\right) -\left( -\frac{9}{2}-\frac{9}{8}+%
    \frac{2}{3}\times \frac{27}{8}\right) \\
    &=\frac{125}{24}\approx5.2083\text{ unidades cuadradas}%
  \end{align*}%
}

\begin{ejemplo}
  Halle el área de la región limitada por las curvas dadas por $y=\cfrac{10x^{2}}{x^{2}+1}$ y $y=x^{2}.$
\end{ejemplo}

\emph{Solución.} Al igualar las expresiones, tenemos
\begin{align*}
  x^{2}&=\frac{10x^{2}}{x^{2}+1}\rightarrow x^{4}+x^{2}=10x^{2}\rightarrow
  x^{2}\left( x^{2}-9\right) =0\\
  &\rightarrow x^{2}\left( x-3\right) \left(
  x+3\right) =0\rightarrow x=0,\ x=-3,\ x=3
\end{align*}
Las figuras, entonces, tienen $3$ puntos de corte\footnote{ En esta ecuación suele presentarse que, al solucionarla, se tiende a cancelar los $x^2$. Debe tener en cuenta que puede cancelarse (con división) alguna expresión, siempre y cuando haya certeza de que esa expresión no es cero, pues, en caso contrario, estaríamos dividiendo entre cero. Para este caso $x^2$, puede ser cero; por tanto, no es correcto cancelarlos cuando solucionamos la ecuación.} dados por $x=0,\ x=-3$ y $x=3$.

Tomando algunos valores adicionales a aquellos, puede identificarse que las figuras tienen la siguiente forma:

\begin{figure}[H]
  \centering
  \caption{Región cuya área se va a medir}
  \includegraphics[width=0.35\textwidth]{Gráficas/GrafAreaEntreCurv5.pdf}
  \propia
\end{figure}

%\begin{center}----------
%\centerline{ \graf{Región cuya área se va a medir}}
%\includegraphics[scale=0.4]{GrafAreaEntreCurv5.pdf}
%\textbf{Fuente:} elaboración propia
%\end{center}

Y por otro lado,
\begin{align*}
  A&=\int_{-3}^{3}\left( \frac{10x^{2}}{x^{2}+1}-x^{2}\right) dx\ \text{(note que el integrando es par)}\\
  &=2\int_{0}^{3}\left( \frac{10x^{2}}{x^{2}+1}-x^{2}\right) dx\ \text{(simplificamos con división)}\\
  &=2\int_{0}^{3}\left( 10-\frac{10}{x^{2}+1}-x^{2}\right) dx\\
  &=2\left( 10x-\frac{1}{3}x^{3}-10\tan^{-1} x\right)\Bigg\vert_{0}^{3}\\
  &=2\left(30-9-10\tan^{-1} 3\right) =42-20\tan^{-1} 3\approx 17.019\text{ unidades cuadradas}%
\end{align*}

\begin{ejemplo}
  Halle el área de la región limitada por las figuras de $y=\sqrt{\frac{1}{2}x+4}$ y $y=\sqrt{100-x}$ para $y\geq 0$.
\end{ejemplo}
\emph{Solución.} Encontremos los puntos de corte de las curvas dadas: si $\sqrt{\frac{1}{2}x+4}=\sqrt{100-x}$, entonces $\cfrac{1}{2}x+4=100-x$, es decir, $\cfrac{3}{2}x=96$ o también $x=64$, con lo cual $y=\sqrt{100-64}=6.$ La figura de la región, se muestra a continuación:

\begin{figure}[H]
  \centering
  \caption{Región cuya área se va a medir}
  \includegraphics[width=0.38\textwidth]{Gráficas/GrafAreaEntreCurv6.pdf}
  \propia
\end{figure}

%\begin{center}---------
%\centerline{ \graf{Región cuya área se va a medir}}
%\includegraphics[scale=0.35]{GrafAreaEntreCurv6.pdf}
%\textbf{Fuente:} elaboración propia
%\end{center}

Por la forma de la región a la que se le quiere medir el área, es decir, que está encerrada por una curva a la derecha y una a la izquierda, conviene integrar con respecto a $y$. Escribamos, entonces, las expresiones en términos de $y$:
\[
  \sqrt{\frac{1}{2}x+4}=y\ \rightarrow \frac{1}{2}x+4=y^{2}\ \rightarrow
  x=2y^{2}-8
\]
y
\[
  \sqrt{100-x}=y\ \rightarrow 100-x=y^{2}\ \rightarrow x=100-y^{2}
\]
Con ello, el área es

\begin{align*}
  A&=\int_{0}^{6}\left( 100-y^{2}-\left( 2y^{2}-8\right) \right)
  dy =\int_{0}^{6}\left( 108-3y^{2}\right) dy\\
  &=\left( 108y-y^{3}\right)\Big\vert_{0}^{6}=108\times 6-6^{3}=432\text{ unidades cuadradas}%
\end{align*}%

\begin{ejemplo}
  Halle el área de la región limitada por $y=\cos (x),y=\sen (2x)$, $x=\frac{\pi }{2}$ y $x=\pi$.
\end{ejemplo}

\emph{Solución.} Hallamos los puntos de corte de las dos curvas. Dado que $\sen(2x)=2\sen x\cos x$, entonces $\sen (2x)=\cos x$ implica que $2\sen x\cos x-\cos x=0$ o, también, $\cos x( 2\sin x-1) =0$, de donde $\cos x=0$ o $\sen x=\frac{1}{2}$, con lo cual $ x=( 2n+1) \frac{\pi }{2}$, $x=\frac{\pi }{6}+2n\pi$ o $x=\frac{5}{6}\pi +2n\pi$ para $n\in \mathbb{Z}$.

Los valores de $x$ en el intervalo $\left[ \frac{\pi }{2},\pi \right] $, intervalo donde se define la región, son: $x=\frac{\pi }{2}$ y $x=\frac{5}{6}\pi $.

Tomemos una tabla de valores para aproximar la figura de la región:

\begin{table}[H]
  \caption{Valores de aproximación}
  \renewcommand{\arraystretch}{1.3}
  \begin{tabular}{|l|r|r|r|r|r|}
    \hline
    $x$ & $\frac{\pi }{2}$ & $\frac{2\pi }{3}$ & $\frac{5\pi }{6}$ & $\frac{11\pi }{12}$ & \multicolumn{1}{l|}{$\pi $} \\
    \hline
    $\cos x$ & $0$ & $-\frac{1}{2}$ & $-\frac{\sqrt{3}}{2}$ & $-0,96$ &
    \multicolumn{1}{l|}{$-1$} \\
    \hline
    $\sen \left( 2x\right) $ & $0$ & $-\frac{\sqrt{3}}{2}$ & $-\frac{\sqrt{3}}{2}
    $ & $-\frac{1}{2}$ & \multicolumn{1}{l|}{$0$}\\
    \hline
  \end{tabular}
  \propia
\end{table}

Abajo la figura:

\begin{figure}[H]
  \centering
  \caption{Área entre seno y coseno}
  \includegraphics[width=0.36\textwidth]{Gráficas/GrafAreaEntreCurv7.pdf}
  \propia
\end{figure}

%\begin{center}----------
%\centerline{ \graf{Área entre seno y coseno}}
%\includegraphics[scale=0.3]{GrafAreaEntreCurv7.pdf}
%\textbf{Fuente:} elaboración propia
%\end{center}

Con ello, el área de la región, usando valor absoluto en la integral, está dada por
\begin{gather*}
  \begin{aligned}
    A&=\int_{\frac{\pi }{2}}^{\pi }\left\vert \cos x-\sen \left( 2x\right)
    \right\vert dx
    =\int_{\frac{\pi }{2}}^{\frac{5}{6}\pi }\left( \cos x-\sen
    \left( 2x\right) \right) dx+\int_{\frac{5}{6}\pi }^{\pi }\left( \sen \left(
    2x\right) -\cos x\right) dx\\
    &=\left( \sen x+\frac{1}{2}\cos \left( 2x\right) \right) _{\frac{\pi }{2}}^{    \frac{5}{6}\pi }+\left( -\frac{1}{2}\cos \left( 2x\right) -\sen x\right)\Bigg\vert_{\frac{5}{6}\pi }^{\pi }\\
    &=\sen \frac{5}{6}\pi +\frac{1}{2}\cos \frac{5}{3}\pi -\left( \sen \frac{\pi}{2}+\frac{1}{2}\cos \pi \right) -\frac{1}{2}\cos \left( 2\pi \right) -\sen \pi+\frac{1}{2}\cos \left( \frac{5}{3}\pi \right) +\sen \frac{5}{6}\pi \\
    &=\frac{1}{2}\text{unidades cuadradas}%
  \end{aligned}
\end{gather*}

\section{Volumen de sólidos con función de área transversal
conocida}

La sección transversal de un sólido $S$ es la región plana
formada por la intersección de $S$ con un plano. Por ejemplo, si tomamos un cono circular recto, sus secciones transversales perpendiculares al eje del cono son círculos {(figura \ref{fig:seccion_cono})}

\begin{figure}[H]
  \centering
  \caption{Sección transversal del cono}
  \label{fig:seccion_cono}
  \includegraphics[width=0.36\textwidth]{Gráficas/SeccTransCono.pdf}
  \propia
\end{figure}

%\begin{center}---------
%\centerline{ \graf{Sección transversal del cono}}
%\includegraphics[scale=0.35]{SeccTransCono.pdf}
%\textbf{Fuente:} elaboración propia
%\end{center}

Y las secciones transversales paralelas a la base cuadrada de la pirámide son cuadrados (figura \ref{seccion_piramide}):

\begin{figure}[H]
  \centering
  \caption{Sección transversal de la pirámide}
  \label{seccion_piramide}
  \includegraphics[width=0.36\textwidth]{Gráficas/SeccTransPiramide.pdf}
  \propia
\end{figure}

%\begin{center}----------
%\centerline{ \graf{Sección transversal de la pirámide}}
%\includegraphics[scale=0.35]{SeccTransPiramide.pdf}
%\textbf{Fuente:} elaboración propia
%\end{center}

Sea $S$ un sólido acotado en el eje $x$ por el intervalo $\left[ a,b\right], $ tal que las secciones transversales perpendiculares al eje $x$ tengan
una función de área conocida A$\left( x\right) ,$ (continua en $\left[
a,b\right] )$; entonces, el volumen del sólido es
{\small
  \[
    V=\int_{a}^{b}A\left( x\right) dx
  \]
}

\begin{figure}[H]
  \centering
  \caption{Área sección transversal}
  \includegraphics[width=0.36\textwidth]{Gráficas/SeccTransDefinic.pdf}
  \propia
\end{figure}

%\begin{center}-------
%\centerline{ \graf{Área sección transversal}}
%\includegraphics[scale=0.35]{SeccTransDefinic.pdf}
%\textbf{Fuente:} elaboración propia
%\end{center}

Para encontrar el volumen usando secciones transversales, podemos seguir los siguientes pasos:

\begin{parrafonumerado}
\item Dibujar el sólido y una sección transversal.
\item Determinar una fórmula para el área de la sección
  transversal $A(x)$, en términos de $x$.
\item Determinar los limites de integración.
\item Integrar $A(x)$.
\end{parrafonumerado}

\begin{ejemplo}
  Encuentre el volumen del sólido cuya base es el triángulo con vértices
  $(0,0)$, $(3,3)$ y $(3,-3)$ y cuyas secciones transversales perpendiculares al eje $x$ son cuadrados con un lado sobre el triángulo.
\end{ejemplo}
%%\begin{multicols}{2}
\emph{Solución.} La base del sólido cuyo volumen quiere calcularse está dada por este triángulo:

\begin{figure}[H]
  \centering
  \caption{Triángulo como base del sólido}
  \includegraphics[width=0.30\textwidth]{Gráficas/GrafArTransvBaseTriang1.pdf}
  \propia
\end{figure}

%\begin{center}--------
%\centerline{ \graf{Triángulo como base del sólido}}
%\includegraphics[scale=0.3]{GrafArTransvBaseTriang1.pdf}
%\textbf{Fuente:} elaboración propia
%\end{center}

El enunciado nos dice que las secciones transversales son cuadrados perpendiculares a $x$ con un lado en el triángulo, es decir, un lado del cuadrado se ve como en la figura:

\begin{figure}[H]
  \centering
  \caption{Lado de un cuadrado}
  \includegraphics[width=0.36\textwidth]{Gráficas/GrafArTransvBaseTriang2.pdf}
  \propia
\end{figure}

%\begin{center}---------
%\centerline{ \graf{Algunos cuadrados}}
%\includegraphics[scale=0.3]{GrafArTransvBaseTriang2.pdf}
%\textbf{Fuente:} elaboración propia
%\end{center}
%\centerline{\graf{Lado del cuadrado}\label{LadoCuad}}
%%\end{multicols}

En la figura \ref{fig:parabola_encerrada}, se ven algunos cuadrados completos seg\'un la información suministrada:

\begin{figure}[H]
  \centering
  \caption{Secciones transversales}
  \label{fig:parabola_encerrada}
  \includegraphics[width=0.56\textwidth]{Gráficas/GrafArTransvBaseTriang3.pdf}
  \propia
\end{figure}

%\begin{center}---------
%\centerline{ \graf{Región encerrada por la parábola}}
%\includegraphics[scale=0.3]{GrafArTransvBaseTriang3.pdf}
%\textbf{Fuente:} elaboración propia
%\end{center}

Ahora bien, el área de un cuadrado de lado $L$ está dada por $A=L^2$. En nuestro caso, la longitud del lado la dan las dos rectas que forman el triángulo, es decir, las rectas $y=x$ y $y=-x$ (véase la figura 38); por tanto, para un $x$ entre $0$ y $3$, se tiene que $L=x-(-x)=2x$, es decir, en términos de $x$, el área está dada por $A(x)=L^{2}=4x^{2}.$ Con ello, el volumen del sólido es
{\small
  \[
    V=\int_{0}^{3}A(x)dx=\int_{0}^{3}4x^{2}dx=\left( \frac{4}{3}x^{3}\right) \Bigg\vert_{0}^{3}=36
  \]
}
El volumen del sólido, entonces, es de $36$ unidades c\'ubicas.

\begin{ejemplo}
  Encuentre el volumen del sólido cuya base es la región limitada por $y=4-x^{2}$ y el eje $x$, y cuyas secciones transversales perpendiculares al eje $x$ son cuadrados con un lado sobre su base.
\end{ejemplo}

\emph{Solución.} La figura de $y=4-x^{2}$ es una parábola que abre hacia abajo con vértice en $(0,4)$. La región encerrada por ella y el eje $x$ se muestra en la figura \ref{fig:region_parabola_2}:

\begin{figure}[H]
  \centering
  \caption{Región encerrada por la parábola}
  \label{fig:region_parabola_2}
  \includegraphics[width=0.55\textwidth]{Gráficas/GrafArTransvBaseParabola.pdf}
  \propia
\end{figure}

%\begin{center}---------
%faltaria ponerle el numero de figura en la parte de arriba
%\includegraphics[scale=0.3]{GrafArTransvBaseParabola.pdf}
%\textbf{Fuente:} elaboración propia
%\end{center}

Una idea del sólido completo puede darse con la figura \ref{fig:solido_generado}:

\begin{figure}[H]
  \centering
  \caption{Sólido generado}
  \label{fig:solido_generado}
  \includegraphics[width=0.56\textwidth]{Gráficas/GrafArTransvBaseParSolid.pdf}
  \propia
\end{figure}

%\begin{center}---------
%\centerline{ \graf{Sólido generado}}
%\includegraphics[scale=0.4]{GrafArTransvBaseParSolid.pdf}
%\textbf{Fuente:} elaboración propia
%\end{center}

Si la sección transversal para un valor $x$ es un cuadrado, entonces el lado mide $L=4-x^{2}$ y $-2\leq x\leq 2$; por tanto,

\resizebox{\textwidth}{!}{$
  \begin{aligned}
    v&=\int_{a}^{b}A\left( x\right) dx=\int_{-2}^{2}\left(4-x^{2}\right)^{2}dx
    %\\&
    =2\int_{0}^{2}\left( 4-x^{2}\right) ^{2}dx=2\int_{0}^{2}\left(
    16-8x^{2}+x^{4}\right) dx \\
    &=2\left( 16x-\frac{8}{3}x^{3}+\frac{1}{5}x^{5}\right)\Bigg\vert _{0}^{2}=2\left( 32-%
    \frac{64}{3}+\frac{32}{5}\right) %\\&
    =\frac{512}{15}\approx34.133?\text{ unidades c\'ubicas.}%
  \end{aligned}
$}

El siguiente enlace contiene la simulación de la generación del sólido y las secciones transversales:

\begin{center}
  \verb"https://www.geogebra.org/m/bd9ueygs"
\end{center}

%\begin{center}---------
%\centerline{ \graf{Cono circular recto}}%mirar si se puede comentar
%\verb"https://www.geogebra.org/m/bd9ueygs"
%\end{center}

\begin{ejemplo}
  Calcule el volumen del cono circular recto de altura $h$ y radio de base $r.$
\end{ejemplo}
\emph{Solución.} En un cono, podemos identificar que los cortes transversales son círculos concéntricos. Si ubicamos su vértice en el origen, los cortes transversales son perpendiculares al eje $x$ y se verían como en la figura \ref{fig:cono_circular_recto}:

\begin{figure}[H]
  \centering
  \caption{Cono circular recto}
  \label{fig:cono_circular_recto}
  \includegraphics[width=0.47\textwidth]{Gráficas/GrafCono.pdf}
  \propia
\end{figure}

%\begin{center}---------
%\centerline{ \graf{Esfera de radio $r$}}
%\includegraphics[scale=0.35]{GrafCono.pdf}
%\textbf{Fuente:} elaboración propia
%\end{center}

En la figura puede verse también el triángulo rectángulo de donde podemos obtener el valor del radio del corte transversal en función de $x$. Notemos que la pendiente de la recta que une $(0,0)$ con $(h,R)$ es $m=\frac{R}{h}$; luego, la ecuación de la recta que contiene a $(x,r)$ es $y-0=\frac{R}{h}(x-0)$, es decir, $y=\frac{R}{h}x$. Pero $y=r$ para cada $x$; entonces, el área de la sección transversal, es decir, el círculo, es
{\small
  \[
    A(x)=\pi (r)^{2}=\pi y^{2}=\pi\left(\frac{R}{h}x\right)^{2}
  \]
}
entonces el volumen $V$ del cono es
{\small
  \begin{align*}
    V&=\int_{0}^{h}\pi \left(\frac{R}{h}x\right)^{2}dx=\int_{0}^{h}\pi \frac{R^{2}x^{2}}{h^{2}}dx
    %\\&
    =\frac{\pi R^{2}}{h^{2}}\int_{0}^{h}x^{2}dx=\frac{\pi R^{2}}{h^{2}}\int_{0}^{h} x^{2}dx\\
    &=\frac{\pi R^{2}}{h^{2}}\cdot \frac{x^{3}}{3}\Big\vert_{0}^{h}=\frac{\pi
    R^{2}}{h^{2}}\left( \frac{h^{3}}{3}-\frac{0^{3}}{3}\right) =\frac{\pi
    R^{2}h^{3}}{3h^{2}}
    %\\&
    =\frac{\pi R^{2}h}{3}\text{ unidades c\'ubicas.}%
  \end{align*}
}

\begin{ejemplo}
  Calcule el volumen de la esfera de radio $r.$
\end{ejemplo}

\emph{Solución.} Los cortes transversales, en este caso, son círculos centrados en el eje $x$ y con radio dado por la parte superior de la circunferencia $x^2+y^2=r^2$ (figura \ref{fig:esfera_cortes})

\begin{figure}[H]
  \centering
  \caption{Esfera de radio $r$}
  \label{fig:esfera_cortes}
  \includegraphics[width=0.55\textwidth]{Gráficas/GrafEsferaCortes.pdf}
  \propia
\end{figure}

%\begin{center}-----------
%falta ponerle el nombre
%\includegraphics[scale=0.35]{GrafEsferaCortes.pdf}
%\textbf{Fuente:} elaboración propia
%\end{center}

Es decir, el radio del corte tranversal es $y=\sqrt{r^2-x^2}$; así, el área del círculo es $A(x)=\pi y^{2}=\pi (r^{2}-x^{2})$ y entonces el volumen $V$ de la esfera es

\resizebox{\textwidth}{!}{$
  \begin{aligned}
    V&=\int_{-r}^{r}\pi (r^{2}-x^{2})dx=\pi \left(r^{2}x-\frac{x^{3}}{3}\right)\Bigg\vert _{-r}^{r}
    %\\&
    =\pi \left(r^{2}r-\frac{r^{3}}{3}\right)-\pi \left(r^{2}(-r)-\frac{(-r)^{3}}{3}\right)\\
    &=\pi \left(r^{3}-\frac{r^{3}}{3}\right)-\pi \left(-r^{3}+\frac{r^{3}}{3}\right)=\pi \left( r^{3}-
    \frac{r^{3}}{3}+r^{3}-\frac{r^{3}}{3}\right)=\pi \left( 2r^{3}-\frac{2r^{3}%
    }{3}\right) \\
    &=\pi \left( \frac{6r^{3}-2r^{3}}{3}\right) =\pi \left( \frac{4r^{3}}{3}%
    \right)=\frac{4}{3}\pi r^{3}\text{ unidades c\'ubicas.}%
  \end{aligned}
$}

Es decir, el volumen de una esfera de radio $r$ es {\small $\cfrac{4}{3}\pi r^3.$}

\begin{ejemplo}
  Calcule el volumen del cono circular recto truncado de altura $h,$ radio de la base inferior $r$ y radio de la base superior $R.$
\end{ejemplo}

\emph{Solución.} El cono truncado que nos dan puede verse como el sólido que se tiene al girar la región del plano $xy$ acotada por la recta que pasa por los puntos $(0,r)$ y $(h,R)$, el eje $y$ y el eje $x$ para $0\leq x\leq h$. La figura \ref{fig:cono_truncado} nos muestra el cono y la región:

\begin{figure}[H]
  \centering
  \caption{Cono truncado}
  \label{fig:cono_truncado}
  \includegraphics[width=0.47\textwidth]{Gráficas/ConoTruncado.pdf}
  \propia
\end{figure}

%\begin{center}-----------
%\centerline{ \graf{Cono truncado}}
%\includegraphics[scale=0.35]{ConoTruncado.pdf}
%\textbf{Fuente:} elaboración propia
%\end{center}

Encontremos la ecuación de la recta que pasa por los puntos {\small $(0,r)$,} {\small $(h,R)$} y {\small $%
(x,y) $.} La pendiente de la recta es {\footnotesize $m=\cfrac{R-r}{h}$.} Así, la ecuación de la recta para cualquier punto $(x,y)$ es
{\small
  \begin{align*}
    y-r= & \frac{R-r}{h}(x-0) \\
    y-r= & \frac{R-r}{h}x \\
    y=& \frac{R-r}{h}x+r
  \end{align*}
}
Con ello, el área del círculo base es $A(x)=\pi y^{2}=\pi \left[ \left(\frac{R-r}{h}\right)x+r\right] ^{2}$, y con ello,

\resizebox{\textwidth}{!}{$
  \begin{aligned}
    V=&\pi \int_{0}^{h}\left[ \left(\frac{R-r}{h}\right)x+r\right] ^{2}dx=\pi \int_{0}^{h}%
    \left[ \left( \frac{R-r}{h}x\right) ^{2}+2r\left(\frac{R-r}{h}\right)x+r^{2}\right]
    dx\\
    =&\pi \int_0^h \left(\frac{R-r}{h}\right)^{2}x^{2}+2r\frac{R-r}{h}x+r^{2}dx=\pi \left[\left(\frac{R-r}{h}\right)^{2}\frac{x^{3}}{3}+2r\frac{R-r}{h}\frac{x^{2}%
    }{2}+r^{2}x\right] _{0}^{h}\\
    =&\pi \left[ \left(\frac{R-r}{h}\right)^{2}\frac{h^{3}}{3}+2r(\frac{R-r}{h})\frac{h^{2}%
    }{2}+r^{2}h\right]=\pi \left[ (R-r)^{2}\frac{h}{3}+r\left( R-r\right) h+r^{2}h\right] \\
    =&\pi \left[ (R^{2}-2Rr+r^{2})\frac{h}{3}+\left( Rr-r^{2}\right) h+r^{2}h%
    \right]=\pi h\left[ \frac{R^{2}-2Rr+r^{2}}{3}+Rr-r^{2}+r^{2}\right] \\
    =&\pi h\left[ \frac{R^{2}-2Rr+r^{2}+3Rr}{3}\right] =\pi h\left[ \frac{%
    R^{2}+Rr+r^{2}}{3}\right] \text{ unidades c\'ubicas.}%
  \end{aligned}
$}

\begin{ejemplo}
  Calcule el volumen de un casquete de altura (profundidad) $h$ en una esfera de radio \textit{$r$:}
\end{ejemplo}

\begin{figure}[H]
  \centering
  \caption{Casquete esférico profundidad $h$}
  \label{fig:casquete_esferico}
  \includegraphics[width=0.25\textwidth]{Gráficas/CasqueteEsferico.pdf}
  \propia
\end{figure}

%\begin{center}----------
%\centerline{ \graf{Casquete esférico profundidad $h$}}
%\includegraphics[scale=0.3]{CasqueteEsferico.pdf}
%\textbf{Fuente:} elaboración propia
%\end{center}

\emph{Solución.} En la figura del ejemplo, ponemos el casquete en esa posición, pero tengamos en cuenta que su volumen no cambia si su posición es otra. Acá, simplemente se ubica así para seguir con los límites de la integral para el eje $x$.

Así, si ubicamos nuestra esfera en el plano cartesiano,

\begin{figure}[H]
  \centering
  \caption{Casquete esférico en el plano $xy$}
  \includegraphics[width=0.38\textwidth]{Gráficas/CasqueteEsfericoEnPlano.pdf}
  \propia
\end{figure}

%\begin{center}-----------
%\centerline{ \graf{Casquete esférico en el plano $xy$}}
%\includegraphics[scale=0.3]{CasqueteEsfericoEnPlano.pdf}
%\textbf{Fuente:} elaboración propia
%\end{center}

vemos que los límites de la integral van desde $x=r-h$ hasta $r$. Además, la ecuación de la circunferencia es $x^{2}+y^{2}=r^{2}$, es decir, $y^{2}=r^{2}-x^{2}$; así, $A(x)=\pi (r^{2}-x^{2})$ y
{\small
  \begin{align*}
    V&=\int_{r-h}^{r} \pi (r^{2}-x^{2})dx=\pi \left(r^{2}x-\frac{x^{3}}{3}\right)\Bigg|_{r-h}^{r}\\
    &=\pi \left[ r^{2}r-\frac{r^{3}}{3}-\left(r^{2}(r-h)-\frac{(r-h)^{3}}{3}\right)\right]\\
    &=\pi \left( r^{3}-\frac{r^{3}}{3}-r^{3}+r^{2}h+\frac{r^{3}-3r^{2}h+3rh^{2}-h^{3}}{3}\right)\\
    %\end{align*}
    %\begin{align*}
    &=\pi \left( -\frac{r^{3}}{3}+r^{2}h+\frac{r^{3}-3r^{2}h+3rh^{2}-h^{3}}{3}%
    \right) \\
    &=\pi \left( -\frac{r^{3}}{3}+r^{2}h+\frac{r^{3}}{3}-r^{2}h+rh^{2}-%
    \frac{h^{3}}{3}\right)\\
    &=\pi \left( rh^{2}-\frac{h^{3}}{3}\right) =\pi h^{2}\left( r-\frac{h}{3}%
    \right) u^{3}\text{ unidades c\'ubicas.}%
  \end{align*}
}

\begin{ejemplo}
  Halle el volumen del sólido cuya base está dada por la región $R=\left\{ \left( x,y\right) :\frac{x^{2}}{100}+\frac{y^{2}}{36}\leq 1\right\} $ y cuyas secciones transversales perpendiculares al eje $x$ son triángulos rectángulos isósceles con un lado sobre la base.
\end{ejemplo}

\emph{Solución.} La región $R$ del plano $xy$ es una elipse alargada en el eje $x$, como en la figura \ref{fig:elipse_base_solido}:

\begin{figure}[H]
  \centering
  \caption{Elipse base del sólido}
  \label{fig:elipse_base_solido}
  \includegraphics[width=0.38\textwidth]{Gráficas/SeccTransElipTriangBase.pdf}
  \propia
\end{figure}

%\begin{center}----------
%\centerline{ \graf{Elipse base del sólido}}
%\includegraphics[scale=0.3]{SeccTransElipTriangBase.pdf}
%\textbf{Fuente:} elaboración propia
%\end{center}

Además, nos dicen que las secciones transversales perpendiculares al eje $x$ son triángulos rec\-tán\-gu\-los isósceles, es decir, en un punto $x$ tenemos un triángulo de base y altura de igual tamaño. En la figura 48 podemos ver un triángulo y la idea del sólido que se forma:

\begin{figure}[H]
  \centering
  \caption{Secciones transversales}
  \includegraphics[width=0.65\textwidth]{Gráficas/SeccTransElipTriang.pdf}
  \propia
\end{figure}

%\begin{center}----------
%\centerline{ \graf{Secciones transversales}}
%\includegraphics[scale=0.35]{SeccTransElipTriang.pdf}
%\textbf{Fuente:} elaboración propia
%\end{center}

Y, en definitiva, el sólido se ve de la siguiente manera:

\begin{figure}[H]
  \centering
  \caption{El sólido}
  \includegraphics[width=0.65\textwidth]{Gráficas/SeccTransElipTriangSolid.pdf}
  \propia
\end{figure}

%\begin{center}-----------
%\centerline{ \graf{El sólido}}
%\includegraphics[scale=0.35]{SeccTransElipTriangSolid.pdf}
%\textbf{Fuente:} elaboración propia
%\end{center}

Ahora bien, haciendo unos cálculos similares a los ejemplos anteriores, se tiene que el área del triángulo es como se usa en la siguiente integral:
{\small
  \begin{align*}
    v&=\int_{a}^{b}A\left( x\right) dx =\int_{-10}^{10}72\underset{\text{es par}}{\left( \underbrace{1-\frac{x^{2}}{%
    100}}\right) }dx=2\times 72\int_{0}^{10}\left( 1-\frac{x^{2}}{100}\right) dx \\
    &=144\left( x-\frac{x^{3}}{300}\right)\Bigg\vert_{0}^{10}
    %\\&
    =144\left( 10-\frac{1000}{300}\right) =960\text{ unidades c\'ubicas}%
  \end{align*}
}

En el siguiente enlace podemos encontrar la simulación del sólido y de las secciones transversales:

\begin{center}
  \verb"https://www.geogebra.org/m/jtzvcr7m"
\end{center}

%\begin{center}-----------
%\centerline{ \graf{Las curvas}}%¿ese nombre va ahí?
%\verb"https://www.geogebra.org/m/jtzvcr7m"
%\end{center}

\begin{ejemplo}
  Halle el volumen del sólido cuya base es la región limitada por $y=x\ln
  x $ y $y=4\ln x$ y cuyas secciones transversales perpendiculares al eje $x$ son cuadrados con un lado sobre la base.
\end{ejemplo}
%%\begin{multicols}{2}
\emph{Solución.} La figura \ref{fig:secciones_transversales_final} muestra las curvas y la región que encierran las funciones dadas:

\begin{figure}[H]
  \centering
  \caption{Base del solido}
  \label{fig:secciones_transversales_final}
  \includegraphics[width=0.35\textwidth]{Gráficas/Reg4LnxXlnx.pdf}
  \propia
\end{figure}

%\begin{center}----------
%\centerline{ \graf{Las secciones transversales}}
%\includegraphics[scale=0.35]{Reg4LnxXlnx.pdf}
%\textbf{Fuente:} elaboración propia
%\end{center}

Hallemos ahora los puntos de corte de las curvas: si $x\ln x=4\ln x$, entonces $(x-4)\ln x=0$, luego $x=4$ y $x=1.$

Las secciones transversales se ven como en la figura \ref{fig:secciones_transversales_b}:

\begin{figure}[H]
  \centering
  \caption{Las secciones transversales}
  \label{fig:secciones_transversales_b}
  \includegraphics[width=0.45\textwidth]{Gráficas/SeccTrans4LnxXlnx.pdf}
  \propia
\end{figure}

%\begin{center}-----------
%\includegraphics[scale=0.35]{SeccTrans4LnxXlnx.pdf}
%\textbf{Fuente:} elaboración propia
%\end{center}

El sólido completo se vería de la siguiente manera:
%%\end{multicols}

\begin{figure}[H]
  \centering
  \caption{El sólido}
  \label{fig:solido_}
  \includegraphics[width=0.45\textwidth]{Gráficas/SeccTrans4LnXlnSolid.pdf}
  \propia
\end{figure}

%\begin{center}----------
%\centerline{ \graf{El sólido}}
%\includegraphics[scale=0.35]{SeccTrans4LnXlnSolid.pdf}
%\textbf{Fuente:} elaboración propia
%\end{center}

Ahora bien, si un lado del cuadrado está en la región, entonces su longitud es $L=4\ln x-x\ln x=\left( 4-x\right) \ln x$, es decir, el área del cuadrado es $\left( 4-x\right)^2 \ln^2 x$ y, así, el volumen del sólido es:
{\small
  \[
    v =\int_{1}^{4}\left( 4-x\right) ^{2}\ln ^{2}xdx=\int_{1}^{4}\left( 16-8x+x^{2}\right) \ln ^{2}xdx
  \]
}
Si hacemos $y=\ln x$, entonces $x=e^{y}$ y $dx=e^{y}dy$. Y cuando $x=1\ \rightarrow y=0$ y si $x=4\ \rightarrow y=\ln 4$. Es decir, la integral para el volumen se convierte en
{\small
  \begin{align*}
    \int_{1}^{4}&\left( 16-8x+x^{2}\right) \ln ^{2}xdx=\int_{0}^{\ln 4}\left(
    16-8e^{y}+e^{2y}\right) y^{2}e^{y}dy \\
    &=\int_{0}^{\ln 4}\left(
    16y^{2}e^{y}-8y^{2}e^{2y}+y^{2}e^{3y}\right) dy
    %\\&
    %=\int_{0}^{\ln 4}\left( 16y^{2}e^{y}-2y^{2}e^{2y}+y^{2}e^{3y}\right)dy\\
    %&=\left( \frac{y^{2}e^{3y} }{3}-y^{2}e^{2y}+16y^{2}e^{y}-\frac{2ye^{3y}}{9}%
    %+ye^{2y}%\right.\\&\qquad\left.
    %-32ye^{y}+\frac{2 e^{3y}}{27}-\frac{e^{2y} }{2}%
    %+32e^{y}\right)\Bigg\vert_{0}^{\ln 4}\\
    %&=\frac{2272}{9}\ln ^{2}2-\frac{13\,472}{27}\ln 2+\frac{3377}{12}%\\&
    %\approx 56.850.
  \end{align*}
}

Calculando cada integral por el método de partes con integración tabular
obtenemos:
%\begin{center}-No tiene figura
%\centerline{ \graf{Cono como sólido de revolución}}

\begin{tabular}{lll}
  & & $e^{y}$ \\
  $16y^{2}$ & $\overset{+}{\longrightarrow }$ & $e^{y}$ \\
  $32y$ & $\overset{-}{\longrightarrow }$ & $e^{y}$ \\
  $32$ & $\overset{+}{\longrightarrow }$ & $e^{y}$ \\
  $0$ & &
\end{tabular},
\quad\quad
\begin{tabular}{lll}
  & & $e^{2y}$ \\
  $8y^{2}$ & $\overset{+}{\longrightarrow }$ & $\frac{1}{2}e^{2y}$ \\
  $16y$ & $\overset{-}{\longrightarrow }$ & $\frac{1}{4}e^{2y}$ \\
  $16$ & $\overset{+}{\longrightarrow }$ & $\frac{1}{8}e^{2y}$ \\
  $0$ & &
\end{tabular},
\quad\quad
\begin{tabular}{lll}
  & & $e^{3y}$ \\
  $y^{2}$ & $\overset{+}{\longrightarrow }$ & $\frac{1}{3}e^{3y}$ \\
  $2y$ & $\overset{-}{\longrightarrow }$ & $\frac{1}{9}e^{3y}$ \\
  $2$ & $\overset{+}{\longrightarrow }$ & $\frac{1}{27}e^{3y}$ \\
  $0$ & &
\end{tabular}
%\end{center}

y con ello,

\resizebox{\textwidth}{!}{$
  \begin{aligned}
    &\int_{0}^{\ln 4}\left(
    16y^{2}e^{y}-8y^{2}e^{2y}+y^{2}e^{3y}\right) dy\\
    &=\left( e^{y}\left( 16y^{2}-32y+32\right) -e^{2y}\left( 4y^{2}-4y+2\right)
    +e^{3y}\left( \frac{1}{3}y^{2}-\frac{2}{9}y+\frac{2}{27}\right) \right)\Bigg|_{0}^{\ln 4} \\
    &=4\left( 16\ln ^{2}4-32\ln 4+32\right) -16\left( 4\ln ^{2}4-4\ln
    4+2\right) +64\left( \frac{1}{3}\ln ^{2}4-\frac{2}{9}\ln 4+\frac{2}{27}%
    \right)\\
    &\quad -\left( 32-2+\frac{2}{27}\right) \\
    &=\frac{256}{3}\ln ^{2}2-\frac{1408}{9}\ln 2+\frac{212}{3} \\
    &\approx 3,2263\text{ unidades cúbicas.}%
  \end{aligned}
$}

En el siguiente enlace podemos encontrar la simulación del sólido y de las secciones transversales:

\begin{center}
  \verb"https://www.geogebra.org/m/buckfsn4"
\end{center}

%\begin{center}----------
%\centerline{ \graf{La región $R$}}
%\verb"https://www.geogebra.org/m/buckfsn4"
%\end{center}

\section{Volumen de sólidos de revolución}

Un sólido de revolución es aquel que se obtiene al hacer girar una región plana y acotada alrededor de una recta fija. Por ejemplo, cuando hacemos girar el triángulo con vértices $(0,0),\ (3,0)$ y $(3,2)$ alrededor del eje $x$, tenemos un cono circular recto de radio $2$ y altura $3$ (figura \ref{fig:cono_solido_revolucion}):

\begin{figure}[H]
  \centering
  \caption{Cono como sólido de revolución}
  \label{fig:cono_solido_revolucion}
  \includegraphics[width=0.54\textwidth]{Gráficas/SolRevIntroducc.pdf}
  \propia
\end{figure}

%\begin{center}----------
%\centerline{ \graf{Cono como sólido de revolución}}
%\includegraphics[scale=0.35]{SolRevIntroducc.pdf}
%\textbf{Fuente:} elaboración propia
%\end{center}

Nos interesa ahora medir el volumen de sólidos de revolución. Tomemos algunos casos.

\subsection*{Método de discos}
%%\begin{multicols}{2}
Supongamos que se tiene la región $R$ del plano $xy$ acotada por las figuras de $y=f\left( x\right) $, $y=0$, $x=a$, y $x=b $ (figura \ref{fig:la_region}):

\begin{figure}[H]
  \centering
  \caption{La región $R$}
  \label{fig:la_region}
  \includegraphics[width=0.40\textwidth]{Gráficas/SolRevRegiDefin.pdf}
  \propia
\end{figure}

%\begin{center}------------
%\centerline{ \graf{La región $R$}}
%\includegraphics[scale=0.35]{SolRevRegiDefin.pdf}
%\textbf{Fuente:} elaboración propia
%\end{center}

Si la región $R$ se hace girar alrededor del eje $x$, entonces el sólido $S$ así generado es un \textbf{sólido de revolución} (figura \ref{fig:solido_de_revolu}).

\begin{figure}[H]
  \centering
  \caption{Sólido de revolución}
  \label{fig:solido_de_revolu}
  \includegraphics[width=0.25\textwidth]{Gráficas/SolRevDefin.pdf}
  \propia
\end{figure}

%\begin{center}-----------
%\centerline{ \graf{Sólido de revolución}}
%\includegraphics[scale=0.35]{SolRevDefin.pdf}
%\textbf{Fuente:} elaboración propia
%\end{center}

%%\end{multicols}

Una sección transversal que pasa por $x$ y es perpendicular al eje $x$ es un disco\footnote{ Región de un plano delimitada por una circunferencia.} con radio $r=\left\vert y\right\vert =\left\vert f\left( x\right) \right\vert $ y cuya área es
\begin{equation*}
  A\left( x\right) =\pi y^{2}=\pi \left[ f\left( x\right) \right] ^{2}
\end{equation*}

\begin{figure}[H]
  \centering
  \caption{Sección transversal}
  \includegraphics[width=0.28\textwidth]{Gráficas/SolRevSeccTransDefin.pdf}
  \propia
\end{figure}

%\begin{center}--------
%\centerline{ \graf{Sección transversal}}
%\includegraphics[scale=0.35]{SolRevSeccTransDefin.pdf}
%\textbf{Fuente:} elaboración propia
%\end{center}

Y con el uso de la fórmula
\begin{equation}\label{DiscosAlrx}
  v=\int_{a}^{b}A\left( x\right) dx=\int_{a}^{b}\pi %
  \left[ f\left( x\right) \right] ^{2}dx
\end{equation}
tenemos el volumen de $S$. El método con el cual empleamos esta fórmula es usualmente llamado el \textit{método de los discos}.

Si la región es acotada por las curvas $x=g\left( y\right)$, $x=0$, $y=c $ y $y=d$, rota alrededor del eje $y$, entonces el volumen de este sólido de
revolución es
\begin{equation}\label{DiscosAlry}
  v=\int_{c}^{d}\pi \left[
  g\left( y\right) \right] ^{2}dy
\end{equation}

\begin{ejemplo}
  Determine el volumen del sólido generado al hacer girar la región determinada por las curvas dadas por $x+2y=2$, $y=0$ y $x=0$ alrededor del eje $x$.
\end{ejemplo}

\emph{Solución.} Si vemos la recta en términos de $x$, tenemos la función $f(x)=y=\frac{2-x}{2}$. En la figura \ref{fig:solido_revolucion_ejemplo} se ve la región que gira y un bosquejo del sólido de revolución:

\begin{figure}[H]
  \centering
  \caption{Región y sólido del ejemplo}
  \label{fig:solido_revolucion_ejemplo}
  \includegraphics[width=0.63\textwidth]{Gráficas/SolRevDiscEjemRect.pdf}
  \propia
\end{figure}

%\begin{center}------------
%\centerline{ \graf{Región y sólido del ejemplo}}
%\includegraphics[scale=0.35]{SolRevDiscEjemRect.pdf}
%\textbf{Fuente:} elaboración propia
%\end{center}

%%\begin{multicols}{2}

El disco formado en este caso puede verse en la figura \ref{fig:disco_ejemplo}:

\begin{figure}[H]
  \centering
  \caption{Disco del ejemplo}
  \label{fig:disco_ejemplo}
  \includegraphics[width=0.35\textwidth]{Gráficas/SolRevDiscEjemRectDiscos.pdf}
  \propia
\end{figure}

%\begin{center}-----------
%\centerline{ \graf{Disco del ejemplo}}%¿este es el nombre??
%\includegraphics[scale=0.35]{SolRevDiscEjemRectDiscos.pdf}
%\textbf{Fuente:} elaboración propia
%\end{center}

La sección transversal es un disco cuyo radio {\small $r=\cfrac{2-x}{2}$;} as%
í, el área del disco es
{\small
  \begin{align*}
    A(x)&=\pi (\text{radio})^{2}=\pi y^{2} \\
    & =\pi \left(\frac{2-x}{2}\right)^{2}\\
    &=\pi \left(1-\frac{x}{2}\right)^{2}
  \end{align*}
}
%%\end{multicols}

Por tanto, el volumen del sólido de revolución es

\resizebox{\textwidth}{!}{$
  \begin{aligned}
    V&=\int_{0}^{2}\pi \left(1-\frac{x}{2}\right)^{2}dx=\int_{0}^{2}\pi \left(1-x+\frac{x^{2}}{4}\right)dx %\\&
    =\int_{0}^{2} \pi \left(1-x+\frac{x^{2}}{4}\right)dx\\
    &=\pi \left(x-\frac{x^{2}}{2}+\frac{x^{3}}{12}\right)\Bigg|_{0}^{2}=\pi \left[2-\frac{2^{2}}{2}+\frac{2^{3}}{12}-\left(0-\frac{0^{2}}{2}+\frac{0^{3}}{12}\right)\right]
    \\
    %& \\
    &=\pi \left( 2-2+\frac{8}{12}\right) =\pi \cdot \frac{2}{3}=\frac{2\pi }{3}\text{ unidades c\'ubicas.}%
  \end{aligned}
$}

En el siguiente enlace podemos encontrar la simulación del sólido y de los discos que se forman en el proceso:

\begin{center}
  \verb"https://www.geogebra.org/m/k5tzayav"
\end{center}

%\begin{center}----------
%\centerline{ \graf{Región y sólido del ejemplo}}
%\verb"https://www.geogebra.org/m/k5tzayav"
%\end{center}

\begin{ejemplo}
  Halle el volumen del sólido de revolución generado por la región plana delimitada por la curva $x=\tan (\frac{\pi y}{4}),x=0,y=1$ alrededor del eje $y.$
\end{ejemplo}

\emph{Solución.} En la figura \ref{fig:disco_del_ejemplo} podemos ver que la sección transversal es un disco cuyo radio $r=\tan (\frac{\pi y}{4}),$ luego su área es
{\small
  \[
    A(y)=\pi x^{2}=\pi \left(\tan \left(\frac{\pi y}{4}\right)\right)^{2}
  \]
}

\begin{figure}[H]
  \centering
  \caption{Región plana y sólido del ejemplo}
  \label{fig:disco_del_ejemplo}
  \includegraphics[width=0.46\textwidth]{Gráficas/SolRevDiscEjemTang.pdf}
  \propia
\end{figure}

%\begin{center}----------
%\centerline{ \graf{Disco del ejemplo}}
%\includegraphics[scale=0.35]{SolRevDiscEjemTang.pdf}
%\textbf{Fuente:} elaboración propia
%\end{center}

El disco formado en este caso puede verse en la figura \ref{fig:discoo_del_ejemplo}:

\begin{figure}[H]
  \centering
  \caption{Disco del ejemplo}
  \label{fig:discoo_del_ejemplo}
  \includegraphics[width=0.35\textwidth]{Gráficas/SolRevDiscEjemTangDisco.pdf}
  \propia
\end{figure}

%\begin{center}------------
%\includegraphics[scale=0.35]{SolRevDiscEjemTangDisco.pdf}
%\textbf{Fuente:} elaboración propia
%\end{center}

Por tanto,

\resizebox{\textwidth}{!}{$
  \begin{aligned}
    V&=\int_{0}^{1}\pi \left(\tan \left(\frac{\pi y}{4}\right)\right)^{2}dy=\int_{0}^{1} \pi \tan ^{2}\left(\frac{\pi y}{4}\right)dy=\pi \int_{0}^{1} \left( \sec ^{2}\left(\frac{\pi y}{4}\right)-1\right) dy
    %\\ &
    \\
    &= \pi \left( \frac{4}{\pi }\tan \left(\frac{\pi y}{4}\right)-y\right) \Bigg|_{0}^{1}=\pi \left[ \left( \frac{4}{\pi }\tan \left(\frac{\pi }{4}\right)-1\right) -\left(\frac{4}{\pi }\tan \left(\frac{\pi (0)}{4}\right)-(0)\right)\right]\\
    &=\pi \left[ \frac{4}{\pi }\tan \left(\frac{\pi }{4}\right)-1-\frac{4}{\pi }\tan (0)\right] =\pi \left( \frac{4}{\pi }(1)-1-0\right)=\left( 4-\pi \right) \text{ unidades c\'ubicas.}%
  \end{aligned}
$}

En el siguiente enlace puede encontrarse la simulación del sólido y de los discos que se forman en el proceso:

\begin{center}
  \verb"https://www.geogebra.org/m/j3zc2tre"
\end{center}

%\begin{center}-----------
%\centerline{ \graf{Región a girar}}%¿este es el nombre??
%\verb"https://www.geogebra.org/m/j3zc2tre"
%\end{center}

\begin{ejemplo}
  Halle el volumen del sólido de revolución obtenido al girar la región acotada por las figuras dadas por $y=\ln x$, $x=1$ y $x=4$ alrededor del eje $x$.
\end{ejemplo}
\emph{Solución.} La región que forman las curvas es

\begin{figure}[H]
  \centering
  \caption{Región a girar}
  \includegraphics[width=0.45\textwidth]{Gráficas/SolRevDiscEjeLnRecta.pdf}
  \propia
\end{figure}

%\begin{center}--------
%\centerline{ \graf{El radio en el sólido}}
%\includegraphics[scale=0.35]{SolRevDiscEjeLnRecta.pdf}
%\textbf{Fuente:} elaboración propia
%\end{center}

Al girar, se identifica el radio $r=\ln x$:

\begin{figure}[H]
  \centering
  \caption{El radio $r=\ln x$}
  \includegraphics[width=0.40\textwidth]{Gráficas/SolRevDiscEjeLnRecta2.pdf}
  \propia
\end{figure}

%\begin{center}------------
%\includegraphics[scale=0.35]{SolRevDiscEjeLnRecta2.pdf}
%\textbf{Fuente:} elaboración propia
%\end{center}

Y un disco formado en la región puede verse en color rojo:

\begin{figure}[H]
  \centering
  \caption{El radio en el sólido}
  \includegraphics[width=0.34\textwidth]{Gráficas/SolRevDiscEjeLnRecta3.pdf}
  \propia
\end{figure}

%\begin{center}-----------
%\centerline{ \graf{El radio en el sólido}}
%\includegraphics[scale=0.45]{SolRevDiscEjeLnRecta3.pdf}
%\textbf{Fuente:} elaboración propia
%\end{center}

En las figuras podemos ver que el volumen del sólido está dado por
{\small
  \[
    v=\int_{1}^{4}\pi \ln ^{2}xdx
  \]
}
Podemos hacer una sustitución y escribirla en términos de exponenciales, como ya hemos hecho en otros ejemplos, o podemos integrar por partes directamente. Usemos esta vez la segunda opción. Si $u=\ln^2x$, entonces
{\footnotesize $du=\cfrac{2\ln x}{x}dx$,} y si $dv=dx$, entonces $v=x$; con ello,

\resizebox{\textwidth}{!}{$
  \begin{aligned}
    v&=\int_{1}^{4}\pi \ln ^{2}xdx=\pi\left(\ln^2x\times x\Big\vert_1^4-\int_1^4x\times\cfrac{2\ln x}{x}\right) =\pi\left(4\ln^24-0-\int_1^42\ln x\, dx\right)\\
    &=\pi\left(4\ln^24-2(x\ln x-x)\Big\vert_1^4\right)=\pi\left(4\ln^24-2(4\ln 4-4)+2(1\ln 1-1)\right)
    \\
    &=\pi\left(4\ln^24-8\ln 4+6\right)\approx 8.1584\text{ unidades c\'ubicas.}%
  \end{aligned}
$}

que corresponde al volumen del siguiente sólido:

\begin{figure}[H]
  \centering
  \caption{El sólido}
  \includegraphics[width=0.38\textwidth]{Gráficas/SolRevDiscEjeLnRecta4.pdf}
  \propia
\end{figure}

%\begin{center}-----------
%\centerline{ \graf{El sólido}}
%\includegraphics[scale=0.35]{SolRevDiscEjeLnRecta4.pdf}
%\textbf{Fuente:} elaboración propia
%\end{center}

En el siguiente enlace puede encontrarse la simulación del sólido y de los discos que se forman en el proceso:

\begin{center}
  \verb"https://www.geogebra.org/classic/qnjuzpaz"
\end{center}

%\begin{center}-----------
%\centerline{ \graf{La región}}
%\verb"https://www.geogebra.org/classic/qnjuzpaz"
%\end{center}

\subsection*{Método de arandelas}
%%\begin{multicols}{2}
Supongamos ahora que se tiene la región $R$ acotada por las curvas dadas por $y=f\left( x\right) $, $\ y=g\left( x\right) $, $x=a$, y $x=b$ (donde $f\left( x\right) \geq g\left( x\right) \geq 0,$ $a\leq x\leq b$):

\begin{figure}[H]
  \centering
  \caption{Región a girar}
  \includegraphics[width=0.46\textwidth]{Gráficas/SolRevArandRegDefin.pdf}
  \propia
\end{figure}

%\begin{center}-----------
%\centerline{ \graf{El sólido}}
%\includegraphics[scale=0.35]{SolRevArandRegDefin.pdf}
%\textbf{Fuente:} elaboración propia
%\end{center}

Supongamos que la región se hace girar alrededor del eje $x$. El sólido así generado tiene un orificio en su centro:

\begin{figure}[H]
  \centering
  \caption{Orificio en su centro}
  \includegraphics[width=0.33\textwidth]{Gráficas/SolRevArandDefin.pdf}
  \propia
\end{figure}

%\begin{center}------------
%\includegraphics[scale=0.3]{SolRevArandDefin.pdf}
%\textbf{Fuente:} elaboración propia
%\end{center}

o mejor

\begin{figure}[H]
  \centering
  \caption{El sólido}
  \includegraphics[width=0.34\textwidth]{Gráficas/SolRevArandDefin3d.pdf}
  \propia
\end{figure}

%%\end{multicols}
%\begin{center}--------
%\centerline{ \graf{El sólido}}
%\includegraphics[scale=0.35]{SolRevArandDefin3d.pdf}
%\textbf{Fuente:} elaboración propia
%\end{center}

Ahora bien, si $S$ es el sólido de revolución así generado, entonces su volumen está dado por
{\small
  \begin{equation}\label{ArandAlrx}
    v=\int_{a}^{b}\pi \left( \left[ f\left( x\right) \right] ^{2}-\left[
    g\left( x\right) \right] ^{2}\right) dx
  \end{equation}
}
es decir, para un $x$ dado, la sección transversal es una arandela de radios $\vert f(x)\vert$ y $\vert g(x) \vert$. El área de la arandela es
\[
  \pi(f(x))^2-\pi(g(x))^2=\pi \left( (f(x))^2-(g(x))^2 \right)
\]

\begin{figure}[H]
  \centering
  \caption{Una arandela}
  \includegraphics[width=0.36\textwidth]{Gráficas/SolRevArandela.pdf}
  \propia
\end{figure}

%\begin{center}------------
%\centerline{ \graf{Una arandela}}
%\includegraphics[scale=0.3]{SolRevArandela.pdf}
%\textbf{Fuente:} elaboración propia
%\end{center}

Cuando calculamos el volumen de un sólido con estas características (secciones transversales arandelas), decimos que estamos usando el método de las arandelas.

\begin{ejemplo}
  Halle el volumen del sólido de revolución generado al girar la
  región limitada por $y=x^{2}$ y $y=x$ en torno a la recta $y=3$.
\end{ejemplo}

\emph{Solución.} La figura 69 nos permite identificar que, al girar la región alrededor de $y=3$, las secciones transversales tienen forma de arandela con radio mayor $R$, medido desde $y=3$ hasta la parábola $y=x^2$: mejor dicho, $R=3-x^2$ y radio menor $r$ medido desde $y=3$ hasta la recta $y=x$, es decir, $r=3-x$. Además, los puntos de corte de las curvas están dados por $x=0$ y $x=1$.

\begin{figure}[H]
  \centering
  \caption{Región y arandelas del ejemplo}
  \includegraphics[width=0.44\textwidth]{Gráficas/SolRevArandEjeRecParab.pdf}
  \propia
\end{figure}

%\begin{center}-----------
%\centerline{ \graf{Región y arandelas del ejemplo}}
%\includegraphics[scale=0.35]{SolRevArandEjeRecParab.pdf}
%\textbf{Fuente:} elaboración propia
%\end{center}

La arandela, de manera más precisa, la que se forma en este caso, puede verse en la siguiente figura:
%\clearpage

\begin{figure}[H]
  \centering
  \caption{Arandela en el ejemplo}
  \includegraphics[width=0.33\textwidth]{Gráficas/SolRevArandEjeRecParabArandel.pdf}
  \propia
\end{figure}

%\begin{center}------------
%\centerline{ \graf{Arandela en el ejemplo}}
%\includegraphics[scale=0.35]{SolRevArandEjeRecParabArandel.pdf}
%\textbf{Fuente:} elaboración propia
%\end{center}

Por tanto, el volumen del sólido es
{\small
  \begin{align*}
    v&=\int_{a}^{b}\pi \left( R^{2}-r^{2}\right) dx=\pi \int_{0}^{1}\left( \left( 3-x^{2}\right) ^{2}-\left( 3-x\right)
    ^{2}\right) dx \\
    &=\pi \int_{0}^{1}\left( x^{4}-6x^{2}+9-x^{2}+6x-9\right) dx=\pi \int_{0}^{1}\left( 6x-7x^{2}+x^{4}\right) dx \\
    &=\pi \left( 3x^{2}-\frac{7}{3}x^{3}+\frac{1}{5}x^{5}\right)\Bigg\vert_{0}^{1}=\pi
    \left( 3-\frac{7}{3}+\frac{1}{5}\right) \\
    &=\frac{13}{15}\pi\approx2.7227\text{ unidades c\'ubicas.}%
  \end{align*}
}
El sólido completo al que le calculamos el volumen se ve en la figura \ref{fig:region_arandelas}:

\begin{figure}[H]
  \centering
  \caption{El sólido}
  \label{fig:region_arandelas}
  \includegraphics[width=0.32\textwidth]{Gráficas/SolRevArandEjeRecParabSolido.pdf}
  \propia
\end{figure}

%\begin{center}----------
%\centerline{ \graf{Región y arandelas del ejemplo}}
%\includegraphics[scale=0.35]{SolRevArandEjeRecParabSolido.pdf}
%\textbf{Fuente:} elaboración propia
%\end{center}

En el siguiente enlace puede encontrarse la simulación del sólido y de las arandelas que se forman:

\begin{center}
  \verb"https://www.geogebra.org/m/ta7wvjgh"
\end{center}

%\begin{center}----------
%\centerline{ \graf{Región del ejemplo}}
%\verb"https://www.geogebra.org/m/ta7wvjgh"
%\end{center}

\begin{ejemplo}
  Halle el volumen del sólido de revolución generado al girar la
  región limitada por $x=y^{2}$ y $x=y+2$ alrededor de la recta $x=4$.
\end{ejemplo}

\emph{Solución.}

\begin{figure}[H]
  \centering
  \caption{Región a girar}
  \label{fig:aran_dela}
  \includegraphics[width=0.30\textwidth]{Gráficas/SolRevArandEjeRecParabol2.pdf}
  \propia
\end{figure}

%\begin{center}----------
%\centerline{ \graf{Arandelas del ejemplo}}
%\includegraphics[scale=0.35]{Cálculo integral/Gráficas/SolRevArandEjeRecParabol2.pdf}
%\textbf{Fuente:} elaboración propia
%\end{center}

%Por favor, confirmen si debe ser la imagen SolRevArandEjeRecParab2Solido, porque el nombre de la imagen estaba incompleto y no se sabía si era sólido o arándela.

Al igualar las expresiones, tenemos que los puntos de corte de las curvas se dan cuando $y=-1$ y $y=2$, es decir, los puntos de corte son $(1,-1)$ y $(4,2)$. La figura \ref{fig:aran_dela} nos permite identificar que los cortes transversales son arandelas con $R=4-y^{2}$ y $r=4-\left( y+2\right) =2-y$. La arandela que se forma en este caso puede verse en la figura \ref{fig:arandelas_}:

\begin{figure}[H]
  \centering
  \caption{Arandelas del ejemplo}
  \label{fig:arandelas_}
  \includegraphics[width=0.30\textwidth]{Gráficas/SolRevArandEjeRecParab2Arandela.pdf}
  \propia
\end{figure}

%\begin{center}-----------
%\includegraphics[scale=0.45]{SolRevArandEjeRecParab2Arandela.pdf}
%\textbf{Fuente:} elaboración propia
%\end{center}

Entonces, el volumen del sólido es

\resizebox{\textwidth}{!}{$
  \begin{aligned}
    v&=\int_{a}^{b}\pi \left( R^{2}-r^{2}\right) dy =\pi \int_{-1}^{2}\left( \left( 4-y^{2}\right) ^{2}-\left( 2-y\right)
    ^{2}\right) dy \\
    &=\pi \int_{-1}^{2}\left( y^{4}-8y^{2}+16-y^{2}+4y-4\right) dy
    =\pi \int_{-1}^{2}\left( 12+4y-9y^{2}+y^{4}\right) dy \\
    &=\pi \left( 12y+2y^{2}-3y^{3}+\frac{1}{5}y^{5}\right)\Bigg\vert _{-1}^{2}=\pi \left( \left( 24+8-24+\frac{32}{5}\right) -\left( -12+2+3-\frac{1}{5}
    \right) \right) \\
    &=\frac{108}{5}\pi =21.6\pi \approx67.858\text{ unidades c\'ubicas.}%
  \end{aligned}
$}

y el sólido al que le estamos calculando el volumen es como en la figura \ref{fig:solido_completo}:

\begin{figure}[H]
  \centering
  \caption{Sólido completo}
  \label{fig:solido_completo}
  \includegraphics[width=0.30\textwidth]{Gráficas/SolRevArandEjeRecParab2Solido.pdf}
  \propia
\end{figure}

%\begin{center}----------
%\centerline{ \graf{Sólido completo}}
%\includegraphics[scale=0.45]{SolRevArandEjeRecParab2Solido.pdf}
%\textbf{Fuente:} elaboración propia
%\end{center}

En el siguiente enlace puede encontrarse la simulación del sólido y de las arandelas que se forman:

\begin{center}
  \verb"https://www.geogebra.org/m/qtmvn4pj"
\end{center}

%\begin{center}----------
%\centerline{ \graf{Región y arandelas del ejemplo}}
%\verb"https://www.geogebra.org/m/qtmvn4pj"
%\end{center}

\begin{ejemplo}
  Halle el volumen del sólido de revolución generado al girar la región limitada por $y=x\ln x$ y $y=4\ln x$ alrededor del eje $x$.
\end{ejemplo}

\emph{Solución.} Al igualar las expresiones, tenemos que los puntos de corte de las curvas se dan cuando $x=1$ y $x=4$. Luego, los puntos de corte son $(1,0)$ y $(4,4\ln 4)$:

\begin{figure}[H]
  \centering
  \caption{Región y arandela del ejemplo}
  \includegraphics[width=0.43\textwidth]{Gráficas/SolRevArandEjeLn.pdf}
  \propia
\end{figure}

%\begin{center}----------
%\centerline{ \graf{Arandela del ejemplo}}
%\includegraphics[scale=0.35]{SolRevArandEjeLn.pdf}
%\textbf{Fuente:} elaboración propia
%\end{center}

De manera explícita, la arandela se ve como en la figura \ref{fig:arandelas_D}:

\begin{figure}[H]
  \centering
  \caption{Arandela}
  \label{fig:arandelas_D}
  \includegraphics[width=0.32\textwidth]{Gráficas/SolRevArandEjeLnArandel.pdf}
  \propia
\end{figure}

%\begin{center}-----------
%\includegraphics[scale=0.35]{SolRevArandEjeLnArandel.pdf}
%\textbf{Fuente:} elaboración propia
%\end{center}

Apoyándonos en la figura, tenemos que $R=4\ln x$ y $r=x\ln x$. Entonces,
{\small
  \begin{align*}
    v&=\pi \int_{a}^{b}\left( R^{2}-r^{2}\right) dx
    =\pi \int_{1}^{4}\left( \left( 4\ln x\right) ^{2}-\left( x\ln x\right)
    ^{2}\right) dx\\
    &=\pi \int_{1}^{4}\left( 16\ln ^{2}x-x^{2}\ln ^{2}x\right) dx
  \end{align*}
}
Debemos calcular la integral para tener el volumen. Podemos usar integración por partes en este punto o una sustitución para que tengamos exponenciales en lugar de logaritmos. Usemos la \'ultima opción:
si $y=\ln x\rightarrow e^{y}=x\rightarrow e^{y}dy=dx$, y si $x=1\rightarrow y=\ln 1=0$, y si $x=4\rightarrow y=\ln 4$. Con ello,
{\small
  \begin{align*}
    \pi \int_{1}^{4}\left( 16\ln ^{2}x-x^{2}\ln ^{2}x\right) dx&=\pi \int_{0}^{\ln 4}\left( 16y^{2}-\left( e^{y}\right) ^{2}y^{2}\right)
    e^{y}dy \\
    &=\pi \left( 16\int_{0}^{\ln 4}y^{2}e^{y}dy-\int_{0}^{\ln
    4}y^{2}e^{3y}dy\right)
  \end{align*}
}

Y acá usamos integración por partes al estilo de la sección \ref{tabular} del capítulo \ref{antiderivada}. Tomemos aparte las dos integrales:

%\begin{center}
%\centerline{ \graf{Sólido del ejemplo}}
%\renewcommand{\arraystretch}{1.5}
\renewcommand{\arraystretch}{1.5}
\begin{center}
  \begin{tabular}{>{$}c<{$} >{$}c<{$} | >{$}c<{$} >{$}c<{$}}
    \multicolumn{2}{c|}{$\int_{0}^{\ln 4}y^{2}e^{y}dy$} & \multicolumn{2}{c}{$\int_{0}^{\ln 4}y^{2}e^{3y}dy$}\\
    \hline
    \hline
    y^{2} & e^{y} & y^{2} & e^{3y}\\
    \hline
    2y & e^{y} & 2y & \frac{1}{3}e^{3y}\\
    \hline
    2 & e^{y} & 2 & \frac{1}{9}e^{3y}\\
    \hline
    0 & e^{y} & 0 & \frac{1}{27}e^{3y}\\
    \hline
    \multicolumn{2}{c|}{$=e^{y}\left( y^{2}-2y+2\right)$} & \multicolumn{2}{c}{$=e^{3y}\left( \frac{1}{3}y^{2}-\frac{2}{9}y+\frac{2}{27}\right)$}\\
    \hline
  \end{tabular}
\end{center}

%\end{center}

En este sentido,
{\small
  \begin{align*}
    \pi &\left( 16\int_{0}^{\ln 4}y^{2}e^{y}dy-\int_{0}^{\ln
    4}y^{2}e^{3y}dy\right)=16\pi \int_{0}^{\ln 4}y^{2}e^{y}dy-\pi \int_{0}^{\ln
    4}y^{2}e^{3y}dy\\
    &=16\pi \left( e^{y}\left( y^{2}-2y+2\right) \right)\Big\vert _{0}^{\ln 4}
    -\pi e^{3y}\left( \frac{1}{3}y^{2}-\frac{2}{9}y+\frac{2}{27}\right)\Bigg\vert_{0}^{\ln 4}\\
    &=16\pi \left( 4\left( \ln ^{2}4-2\ln 4+2\right) -2\right)-64\pi\left( \frac{1}{3}\ln ^{2}4-\frac{2}{9}\ln 4+\frac{2}{27}\right) -\frac{2\pi}{27} \\
    &=\pi \left( \frac{512}{3}\ln ^{2}2-\frac{2048}{9}\ln 2+\frac{274}{3}\right)
    \approx49.012\text{ unidades c\'ubicas.}%
  \end{align*}
}

El sólido al que le estamos calculando el volumen se ve en la figura \ref{fig:region_arandela}:

\begin{figure}[H]
  \centering
  \caption{Sólido del ejemplo}
  \label{fig:region_arandela}
  \includegraphics[width=0.25\textwidth]{Gráficas/SolRevArandEjeLnSolido.pdf}
  \propia
\end{figure}

%\begin{center}----------
%\centerline{ \graf{Región y arandelas del ejemplo}}
%\includegraphics[scale=0.35]{SolRevArandEjeLnSolido.pdf}
%\textbf{Fuente:} elaboración propia
%\end{center}

En el siguiente enlace puede encontrarse la simulación del sólido y de las arandelas que se forman:

\begin{center}
  \verb"https://www.geogebra.org/m/eq3cbtjt"
\end{center}

\begin{ejemplo}
  Halle el volumen del sólido de revolución generado al girar la región limitada por $y=\cfrac{2x^{2}}{x^{2}+1}$ y $y=x^{2}$ alrededor del eje $x$.
\end{ejemplo}
\emph{Solución.} La información suministrada se resume en la siguiente figura \ref{fig:arandelas_arandela}:

\begin{figure}[H]
  \centering
  \caption{Región y arandela del ejemplo}
  \label{fig:arandelas_arandela}
  \includegraphics[width=0.48\textwidth]{Gráficas/SolRevArandEjeRacion.pdf}
  \propia
\end{figure}

%\begin{center}----------
%\centerline{ \graf{Arandela del ejemplo}}
%\includegraphics[scale=0.35]{SolRevArandEjeRacion.pdf}
%\textbf{Fuente:} elaboración propia
%\end{center}

En tres dimensiones, la arandela puede verse como en la figura \ref{fig:arandelas_del_ejemp_}:

\begin{figure}[H]
  \centering
  \caption{Arandela del ejemplo}
  \label{fig:arandelas_del_ejemp_}
  \includegraphics[width=0.33\textwidth]{Gráficas/SolRevArandEjeRacionArandel.pdf}
  \propia
\end{figure}

%\begin{center}-----------
%\includegraphics[scale=0.35]{SolRevArandEjeRacionArandel.pdf}
%\textbf{Fuente:} elaboración propia
%\end{center}

Con la información de la figura, tenemos

\resizebox{\textwidth}{!}{$
  \begin{aligned}
    v&=\pi \int_{-1}^{1}\left( R^{2}-r^{2}\right) dx
    =\pi \int_{-1}^{1}\underbrace{\left[\left( \frac{%
    2x^{2}}{x^{2}+1}\right) ^{2}-\left( x^{2}\right) ^{2} \right]}_\text{función par}dx=2\pi \int_{0}^{1}\left[ \left( \frac{2x^{2}}{x^{2}+1}\right)
    ^{2}-x^{4}\right] dx \\
    &
    =2\pi \left( \int_{0}^{1}\left( \frac{2x^{2}}{x^{2}+1}\right) ^{2}dx-\left(
    \frac{1}{5}x^{5}\right)\Bigg\vert_{0}^{1}\right)
    =2\pi \left( 5-\frac{3}{2}\pi -\frac{1}{5}\right)
    =2\pi \left( \frac{24}{5}-\frac{3}{2}\pi \right)\\
    &\approx0.55048\text{ unidades c\'ubicas.}%
  \end{aligned}
$}

La primera integral se calculó usando sustitución trigonométrica:
{\small
  \begin{align*}
    \int_{0}^{1}&\left( \frac{2x^{2}}{x^{2}+1}\right) ^{2}dx=4\int_{0}^{1}\frac{x^{4}}{\left( x^{2}+1\right) ^{2}}dx%\\
    %&
    =4\int_{0}^{\frac{\pi }{4}}\frac{\tan ^{4}\theta \sec ^{2}\theta d\theta }{
    \sec ^{4}\theta }\\
    &=4\int_{0}^{\frac{\pi }{4}}\frac{\left( \sec ^{2}\theta
    -1\right) ^{2}d\theta }{\sec ^{2}\theta } %\\
    %&
    =4\int_{0}^{\frac{\pi }{4}}\frac{\sec ^{4}\theta -2\sec ^{2}\theta +1}{\sec
    ^{2}\theta }d\theta \\
    &=4\int_{0}^{\frac{\pi }{4}}\left( \sec ^{2}\theta -2+\cos ^{2}\theta \right)
    d\theta %\\
    %&
    =4\int_{0}^{\frac{\pi }{4}}\left( \sec ^{2}\theta -2+\frac{1}{2}+\frac{1}{2}%
    \cos \left( 2\theta \right) \right)d\theta \\
    &=4\left( \tan\theta -\frac{3}{2}\theta +\frac{1}{4}\sen \left( 2\theta
    \right) \right)\Bigg\vert _{0}^{\frac{\pi }{4}} %\\
    %&
    =4\left( \tan \frac{\pi }{4}-\frac{3}{2}\frac{\pi }{4}+\frac{1}{4}\sen
    \left( \frac{\pi }{2}\right) \right) \\
    &=4-\frac{3\pi }{2}+1=5-\frac{3\pi}{2}
  \end{align*}
}
El sólido completo se ve adelante:

\begin{figure}[H]
  \centering
  \caption{Sólido completo del ejemplo}
  \includegraphics[width=0.35\textwidth]{Gráficas/SolRevArandEjeRacionSolido.pdf}
  \propia
\end{figure}

%\begin{center}----------
%\centerline{ \graf{Sólido completo del ejemplo}}
%\includegraphics[scale=0.35]{SolRevArandEjeRacionSolido.pdf}
%\textbf{Fuente:} elaboración propia
%\end{center}

En el siguiente enlace puede encontrarse la simulación del sólido y de las arandelas que se forman:

\begin{center}
  \verb"https://www.geogebra.org/m/zfytfxrj"
\end{center}

%\begin{center}----------
%\centerline{ \graf{La región del ejemplo}}
%\verb"https://www.geogebra.org/m/zfytfxrj"
%\end{center}

\begin{ejemplo}
  Determine el volumen del sólido de revolución que se tiene al girar alrededor del eje $x$ la región limitada por $y=9-x^{2}$ y \textit{$y=x+3$.}
\end{ejemplo}
%%\begin{multicols}{2}
\emph{Solución.} Los puntos de corte de las curvas se tienen al solucionar $9-x^{2}=x+3$, es decir, $x^2+x-6=(x+3)(x-2)=0$, o, mejor, $x=-3$ y $x=2$; así, los puntos de corte son {\small $(-3,0)$} y {\small $(2,5)$}. La región que gira se ve de la siguiente manera:

\begin{figure}[H]
  \centering
  \caption{Región a girar}
  \includegraphics[width=0.28\textwidth]{Gráficas/SolRevArandParabRecta.pdf}
  \propia
\end{figure}

%\begin{center}---------
%\centerline{ \graf{Radio mayor y radio menor}}
%\includegraphics[scale=0.35]{SolRevArandParabRecta.pdf}
%\textbf{Fuente:} elaboración propia
%\end{center}

En la figura \ref{fig:alrededor_eje} puede identificarse el radio mayor y radio menor al girar la región anterior alrededor del eje $x$:

\begin{figure}[H]
  \centering
  \caption{Radios mayor y menor}
  \label{fig:alrededor_eje}
  \includegraphics[width=0.28\textwidth]{Gráficas/SolRevArandParabRecta2.pdf}
  \propia
\end{figure}

%\begin{center}
%\includegraphics[scale=0.35]{SolRevArandParabRecta2.pdf}
%\textbf{Fuente:} elaboración propia
%\end{center}
%
La arandela formada se ve en la figura \ref{fig:radio_mayor_menor_}:

\begin{figure}[H]
  \centering
  \caption{Arandela}
  \label{fig:radio_mayor_menor_}
  \includegraphics[width=0.20\textwidth]{Gráficas/SolRevArandParabRecta3.pdf}
  \propia
\end{figure}

%\begin{center}-----
%\centerline{ \graf{Radio mayor y radio menor}}
%\includegraphics[scale=0.35]{SolRevArandParabRecta3.pdf}
%\textbf{Fuente:} elaboración propia
%\end{center}

El sólido completo se ve como en la figura \ref{fig:solido_ejem}:

\begin{figure}[H]
  \centering
  \caption{El sólido del ejemplo}
  \label{fig:solido_ejem}
  \includegraphics[width=0.30\textwidth]{Gráficas/SolRevArandEjeParabRect.pdf}
  \propia
\end{figure}

%\begin{center}---------
%\centerline{ \graf{El sólido del ejemplo}}
%\includegraphics[scale=0.35]{SolRevArandEjeParabRect.pdf}
%\textbf{Fuente:} elaboración propia
%\end{center}

En el siguiente enlace puede encontrarse la simulación del sólido y de las arandelas que se forman:

\begin{center}
  \verb"https://www.geogebra.org/m/nhfsje2v"
\end{center}

%\begin{center}-----------
%\centerline{ \graf{La región}} %este nombre va ahí??
%\verb"https://www.geogebra.org/m/nhfsje2v"
%\end{center}

Con ello, puede calcularse su volumen, que corresponde a $\frac{625}{3}\pi $ unidades c\'ubicas, como lo muestran los siguientes cálculos:
%%\end{multicols}

\resizebox{\textwidth}{!}{$
  \begin{aligned}
    v&=\int_{-3}^{2}\pi\left[ \left( 9-x^{2}\right) ^{2}-\left( x+3\right)
    ^{2}\right] dx
    %\\ &
    =\pi\int_{-3}^{2}\left( x^{4}-19x^{2}-6x+72\right) dx\\
    &=\pi\left(
    72x-3x^{2}-\frac{19}{3}x^{3}+\frac{1}{5}x^{5}\right)\Bigg\vert_{-3}^{2}\\
    &=\pi \left( 72\left( 2\right) -3\left( 2\right) ^{2}-\frac{19}{3}\left(
      2\right) ^{3}+\frac{1}{5}\left( 2\right) ^{5}-\left( 72\left( -3\right)
        -3\left( -3\right) ^{2}-\frac{19}{3}\left( -3\right) ^{3}+\frac{1}{5}\left(
    -3\right) ^{5}\right) \right)\\
    &=\pi \left( 144-12-\frac{152}{3}+\frac{32}{5}-\left( -216-27+171-\frac{243}{%
    5}\right) \right) \\
    &=\pi \left( 132-\frac{152}{3}+\frac{32}{5}+72+\frac{243}{5}=204-\frac{152}{3%
    }+\frac{275}{5}\right) \\
    &=\pi \left( \frac{204\left( 15\right) -152\left( 5\right) +275\left(
    3\right) }{15}\right) \\
    &=\frac{3125}{15}\pi \\
    &=\frac{625}{3}\pi \text{ unidades c\'ubicas.}%
  \end{aligned}
$}

\begin{ejemplo}
  Halle el volumen del sólido de revolución obtenido al girar la
  región acotada por $y=x^{2}$ y $y=2x+3$ alrededor del eje $x$.
\end{ejemplo}
\emph{Solución.} Los puntos de corte de las curvas se encuentran haciendo $x^{2}=2x+3$, es decir, $x^{2}-2x-3=0$, o también $\left( x+1\right)\left( x-3\right) =0$; por tanto, $x=-1$ y $x=3$; al reemplazar esos valores en cualquiera de las expresiones igualadas, tenemos el valor de $y$, es decir, los puntos de corte son $(-1,1)$ y $(3,9)$. La figura \ref{fig:rotar_region} nos muestra la región que encierran las curvas:
%%\begin{multicols}{2}

\begin{figure}[H]
  \centering
  \caption{Región de giro}
  \label{fig:rotar_region}
  \includegraphics[width=0.28\textwidth]{Gráficas/SolRevArandEjeRecParabol.pdf}
  \propia
\end{figure}

%\begin{center}-----------
%\centerline{ \graf{Al rotar la región}}
%\includegraphics[scale=0.4]{SolRevArandEjeRecParabol.pdf}
%\textbf{Fuente:} elaboración propia
%\end{center}

Al rotar la región, pueden identificarse los radios como en la siguiente figura \ref{fig:identifi_radio}:

\begin{figure}[H]
  \centering
  \caption{Identificación de radios}
  \label{fig:identifi_radio}
  \includegraphics[width=0.29\textwidth]{Gráficas/SolRevArandEjeRecParabol2.pdf}
  \propia
\end{figure}

%\begin{center}------------
%\includegraphics[scale=0.4]{SolRevArandEjeRecParabol2.pdf}
%\textbf{Fuente:} elaboración propia
%\end{center}

La arandela que se forma al girar la región se ve como en la figura \ref{fig:arandela_formada}:

\begin{figure}[H]
  \centering
  \caption{Arandela formada}
  \label{fig:arandela_formada}
  \includegraphics[width=0.28\textwidth]{Gráficas/SolRevArandEjeRecParabol3.pdf}
  \propia
\end{figure}

%\begin{center}-----------
%\centerline{ \graf{Arandela formada}}
%\includegraphics[scale=0.45]{SolRevArandEjeRecParabol3.pdf}
%\textbf{Fuente:} elaboración propia
%\end{center}

El sólido es el siguiente:

\begin{figure}[H]
  \centering
  \caption{El sólido}
  \includegraphics[width=0.28\textwidth]{Gráficas/SolRevArandEjeRecParabol4.pdf}
  \propia
\end{figure}

%\begin{center}----------
%\centerline{ \graf{El sólido}}
%\includegraphics[scale=0.45]{SolRevArandEjeRecParabol4.pdf}
%\textbf{Fuente:} elaboración propia
%\end{center}

Y en el siguiente enlace puede encontrarse la simulación del sólido y de las arandelas que se forman:

\begin{center}
  \verb"https://www.geogebra.org/m/pkwxxxv3"
\end{center}

%\begin{center}-----------
%\centerline{ \graf{La región}}
%\verb"https://www.geogebra.org/m/pkwxxxv3"
%\end{center}

%%\end{multicols}
Los siguientes cálculos encuentran el volumen del sólido:

\resizebox{\textwidth}{!}{$
  \begin{aligned}
    v&=\pi \int_{-1}^{3}\left( \left( 2x+3\right)^{2}-\left( x^{2}\right)
    ^{2}\right) dx %\\
    %&
    =\pi \int_{-1}^{3}\left( 4x^{2}+12x+9-x^{4}\right) dx\\
    &=\pi \left( 9x+6x^{2}+\frac{4}{3}x^{3}-\frac{1}{5}x^{5}\right)\Bigg\vert_{-1}^{3}\\
    &=\pi \left( 9\left( 3\right) +6\left( 3\right) ^{2}+\frac{4}{3}\left(
      3\right) ^{3}-\frac{1}{5}\left( 3\right) ^{5}-\left( 9\left( -1\right)
        +6\left( -1\right) ^{2}+\frac{4}{3}\left( -1\right) ^{3}-\frac{1}{5}\left(
    -1\right) ^{5}\right) \right) \\
    &=\pi \left( 27+54+36-\frac{243}{5}+9-6+\frac{4}{3}-\frac{1}{5}\right) \\
    &=\pi \left( 120+\frac{4}{3}-\frac{244}{5}\right) \\
    &=\pi \left( \frac{120\left( 15\right) +20-732}{15}\right) \\
    &=\frac{1088}{15}\pi\text{ unidades c\'ubicas.}%
  \end{aligned}
$}

\begin{ejemplo}
  Halle el volumen del sólido de revolución obtenido al girar alrededor del eje $x$ la región acotada por las figuras de $y=4xe^{x}$ y \textit{$y=x^{2}e^{x}$.}
\end{ejemplo}

%%\begin{multicols}{2}

\emph{Solución.} La figura de la región del ejemplo es la siguiente:

\begin{figure}[H]
  \centering
  \caption{Región a girar}
  \includegraphics[width=0.25\textwidth]{Gráficas/SolRevArandExponen.pdf}
  \propia
\end{figure}

%\begin{center}----------
%\centerline{ \graf{Los radios}}
%\includegraphics[scale=0.35]{SolRevArandExponen.pdf}
%\textbf{Fuente:} elaboración propia
%\end{center}

Al girar alrededor del eje $x$ podemos ver esto:

\begin{figure}[H]
  \centering
  \caption{Los radios}
  \includegraphics[width=0.29\textwidth]{Gráficas/SolRevArandExponen2.pdf}
  \propia
\end{figure}

%\begin{center}-----------
%\includegraphics[scale=0.35]{SolRevArandExponen2.pdf}
%\textbf{Fuente:} elaboración propia
%\end{center}

La arandela puede identificarse en la figura \ref{fig:la_arandela}:

\begin{figure}[H]
  \centering
  \caption{La arandela}
  \label{fig:la_arandela}
  \includegraphics[width=0.24\textwidth]{Gráficas/SolRevArandExponen3.pdf}
  \propia
\end{figure}

%\begin{center}------------
%\centerline{ \graf{La arandela}}
%\includegraphics[scale=0.35]{SolRevArandExponen3.pdf}
%\textbf{Fuente:} elaboración propia
%\end{center}

Precisemos los puntos de corte: si $x^{2}e^{x}=4xe^{x}$, entonces $x(x-4)e^{x}=0$, pero $e^x$ nunca se hace cero; por tanto, no le aporta soluciones a la igualdad; entonces, $x(x-4)=0$ y, así, $x=0$ y $x=4.$ Con ello, la integral (con arandelas) que representa el volumen del sólido es
{\small
  \begin{align*}
    v&=\pi \int_{0}^{4}\left[\left( 4xe^{x}\right)^{2}-\left(x^{2}e^{x}\right)^{2}\right]\, dx\\
    &=\pi \int_{0}^{4}\left(16x^{2}e^{2x}-x^{4}e^{2x}\right)\, dx \\
    &=\pi \int_{0}^{4}\left(16x^{2}-x^{4}\right) e^{2x}dx
  \end{align*}
}
Esta es una integral que podemos solucionar tabulando, como se hizo en la sección \ref{tabular} del capítulo \ref{antiderivada}:
\renewcommand{\arraystretch}{1.5}
%\begin{center}------------
%\centerline{ \graf{El sólido}}

\begin{center}
  \renewcommand{\arraystretch}{1.2} % Mejora el espaciado vertical
  \begin{tabular}{|c|c|c|c|}
    \hline
    \textbf{Funciones} & $16x^{2}-x^{4}$ & $e^{2x}$ & \textbf{Operación}\\
    \hline
    \text{Derivada} $\rightarrow$ & $32x-4x^{3}$ & $\frac{1}{2}e^{2x}$ & $\leftarrow$ \text{Integral} \\
    \hline
    \text{Derivada} $\rightarrow$ & $32-12x^{2}$ & $\frac{1}{4}e^{2x}$ & $\leftarrow$ \text{Integral} \\
    \hline
    \text{Derivada} $\rightarrow$ & $-24x$ & $\frac{1}{8}e^{2x}$ & $\leftarrow$ \text{Integral} \\
    \hline
    \text{Derivada} $\rightarrow$ & $-24$ & $\frac{1}{16}e^{2x}$ & $\leftarrow$ \text{Integral} \\
    \hline
    \text{Derivada} $\rightarrow$ & $0$ & $\frac{1}{32}e^{2x}$ & $\leftarrow$ \text{Integral} \\
    \hline
  \end{tabular}
\end{center}

%\end{center}

Y así, el volumen del sólido de la figura \ref{fig:region_y_solido_del_ejemplo}:

\begin{figure}[H]
  \centering
  \caption{El sólido del ejemplo}
  \label{fig:region_y_solido_del_ejemplo}
  \includegraphics[width=0.21\textwidth]{Gráficas/SolRevArandExponen4.pdf}
  \propia
\end{figure}

%\begin{center}---------
%\centerline{ \graf{La región y el sólido del ejemplo}}
%\includegraphics[scale=0.35]{SolRevArandExponen4.pdf}
%\textbf{Fuente:} elaboración propia
%\end{center}

está dado por
%%\end{multicols}

\resizebox{\textwidth}{!}{$
  \begin{aligned}
    v&=\left. \pi \frac{1}{2}e^{2x}\left( 16x^{2}-x^{4}-\frac{1}{2}\left(
      32x-4x^{3}\right) +\frac{1}{4}\left( 32-12x^{2}\right) -\frac{1}{8}\left(
    -24x\right) +\frac{1}{16}\left( -24\right) \right) \right\vert _{0}^{4}\\
    &=\left. \frac{\pi e^{2x}}{4}\left( -2x^{4}+4x^{3}+26x^{2}-26x+13\right)
    \right\vert _{0}^{4}\\
    &=\frac{\pi }{4}\left( e^{8}\left( -2\times 4^{4}+4\times 4^{3}+26\times
    4^{2}-26\times 4+13\right) -13\right) \\
    &=\frac{\pi }{4}\left( 69e^{8}-13\right)\text{ unidades c\'ubicas.}%
  \end{aligned}
$}

En el siguiente enlace puede encontrarse la simulación del sólido y de las arandelas que se forman:

\begin{center}
  \verb"https://www.geogebra.org/m/pvb2a4cs"
\end{center}

%\begin{center}---------
%\verb"https://www.geogebra.org/m/pvb2a4cs"
%\end{center}

%\subsection{Método de capas cilíndricas\footnote{Este método, en lugar de \textit{capas cilíndricas}, suele llamarse también de \textit{casquetes} o \textit{cascarones cilíndricos}.}}

\subsection*{Método de capas cilíndricas\protect\footnote{Este método, en lugar de \emph{capas cilíndricas}, suele llamarse también de \emph{casquetes} o \emph{cascarones cilíndricos}.}}

Con este método tenemos una herramienta adicional a las dos anteriores para calcular el volumen de un sólido de revolución. Tengamos en cuenta que no todos se usan necesariamente para la misma situación. Ilustremos la necesidad de este método por medio de un ejemplo:

\begin{ejemplo}
  Halle el volumen del sólido de revolución que se genera al
  girar alrededor del eje $y$, la región limitada por las figuras de $y=\frac{1}{16}x^{5}+x+1$, el eje $x$, el eje $y$ y la recta $x=2$.
\end{ejemplo}

\emph{Solución.} La región y el sólido que describen el contexto son los siguientes:

\begin{figure}[H]
  \centering
  \caption{Región y sólido del ejemplo}
  \includegraphics[width=0.42\textwidth]{Gráficas/CasquillosGradCinco.pdf}
  \propia
\end{figure}

%\begin{center}------------
%\centerline{ \graf{Casquillo del ejemplo}}
%\includegraphics[scale=0.35]{CasquillosGradCinco.pdf}
%\textbf{Fuente:} elaboración propia
%\end{center}

%%%\begin{multicols}{2}
El casquillo solo puede verse en la figura \ref{fig:casquillo_ejemplo}:

\begin{figure}[H]
  \centering
  \caption{Casquillo del ejemplo}
  \label{fig:casquillo_ejemplo}
  \includegraphics[width=0.41\textwidth]{Gráficas/CasquillosGradCincoCasquill.pdf}
  \propia
\end{figure}

%\begin{center}-----------
%\includegraphics[scale=0.45]{CasquillosGradCincoCasquill.pdf}
%\textbf{Fuente:} elaboración propia
%\end{center}

O también podemos ver el sólido de la siguiente manera:

\begin{figure}[H]
  \centering
  \caption{El sólido del ejemplo en 3D}
  \includegraphics[width=0.32\textwidth]{Gráficas/CasquillosGradCinco3D.pdf}
  \propia
\end{figure}

%\begin{center}-----------
%\centerline{ \graf{El sólido del ejemplo en 3D}}
%\includegraphics[scale=0.35]{CasquillosGradCinco3D.pdf}
%\textbf{Fuente:} elaboración propia
%\end{center}

%%%\end{multicols}

Notemos que, para poder calcular el volumen del sólido por el método de las secciones transversales perpendiculares al eje $y$, que en este caso son arandelas, tendríamos que hallar la distancia desde $\left( 0,y\right)$ hasta {\small \[\left( x,\frac{1}{16}x^{5}+x+1\right)\]} para un n\'umero $y\in \left[ 1,5\right] $. Para ello, debería escribirse $x$ en términos de $y$ en la ecuación $y=\frac{1}{16}x^{5}+x+1,$ lo cual es imposible en términos de radicales. Es decir, el método de las arandelas no puede aplicarse en este caso particular. Se requiere, por tanto, un nuevo método para medir el volumen de este tipo de sólidos.

Con el siguiente enlace pueden simularse el sólido y los casquillos en el ejemplo:

\begin{center}
  \verb"https://www.geogebra.org/m/xsk7h7bz"
\end{center}

%\begin{center}-----------
%\centerline{ \graf{Región y eje de rotación}}
%\verb"https://www.geogebra.org/m/xsk7h7bz"
%\end{center}

Veamos, entonces, el nuevo método. Sea $S$ el sólido obtenido al girar alrededor del eje $y$ la región acotada por las figuras de $y=f\left( x\right) $ (donde $f\left( x\right) \geq 0$), $\ y=0$, $x=a$, y $x=b$, con $b>a\geq 0$).

Con ello, el volumen del sólido \textit{S} es

\begin{equation}\label{casquetes}
V\left( S\right) =\int_{a}^{b}2\pi xf\left( x\right) dx\text{ donde }0\leq
a\leq b
\end{equation}

Demos una explicación del funcionamiento de la fórmula \eqref{casquetes}: para cada $x$ en el intervalo $\left[ a,b\right] $, supongamos que $C_{x}$ es la capa cilíndrica generada al girar alrededor del eje\textbf{ $y$} el segmento que une los puntos $(x,0)$ con $\left( x,f\left( x\right) \right) $. Así que el s%
ólido \textit{S} es $S=\cup _{x\in \left[ a,b\right] }C_{x}.$

Recortando $C_{x}$ por el segmento generador, el volumen $V\left(
C_{x}\right) $ de la capa cilíndrica C$_{x}$ es
\begin{align*}
V\left( C_{x}\right) &=2\pi (radio)(altura)\left( espesor\right)\\
& =2\pi xf\left( x\right) dx
\end{align*}
y el volumen del sólido es
\begin{align*}
V\left( S\right) &=V\left( \cup_{x\in \left[ a,b\right] }C_{x}\right)
\\
&=\sum_{x\in \left[ a,b\right] }V\left( C_{x}\right) \\
&= \int_{a}^{b}2\pi
(radio)(altura)\left( espesor\right) \\
&=\int_{a}^{b}2\pi xf\left( x\right) dx
\end{align*}
Tenga en cuenta que, si el eje de rotación es diferente a uno de los ejes coordenados, entonces el radio es una función de $x$ en lugar de $x$.

\begin{ejemplo}
Halle el volumen del sólido de revolución que se tiene al girar la
región limitada por las figuras de $y=x^{4}+x $, el eje $x$ y la recta $x=1$, en torno a la recta $x=3$.
\end{ejemplo}

\emph{Solución.} La figura \ref{fig:radio_altura_casquillo} muestra la región y el eje de rotación del ejemplo:

\begin{figure}[H]
\centering
\caption{Región y eje de rotación}
\label{fig:radio_altura_casquillo}
\includegraphics[width=0.34\textwidth]{Gráficas/CasquillAlredx3.pdf}

\propia
\end{figure}

%\begin{center}-----------
%\centerline{ \graf{Radio y altura del casquillo}}
%\includegraphics[scale=0.35]{CasquillAlredx3.pdf}
%\textbf{Fuente:} elaboración propia
%\end{center}

Al girar un rectángulo, tenemos

\begin{figure}[H]
\centering
\caption{Radio y altura del casquillo}
\label{fig:radio_casquillo}
\includegraphics[width=0.34\textwidth]{Gráficas/CasquillAlredx3_2.pdf}

\propia
\end{figure}

%\begin{center}-----------
%\includegraphics[scale=0.35]{CasquillAlredx3_2.pdf}
%\textbf{Fuente:} elaboración propia
%\end{center}

y el casquillo solo lo podemos ver como

\begin{figure}[H]
\centering
\caption{Casquillo}
\includegraphics[width=0.34\textwidth]{Gráficas/CasquillAlredx3_3.pdf}

\propia
\end{figure}

%\begin{center}-----------
%\centerline{ \graf{Casquillo}}
%\includegraphics[scale=0.35]{CasquillAlredx3_3.pdf}
%\textbf{Fuente:} elaboración propia
%\end{center}

y el sólido completo es

\begin{figure}[H]
\centering
\caption{Sólido}
\includegraphics[width=0.34\textwidth]{Gráficas/CasquillAlredx3Solid.pdf}

\propia
\end{figure}

%\begin{center}------------
%\centerline{ \graf{Sólido}}
%\includegraphics[scale=0.35]{CasquillAlredx3Solid.pdf}
%\textbf{Fuente:} elaboración propia
%\end{center}

Con la información de la figura, podemos confirmar que el radio del casquillo cilíndrico es $3-x$ y su altura, $x^{4}+x$, entonces

\resizebox{\textwidth}{!}{$
\begin{aligned}
  V&=2\pi \int_{a}^{b}(radio)(altura)\, dx=2\pi \int_{0}^{1}\left( 3-x\right) \left(x^{4}+x\right) dx \\
  &=2\pi \int_{0}^{1}\left( -x^{5}+3x^{4}-x^{2}+3x\right) dx %\\
  =2\pi\left. \left(-\frac{1}{6}x^{6}+\frac{3}{5}x^{5}-\frac{1}{3}x^{3}+\frac{3}{2}%
  x^{2}\right)\right\vert _{0}^{1} \\
  &=2\pi \left( -\frac{1}{6}+\frac{3}{5}-\frac{1}{3}+\frac{3}{2}\right) =\frac{16}{5}\pi \approx10.053 \text{ unidades c\'ubicas.}%
\end{aligned}
$}

Con el siguiente enlace pueden simularse el sólido y los casquillos en el ejemplo:

\begin{center}
\verb"https://www.geogebra.org/m/c6tw4pm6"
\end{center}

%\begin{center}-----------
%\centerline{ \graf{Región que gira}}
%\verb"https://www.geogebra.org/m/c6tw4pm6"
%\end{center}

\begin{ejemplo}
Encuentre el volumen del sólido de revolución obtenido por el giro de la región limitada por las figuras de $y=x^{4}+x$, $x=3$ y el eje $x$ en torno al eje \textit{$y$.}
\end{ejemplo}

%%\begin{multicols}{2}

\emph{Solución.} En la figura se tiene la región que gira:
%\clearpage

\begin{figure}[H]
\centering
\caption{Región a girar}
\includegraphics[width=0.32\textwidth]{Gráficas/CasquiAlredEjey.pdf}

\propia
\end{figure}

%\begin{center}------------
%\centerline{ \graf{Sólido}}
%\includegraphics[scale=0.35]{CasquiAlredEjey.pdf}
%\textbf{Fuente:} elaboración propia
%\end{center}

El sólido se ve como en la figura \ref{fig:solido_casquillo_1}:

\begin{figure}[H]
\centering
\caption{Sólido}
\label{fig:solido_casquillo_1}
\includegraphics[width=0.30\textwidth]{Gráficas/CasquillAlredyx4Solid.pdf}

\propia
\end{figure}

%\begin{center}-----------
%\includegraphics[scale=0.5]{CasquillAlredyx4Solid.pdf}
%\textbf{Fuente:} elaboración propia
%\end{center}

%%\end{multicols}

en donde se forma el casquillo:

\begin{figure}[H]
\centering
\caption{Casquillo}
\includegraphics[width=0.29\textwidth]{Gráficas/CasquillAlredyx4_1.pdf}

\propia
\end{figure}

%\begin{center}------------
%\centerline{ \graf{Casquillo}}
%\includegraphics[scale=0.5]{CasquillAlredyx4_1.pdf}
%\textbf{Fuente:} elaboración propia
%\end{center}

Notemos que en este caso el radio de la capa cilíndrica es $x$; por tanto, su volumen es
{\small
\begin{align*}
  V&=2\pi \int_{a}^{b}(radio)(altura)\, dx=2\pi \int_{0}^{3}x\left( x^{4}+x\right) dx=2\pi \left( \frac{1}{3}x^{3}+\frac{1}{6}x^{6}\right) _{0}^{3}\\
  &=2\pi \left( 9+\frac{243}{2}\right)=261\pi\text{ unidades c\'ubicas.}%
\end{align*}
}

Con el siguiente enlace pueden simularse el sólido y los casquillos en el ejemplo:

\begin{center}
\verb"https://www.geogebra.org/m/cng7rxgc"
\end{center}

%\begin{center}-----------
%\centerline{ \graf{Región que gira}}
%\verb"https://www.geogebra.org/m/cng7rxgc"
%\end{center}

\begin{ejemplo}
Encuentre el volumen del sólido que se tiene al hacer girar la región plana limitada por la curva $x=y^{4}-2y+1$ y la recta $x=1,$ alrededor del eje $x.$
\end{ejemplo}
%%\begin{multicols}{2}

\emph{Solución.} La región que se gira se muestra en la figura \ref{fig:solido_ejemploo_}:

\begin{figure}[H]
\centering
\caption{Región a girar}
\label{fig:solido_ejemploo_}
\includegraphics[width=0.29\textwidth]{Gráficas/CasquiAlredEjey4.pdf}

\propia
\end{figure}

%\begin{center}-----------
%\centerline{ \graf{El sólido}}
%\includegraphics[scale=0.4]{CasquiAlredEjey4.pdf}
%\textbf{Fuente:} elaboración propia
%\end{center}

El sólido se ve como en la figura \ref{fig:solido_casquillo_2}:

\begin{figure}[H]
\centering
\caption{El sólido}
\label{fig:solido_casquillo_2}
\includegraphics[width=0.27\textwidth]{Gráficas/CasquiAlredEjey4Solid.pdf}

\propia
\end{figure}

%\begin{center}-----------
%\includegraphics[scale=0.35]{CasquiAlredEjey4Solid.pdf}
%\textbf{Fuente:} elaboración propia
%\end{center}

%%\end{multicols}
%\clearpage

Y en la figura \ref{fig:casquillo_cilindro} se muestra el casquillo cilíndrico generado:

\begin{figure}[H]
\centering
\caption{El casquillo cilíndrico}
\label{fig:casquillo_cilindro}
\includegraphics[width=0.28\textwidth]{Gráficas/CasquiAlredEjey4Casqui.pdf}

\propia
\end{figure}

%\begin{center}-----------
%\centerline{ \graf{El casquillo cilíndrico}}
%\includegraphics[scale=0.4]{CasquiAlredEjey4Casqui.pdf}
%\textbf{Fuente:} elaboración propia
%\end{center}

Con ello, el volumen está dado por la integral:
{\small
\begin{eqnarray*}
  V &=&\int_{0}^{\sqrt[3]{2}}2\pi \left( 1-y^{4}+2y-1\right) ydy=\int_{0}^{%
  \sqrt[3]{2}}2\pi \left( -y^{4}+2y\right) ydy \\
  &=&2\pi \int_{0}^{\sqrt[3]{2}}\left( -y^{5}+2y^{2}\right) dy=2\pi \left( -%
  \frac{\sqrt[3]{2}^{\, 6}}{6}+2\frac{\sqrt[3]{2}^{\, 3}}{3}\right) \\
  &=&2\pi \left( -\frac{4}{6}+\frac{4}{3}\right) =2\pi \left( \frac{4}{6}%
  \right) =\frac{4\pi }{3}\text{ unidades c\'ubicas.}%
\end{eqnarray*}
}
Con el siguiente enlace pueden simularse el sólido y los casquillos en el ejemplo:

\begin{center}
\verb"https://www.geogebra.org/m/bjm64xtr"
\end{center}

%\begin{center}-----------
%\centerline{ \graf{Región que gira}}
%\verb"https://www.geogebra.org/m/bjm64xtr"
%\end{center}

\begin{ejemplo}
Encuentre el volumen del sólido que se tiene al hacer girar la región plana limitada por la curva $x=y^{3}+3y+4,$ la recta $y=1 $ y el eje $y$ alrededor de la recta $y=-2.$
\end{ejemplo}
%%\begin{multicols}{2}

\emph{Solución.} La figura \ref{fig:region_figura} muestra la región del plano que se está girando:

\begin{figure}[H]
\centering
\caption{Región a girar}
\label{fig:region_figura}
\includegraphics[width=0.35\textwidth]{Gráficas/CasquiAlredEjey2.pdf}

\propia
\end{figure}

%\begin{center}----------
%\centerline{ \graf{El sólido}}
%\includegraphics[scale=0.35]{CasquiAlredEjey2.pdf}
%\textbf{Fuente:} elaboración propia
%\end{center}

El sólido se ve como en la figura \ref{fig:solido_casquillo_ejemplo}:

\begin{figure}[H]
\centering
\caption{El sólido}
\label{fig:solido_casquillo_ejemplo}
\includegraphics[width=0.30\textwidth]{Gráficas/CasquiAlredEjey2Solid.pdf}

\propia
\end{figure}

%\begin{center}------------
%\includegraphics[scale=0.35]{CasquiAlredEjey2Solid.pdf}
%\textbf{Fuente:} elaboración propia
%\end{center}

%%\end{multicols}

Y en la figura \ref{fig:casquillo_cilindrico_figura} se muestra el casquillo cilíndrico generado:

\begin{figure}[H]
\centering
\caption{El casquillo cilíndrico}
\label{fig:casquillo_cilindrico_figura}
\includegraphics[width=0.40\textwidth]{Gráficas/CasquiAlredEjey2Casqui.pdf}

\propia
\end{figure}

%\begin{center}---------
%\centerline{ \graf{El casquillo cilíndrico}}
%\includegraphics[scale=0.35]{CasquiAlredEjey2Casqui.pdf}
%\textbf{Fuente:} elaboración propia
%\end{center}

Con ello, el volumen está dado por la integral:

\resizebox{\textwidth}{!}{$
\begin{aligned}
  V &=&\int_{-1}^{1}2\pi \left( y^{3}+3y+4\right) \left( y+2\right) dy=2\pi
  \int_{-1}^{1}\left( y^{4}+2y^{3}+3y^{2}+6y+4y+8\right) dy \\
  &=&2\pi \int_{-1}^{1}\left( y^{4}+2y^{3}+3y^{2}+10y+8\right) dy \\
  &=&2\pi \left( \int_{-1}^{1}\underset{\text{función par}}{\underbrace{%
    \left( y^{4}+3y^{2}+8\right) }}dy+\int_{-1}^{1}\underset{\text{ función
  impar}}{\underbrace{\left( 2y^{3}+10y\right) }}dy\right) \\
  &=&2\pi \left( 2\int_{0}^{1}\left( y^{4}+3y^{2}+8\right) dy+0\right)=4\pi \left( \frac{1}{5}+3\left( \frac{1}{3}\right) +8\right) =4\pi \left(
  \frac{46}{5}\right) \\
  &=&\frac{184\pi }{5}\text{ unidades c\'ubicas.}%
\end{aligned}
$}

Con el siguiente enlace pueden simularse el sólido y los casquillos en el ejemplo:

\begin{center}
\verb"https://www.geogebra.org/m/wpvxybav"
\end{center}

%\begin{center}---------
%\centerline{ \graf{Parte de la función por medir}}
%\verb"https://www.geogebra.org/m/wpvxybav"
%\end{center}

\section{Longitud de arco}

Si una curva $C$ está dada por la figura de la función $y=f\left(
x\right) $, $a\leq x\leq b$, donde $f^\prime(x)$ es continua en $\left(
a,b\right) $, entonces la longitud de $C$ está dada por
{\small
\[
  L=\int_{a}^{b}\sqrt{1+\left( f^\prime(x)\right)^{2}}\, dx%
\]
}

Demos una idea del origen de la anterior expresión:

\begin{figure}[H]
\centering
\caption{Longitud a medir}
\includegraphics[width=0.36\textwidth]{Gráficas/LongCurvaDef.pdf}

\propia
\end{figure}

%\begin{center}-----------
%\centerline{ \graf{Longitud a medir}}
%\includegraphics[scale=0.4]{LongCurvaDef.pdf}
%\textbf{Fuente:} elaboración propia
%\end{center}

Notemos que para cada n\'umero $x$ del intervalo $\left[ a,b\right],$ el punto $\left(x,f\left( x\right) \right)$ puede verse como un rectángulo diferencial de lados $dx$ y $dy$. Sea $dz$ la diagonal del rectángulo, que es aproximadamente la longitud de la curva alrededor del punto $\left(x,f\left( x\right) \right)$, así que la curva $C$ es $C=\cup _{x\in \left[ a,b\right] }dz.$

Usando el teorema de Pitágoras, se tiene, entonces, que la longitud de $dz$ es
{\small
\[
  L\left( dz\right) =\sqrt{\left( dx\right) ^{2}+\left( dy\right) ^{2}}=%
  \sqrt{\left( dx\right) ^{2}\left( 1+\left( \frac{dy}{dx}\right) ^{2}\right) }%
  =\sqrt{1+\left( \frac{dy}{dx}\right) ^{2}}dx
\]
}
Con ello, la longitud de C es
{\small
\[
  L=L\left( \cup _{x\in \left[ a,b\right] }dz\right)
  =\sum_{x\in \left[ a,b\right] }L\left( dz\right) =\int_{a}^{b}\sqrt{1+\left(
  \frac{dy}{dx}\right) ^{2}}dx
\]
}

\begin{ejemplo}
Halle la longitud de la curva dada por $y=\frac{2}{3}\left( x+1\right) ^{\frac{3}{2}}$ para $2\leq x\leq 7$.
\end{ejemplo}

\emph{Solución.} Encontremos inicialmente la expresión que vamos a integrar:

{\small
\[
  \sqrt{1+\left( y^{\prime}\right) ^{2}}=\sqrt{1+\left( \left( x+1\right) ^{\frac{1%
  }{2}}\right) ^{2}}=\sqrt{1+x+1}=\sqrt{x+2}
\]
}

Notemos que, si $u=\sqrt{x+2}$, entonces $u^{2}=x+2$ y $2udu=dx$. Ahora, si $x=2$, por tanto $u=\sqrt{x+2}=\sqrt{2+2}=2$, y si $x=7$, entonces $u=\sqrt{x+2}=\sqrt{7+2}=3$. Con ello, la longitud de la figura dada por

\begin{figure}[H]
\centering
\caption{Longitud}
\includegraphics[width=0.44\textwidth]{Gráficas/LongCurvaRecta.pdf}

\propia
\end{figure}

%\begin{center}-----------
%\includegraphics[scale=0.4]{LongCurvaRecta.pdf}
%\textbf{Fuente:} elaboración propia
%\end{center}

por ello
{\small
\begin{align*}
  L&=\int_{2}^{7}\sqrt{x+2}dx=\int_2^3 u\times2u\, du
  =\left. \frac{2}{3}u^{3}\right\vert_{2}^{3}=18-\frac{16}{3}=\frac{38}{3}
\end{align*}
}
Es decir, la longitud de la parte de la curva dada es $\frac{38}{3}\text{ unidades.}$

\begin{ejemplo}
Halle la longitud de la curva dada por $y=\ln \left( \sen x\right)$ para $\frac{\pi }{4}\leq x\leq \frac{\pi }{3}.$
\end{ejemplo}
%%\begin{multicols}{2}

\emph{Solución.} En este caso,
\[
\frac{dy}{dx}=\frac{d(\ln \left( \sen x\right))}{dx}=\frac{1}{\sin x}\cos x=\cot x
\]

entonces,
{\small
\begin{align*}
  \sqrt{1+\left(\frac{dy}{dx}\right) ^{2}}&=\sqrt{1+\cot ^{2}x}\\
  &=\sqrt{\csc^{2}x}=\csc x
\end{align*}
}
y la longitud de la parte de la curva

\begin{figure}[H]
\centering
\caption{Longitud por medir}
\includegraphics[width=0.35\textwidth]{Gráficas/LongCurvalnSen.pdf}

\propia
\end{figure}

%\begin{center}-----------
%\centerline{ \graf{Longitud por medir}}
%\includegraphics[scale=0.4]{LongCurvalnSen.pdf}
%\textbf{Fuente:} elaboración propia
%\end{center}

es
%%\end{multicols}
{\small
\begin{align*}
  L &= \int_{\frac{\pi }{4}}^{\frac{\pi }{3}}\csc x\, dx \\
  &= \left( \ln \left\vert \csc x-\cot x\right\vert \right)_{\frac{\pi }{4}}^{\frac{\pi }{3}} \\
  &= \ln \left\vert \csc \frac{\pi }{3}-\cot \frac{\pi }{3}\right\vert - \ln \left\vert \csc \frac{\pi }{4}-\cot \frac{\pi }{4}\right\vert \\
  &= \ln \left\vert \frac{2\sqrt{3}}{3}-\frac{1/2}{\sqrt{3}/2} \right\vert - \ln \left\vert \sqrt{2}-1 \right\vert \\
  &= \ln \left\vert \frac{2\sqrt{3}}{3}-\frac{\sqrt{3}}{3} \right\vert - \ln \left\vert \sqrt{2}-1 \right\vert \\
  &= \ln \left( \frac{\sqrt{3}}{3} \right) - \ln \left( \sqrt{2}-1 \right) \\
  &= \frac{1}{2} \ln 3 - \ln 3 - \ln \left( \sqrt{2}-1 \right) \\
  &= -\frac{1}{2} \ln 3 - \ln \left( \sqrt{2}-1 \right) \\
  &\approx 0,33207\ \text{unidades}
\end{align*}
}

\begin{ejemplo}
Halle la longitud de la curva dada por $y=\frac{1}{4}x^{2}-\frac{1}{2}\ln x$ desde el punto $\left( 1,\frac{1}{4}\right)$ hasta el punto $\left( 4,4-\ln
2\right) $.
\end{ejemplo}

%\begin{center}------------
%\centerline{ \graf{Longitud por medir}}
%\includegraphics[scale=0.4]{LongCurvax2ln.pdf}
%\textbf{Fuente:} elaboración propia
%\end{center}

\emph{Solución.} La figura en este caso es
%\clearpage

\begin{figure}[H]
\centering
\caption{Longitud por medir}
\includegraphics[width=0.48\textwidth]{Gráficas/LongCurvax2ln.pdf}

\propia
\end{figure}

Usando la fórmula de la longitud, tenemos
{\small
\begin{align*}
  L &= \int_{1}^{4} \sqrt{1 + \left( \frac{d}{dx}\left( \frac{1}{4}x^{2} - \frac{1}{2} \ln x \right) \right)^2} \, dx \\
  &= \int_{1}^{4} \sqrt{1 + \left( \frac{1}{2}x - \frac{1}{2x} \right)^2} \, dx \\
  &= \int_{1}^{4} \sqrt{1 + \frac{1}{4}x^2 - \frac{1}{2} + \frac{1}{4x^2}} \, dx \\
  &= \int_{1}^{4} \sqrt{\frac{1}{2} + \frac{1}{4}x^2 + \frac{1}{2x^2}} \, dx \\
  &= \int_{1}^{4} \sqrt{\frac{1}{2} + \frac{1}{4}x^2 + \frac{1}{4x^2}} \, dx \\
  &= \int_{1}^{4} \sqrt{ \frac{1}{4x^2}(x^2 + 1)^2 } \, dx \\
  &= \int_{1}^{4} \frac{1}{2x}(x^2 + 1) \, dx \\
  &= \int_{1}^{4} \left( \frac{x}{2} + \frac{1}{2x} \right) \, dx \\
  &= \left( \frac{x^2}{4} + \frac{1}{2} \ln |x| \right) \Bigg|_{1}^{4} \\
  &= \left( \frac{16}{4} + \frac{1}{2} \ln 4 \right) - \left( \frac{1}{4} + \frac{1}{2} \ln 1 \right) \\
  &= 4 + \frac{1}{2} \ln(2^2) - \frac{1}{4} \\
  &= \ln 2 + \frac{15}{4} \\
  &\approx 4.4431
\end{align*}
}
\begin{ejemplo}
Calcule la longitud de la curva dada por $y=e^{x}$ desde el punto $\left(0,1\right)$, hasta el punto \textit{$\left(\ln 8,8\right)$.}
\end{ejemplo}

\emph{Solución.} La raíz dentro de la integral que debemos calcular es $\sqrt{1+(e^x)^2}$, que no se integra directamente. Hagamos una sustitución trigonométrica con $e^x$ como la variable, es decir, supongamos que $e^{x}=\tan \theta$; entonces, $e^{x}dx=\sec ^{2}\theta d\theta$ y $dx=\cfrac{\sec ^{2}\theta d\theta }{\tan \theta }$, es decir,
\[
\sqrt{1+\left( e^{x}\right) ^{2}}=\sqrt{\sec ^{2}\theta }=\sec \theta
\]
y también $\theta =\tan^{-1} \left( e^{x}\right)$. Con ello, la longitud de la curva para $0\leq x\leq \ln 8$:

\begin{figure}[H]
\centering
\caption{Longitud por medir}
\includegraphics[width=0.32\textwidth]{Gráficas/LongCurvEx.pdf}

\propia
\end{figure}

%\begin{center}-----------
%\centerline{ \graf{Longitud por medir}}
%\includegraphics[scale=0.4]{LongCurvEx.pdf}
%\textbf{Fuente:} elaboración propia
%\end{center}

es
\[
L=\int_0^{\ln 8}\sqrt{1+\left( e^{x}\right) ^{2}}dx
\]
Calculemos la integral sin los límites:

\resizebox{\textwidth}{!}{$
\begin{aligned}
  L&=\int\sqrt{1+\left( e^{x}\right) ^{2}}dx=\int \sec \theta
  \frac{\sec ^{2}\theta }{\tan \theta }\,d\theta=\int
  \,\sec \theta \frac{\tan ^{2}\theta +1}{\tan \theta }\,d\theta=\int \,\left( \sec \theta \tan \theta +\csc \theta \right)\, d\theta \\
  &=\sec \theta +\ln \left( \csc \theta -\cot \theta \right)=\sqrt{1+\left( e^{x}\right) ^{2}}+\ln \left( \frac{\sqrt{1+\left( e^{x}\right) ^{2}}}{e^{x}}-\frac{1}{e^{x}}\right)
\end{aligned}
$}

y entonces,
{\footnotesize
\begin{align*}
  L&=\int_0^{\ln 8}\sqrt{1+\left( e^{x}\right) ^{2}}dx=\left.\left( \sqrt{1+\left( e^{x}\right) ^{2}}+\ln \left( \frac{\sqrt{1+\left( e^{x}\right) ^{2}}}{e^{x}}-\frac{1}{e^{x}}\right) \right)\right\vert_{0}^{\ln 8}\\
  &=\sqrt{65}+\ln \left\vert \frac{\sqrt{65}}{8}-\frac{1}{8}\right\vert -\sqrt{2}-\ln \left\vert \sqrt{2}-1\right\vert
\end{align*}
}
Así, la longitud de la curva es de aproximadamente $7.4047$ unidades.

\begin{ejemplo}
Encuentre la longitud de la curva dada por $y=21\ln x$ para $1\leq x\leq a.$
\end{ejemplo}

\emph{Solución.} La figura que corresponde a la función es

\begin{figure}[H]
\centering
\caption{Longitud por medir}
\includegraphics[width=0.40\textwidth]{Gráficas/LongCurva21ln.pdf}

\propia
\end{figure}

%\begin{center}----------
%\centerline{ \graf{Longitud por medir}}
%\includegraphics[scale=0.4]{LongCurva21ln.pdf}
%\textbf{Fuente:} elaboración propia
%\end{center}

Al plantear la integral de la longitud, tenemos que
{\small
\[
  L=\int_{1}^{a}\sqrt{1+\frac{441}{x^{2}}}dx=\int_{1}^{a}\frac{\sqrt{x^{2}+441%
  }}{x}dx
\]
}
Calculemos la integral sin tener en cuenta los límites.

Si $x=21\tan \theta$, entonces $dx=21\sec ^{2}\theta d\theta$ y también $\sqrt{x^{2}+441}=21\sec \theta$. Con eso, para $1\leq x\leq a$,

\resizebox{\textwidth}{!}{$
\begin{aligned}
  \int\frac{\sqrt{x^{2}+441}}{x}dx&=\int \frac{21\sec \theta \times 21\sec ^{2}\theta d\theta }{21 \tan\theta }=21\int \frac{\sec \theta \left( \tan ^{2}\theta +1\right)
  d\theta }{\tan \theta } \\
  &=21\left( \int \sec \theta \tan \theta d\theta
  +\int \csc \theta d\theta \right)=21\left( \sec \theta +\ln \left\vert \csc \theta -\cot \theta\right\vert \right)\\
  &=21\left( \frac{\sqrt{x^{2}+441}}{21}+\ln \left\vert \frac{\sqrt{x^{2}+441}}{x}-\frac{21}{x}\right\vert \right)\\
\end{aligned}
$}

Con ello, la longitud de la curva es
{\small
\begin{align*}
  L&=\int_1^a\frac{\sqrt{x^{2}+441}}{x}dx\\
  &=21\left.\left( \frac{\sqrt{x^{2}+441}}{21}+\ln \left\vert \frac{\sqrt{x^{2}+441}}{x}-\frac{21}{x}\right\vert \right)\right\vert_{1}^{a}\\
  &=21\left( \frac{\sqrt{a^{2}+441}}{21}+\ln \left\vert \frac{\sqrt{%
    a^{2}+441}}{a}-\frac{21}{a}\right\vert - \frac{\sqrt{442}}{21}-\ln
  \left\vert \sqrt{442}-21\right\vert \right)
\end{align*}
}

\section{Área de una superficie de revolución}

Si la curva dada por la ecuación $y=f\left( x\right) $ para $a \leq x\leq b$ se rota alrededor del eje $x$ y $f^\prime $ es continua en $a \leq x\leq b $, entonces el área de la superficie $S$ de
revolución resultante $a\left( S\right) $ está dada por
{\small
\begin{equation}\label{AreaSuperfAlrx}
  a\left( S\right) =\int_{a}^{b}2\pi \vert f\left( x\right)\vert \sqrt{1+\left( f^\prime(x)\right) ^{2}}dx
\end{equation}
}
Ahora bien, si escribimos la curva dada por $y=f\left( x\right) $ en la forma $g(y)=x $ para $c \leq y\leq d$, se tiene que el área superficial, al girar la misma porción alrededor del eje $x$, también está dada por
{\small
\begin{equation}\label{AreaSuperfAlrxy}
  a\left( S\right) =\int_{c}^{d}2\pi \vert y\vert \sqrt{1+\left( g^\prime(y)\right)^{2}}dy
\end{equation}
}
%%\begin{multicols}{2}
La figura muestra la curva que gira alrededor del eje $x$:

\begin{figure}[H]
\centering
\caption{Curva que gira}
\includegraphics[width=0.38\textwidth]{Gráficas/SupRevolDef.pdf}

\propia
\end{figure}

%\begin{center}----------
%\centerline{ \graf{Curva que gira}}
%\includegraphics[scale=0.3]{SupRevolDef.pdf}
%\textbf{Fuente:} elaboración propia
%\end{center}

Demos una idea intuitiva de la demostración. Cada punto $\left( x,f\left( x\right) \right)$ de la curva puede verse como un rectángulo diferencial de lados $dx$ y $dy$.

\begin{figure}[H]
\centering
\caption{Rectángulo diferencial}
\includegraphics[width=0.35\textwidth]{Gráficas/SupRevolDef2.pdf}

\propia
\end{figure}

%\begin{center}------------
%\centerline{ \graf{Rectángulo diferencial}}
%\includegraphics[scale=0.3]{SupRevolDef2.pdf}
%\textbf{Fuente:} elaboración propia
%\end{center}

Sea $dz$ la diagonal del rectángulo. Para cada $x$ en el intervalo $\left[ a,b\right],$ sea $C_{x}$ la cinta generada por el giro del segmento $dz$ alrededor del eje $x$:

\begin{figure}[H]
\centering
\caption{Cinta generada por el giro}
\includegraphics[width=0.30\textwidth]{Gráficas/SupRevolDefBand.pdf}

\propia
\end{figure}

%\begin{center}----------
%\centerline{ \graf{Cinta generada por el giro}}
%\includegraphics[scale=0.4]{SupRevolDefBand.pdf}
%\textbf{Fuente:} elaboración propia
%\end{center}

Así que la superficie de revolución $S$ obtenida por el giro de la curva $C:\, y=f\left( x\right),\ a\leq x\leq b$ puede expresarse como $S=\cup _{x\in \left[ a,b\right] }C_{x}.$

\begin{figure}[H]
\centering
\caption{Superficie de revolución}
\includegraphics[width=0.28\textwidth]{Gráficas/SupRevolDefCompl.pdf}

\propia
\end{figure}

%\begin{center}------------
%\centerline{ \graf{Superficie de revolución}}
%\includegraphics[scale=0.45]{SupRevolDefCompl.pdf}
%\textbf{Fuente:} elaboración propia
%\end{center}

%%\end{multicols}

Como {\small $dz=\sqrt{\left( dx\right) ^{2}+\left( dy\right) ^{2}}=\sqrt{\left(
dx\right) ^{2}\left( 1+\left( \frac{dy}{dx}\right) ^{2}\right) }=\sqrt{%
1+\left( \frac{dy}{dx}\right) ^{2}}dx$,} el área de $S$, $a\left(
S\right) $ es
{\small
\begin{align*}
  a\left( S\right)&=a\left( \cup _{x\in \left[ a,b\right] }C_{x}\right)
  =\sum_{x\in \left[ a,b\right] }a\left( C_{x}\right) \\
  &=\int_{a}^{b}2\pi \times \text{radio}\times dz=\int_{a}^{b}2\pi \vert f\left( x\right)\vert \sqrt{1+\left( \frac{dy}{dx}\right)^{2}}dx
\end{align*}
}

Por otro lado, si se tiene la curva $x=g(y)$, con $c\leq y\leq d$, y se hace girar alrededor del eje $y$, entonces el área superficial de la superficie así generada puede calcularse con la integral
{\small
\begin{equation}\label{AreaSuperfAlry}
  a\left( S\right) =\int_{c}^{d}2\pi \vert g\left( y\right)\vert \sqrt{1+\left( g^\prime(y)\right) ^{2}}dy
\end{equation}
}
Similarmente, en el caso alrededor de $x$, si escribimos la curva dada por $x=g\left( y\right) $ en la forma $h(x)=y $ para $a \leq x\leq b$, se tiene que el área superficial al girar la misma porción alrededor del eje $y$ también está dada por
{\small
\begin{equation}\label{AreaSuperfAlryx}
  a\left( S\right) =\int_{a}^{b}2\pi\vert x\vert \sqrt{1+\left( h^\prime(x)\right)^{2}}dx
\end{equation}
}

\begin{ejemplo}
Halle el área superficial de la superficie de revolución generada al girar la curva dada por $y=\sqrt{8x-x^{2}}$ para $1\leq x\leq 6$ alrededor del eje $x$.
\end{ejemplo}
\emph{Solución.} La sección de la figura de la función dada es

\begin{figure}[H]
\centering
\caption{Curva por girar}
\includegraphics[width=0.41\textwidth]{Gráficas/SuperRevolAlrxCurva.pdf}

\propia
\end{figure}

%\begin{center}-----------
%\centerline{ \graf{Curva por girar}}
%\includegraphics[scale=0.35]{SuperRevolAlrxCurva.pdf}
%\textbf{Fuente:} elaboración propia
%\end{center}

Si $y=\sqrt{8x-x^{2}}$, entonces $y^\prime=\cfrac{\left( 8-2x\right) }{2\sqrt{8x-x^{2}}}$ y la expresión dentro de la raíz en \eqref{AreaSuperfAlrx} queda como
{\small
\begin{align*}
  1+\left( y^{\prime}\right)^{2}&=1+\frac{\left( 4-x\right) ^{2}}{\sqrt{8x-x^{2}}^{2}}=1+\frac{16-8x+x^{2}}{8x-x^{2}} \\
  &=\frac{8x-x^{2}+16-8x+x^{2}}{8x-x^{2}}=\frac{16}{8x-x^{2}}
\end{align*}
}
y también
{\small
\[
  \sqrt{1+\left(y^{\prime}\right) ^{2}}=\frac{4}{\sqrt{8x-x^{2}}}
\]
}
Con ello,
{\small
\[
  a(S)=\int_{1}^{6}2\pi f\left( x\right) \sqrt{1+\left( f^{\prime}(x)\right) ^{2}}dx
  =\int_{1}^{6}2\pi \sqrt{8x-x^{2}}\frac{4}{\sqrt{8x-x^{2}}}dx=40\pi
\]
}
Es decir, el área superficial de
%\clearpage

\begin{figure}[H]
\centering
\caption{Superficie}

\includegraphics[width=0.35\textwidth]{Gráficas/SuperRevolAlrxSuperf.pdf}

\propia
\end{figure}

%\begin{center}-----------
%\centerline{ \graf{Curva por girar}}
%\includegraphics[scale=0.45]{SuperRevolAlrxSuperf.pdf}
%\textbf{Fuente:} elaboración propia
%\end{center}

es $40\pi.$ En el siguiente enlace puede verse la superficie generada:

\begin{center}
\verb"https://www.geogebra.org/m/zstmbxcs"
\end{center}

%\begin{center}----------
%\centerline{ \graf{Parte de la curva por girar}}
%\verb"https://www.geogebra.org/m/zstmbxcs"
%\end{center}

\begin{ejemplo}
Determine el área de la superficie generada al hacer girar la curva dada por $x=2\sqrt{4-y}$ para $0\leq y\leq \frac{15}{4}$ alrededor del eje $y$.
\end{ejemplo}

\emph{Solución.} La curva que gira en este caso se muestra en la figura:

\begin{figure}[H]
\centering
\caption{Curva a girar}
\includegraphics[width=0.30\textwidth]{Gráficas/SuperRevol2Raiz.pdf}

\propia
\end{figure}

%\begin{center}-----------
%\centerline{ \graf{Superficie generada}}
%\includegraphics[scale=0.3]{SuperRevol2Raiz.pdf}
%\textbf{Fuente:} elaboración propia
%\end{center}

y la superficie formada en este caso es la siguiente:

\begin{figure}[H]
\centering
\caption{Superficie generada}
\includegraphics[width=0.32\textwidth]{Gráficas/SuperRevol2RaizCompl.pdf}

\propia
\end{figure}

%\begin{center}-----------
%\includegraphics[scale=0.45]{SuperRevol2RaizCompl.pdf}
%\textbf{Fuente:} elaboración propia
%\end{center}

Al usar la fórmula dada por \eqref{AreaSuperfAlrxy}, tenemos

{\small
\begin{align*}
  a(S) &= \int_{0}^{\frac{15}{4}} 2\pi \left( 2\sqrt{4 - y} \right)
  \sqrt{1 + \left( \frac{dx}{dy} \right)^2} \, dy \\
  &= \int_{0}^{\frac{15}{4}} 2\pi \left( 2\sqrt{4 - y} \right)
  \sqrt{1 + \left( \frac{-1}{\sqrt{4 - y}} \right)^2} \, dy \\
  &= 4\pi \int_{0}^{\frac{15}{4}} \sqrt{4 - y} \cdot
  \sqrt{1 + \frac{1}{4 - y}} \, dy
\end{align*}
}
Como $0\leq y\leq \frac{15}{4}$, entonces $4-y>0$.
{\small
\[
  \begin{aligned}
    &= 4\pi \int_{0}^{\frac{15}{4}} \sqrt{4 - y} \cdot \sqrt{1 + \frac{1}{4 - y}} \, dy \\
    &= 4\pi \int_{0}^{\frac{15}{4}} \sqrt{4 - y} \cdot \sqrt{\frac{4 - y + 1}{4 - y}} \, dy \\
    &= 4\pi \int_{0}^{\frac{15}{4}} \sqrt{4 - y} \cdot \sqrt{\frac{5 - y}{4 - y}} \, dy
  \end{aligned}
\]
}

pues, $y\in \left[ 0,\frac{15}{4}\right] $, lo que implica que $4-y>0$ y $5-y>0$.
{\small
\begin{align*}
  &= 4\pi \int_{0}^{\frac{15}{4}} \left( \sqrt{4 - y} \right) \frac{\sqrt{5 - y}}{\sqrt{4 - y}} \, dy \\
  &= 4\pi \int_{0}^{\frac{15}{4}} \sqrt{5 - y} \, dy
\end{align*}
}

Ahora, si hacemos $u=5-y,  du=-1dy  dy=-du$ si $y=0,  u=5$  si $y=\frac{15}{4}, u=5-\frac{15}{4}=\frac{20-15} {4}=\frac{5}{4}$;  entonces
{\small
\begin{align*}
  &= -4\pi \int_{5}^{\frac{5}{4}} \sqrt{u} \, du = -4\pi \left( \frac{2}{3} u^{3/2} \right) \Bigg|_{5}^{\frac{5}{4}} \\
  &= -\frac{8}{3} \pi \left( u^{3/2} \right) \Bigg|_{5}^{\frac{5}{4}} = -\frac{8}{3} \pi \left( \left( \frac{5}{4} \right)^{3/2} - 5^{3/2} \right) \\
  &= -\frac{8}{3} \pi \left( \frac{5}{4} \cdot \sqrt{ \frac{5}{4} } - 5 \sqrt{5} \right) = -\frac{8}{3} \pi \left( \frac{5 \sqrt{5}}{8} - 5 \sqrt{5} \right) \\
  &= -\frac{8}{3} \pi \left( \frac{5\sqrt{5} - 40\sqrt{5}}{8} \right) = -\frac{8}{3} \pi \left( \frac{-35\sqrt{5}}{8} \right) \\
  &= \frac{35\sqrt{5} \pi}{3} \quad \text{unidades cuadradas}
\end{align*}
}

En el siguiente enlace puede verse la superficie generada:

\begin{center}
\verb"https://www.geogebra.org/m/mtnbpvk9"
\end{center}

%\begin{center}------------
%\centerline{ \graf{Curva por girar}}
%\verb"https://www.geogebra.org/m/mtnbpvk9"
%\end{center}

%\begin{ejemplo}
%Calcule el área de la superficie de revolución obtenida por el giro
%de la curva dada por $y=2\sqrt{x}$, $3\leq x\leq 8,$ alrededor del eje $x$.
%\end{ejemplo}
%\emph{Solución.} Derivamos dentro de la raíz,
%\[
%\sqrt{1+\left( y^{\prime}\right) ^{2}}=\sqrt{1+\left( \frac{1}{\sqrt{x}}\right)^{2}}=\sqrt{1+\frac{1}{x}}=\frac{\sqrt{x+1}}{\sqrt{x}}
%\]
%Esa expresión debemos integrarla entre $3$ y $8$. Si $ u=\sqrt{x+1}$ entonces $u^{2}=x+1$ y $2udu=dx$. Y entonces,
%\begin{align*}
%a(S)&=2\pi \int_{3}^{8}f(x)\sqrt{1+\left( f^{\prime}(x)\right) ^{2}}%
%dx=2\pi \int_{3}^{8}2\sqrt{x}\frac{\sqrt{x+1}}{\sqrt{x}}dx \\
%&=4\pi \int_{3}^{8}\sqrt{x+1}dx=4\pi \int u2udu\quad (3\leq x\leq8)\footnote{recuerde que $u=\sqrt{x+1}$, entonces cuando $x=3$ se tiene que $u=\sqrt{4}=2$ y cuando $x=8$ entonces $u=\sqrt{9}=3$}\\
%&=4\pi \int_2^3 2u^2du=4\pi \left.\left( \frac{2}{3}y^{3}\right)\right\vert_{2}^{3} \\
%&=\frac{8\pi }{3}\left(27-8\right) =\frac{152}{3}\pi.$
%\end{align*}
%
%\begin{ejemplo}
%Hallar el área de la superficie de revolución obtenida al hacer girar alrededor del eje $x$ la curva dada por $y=\sqrt{100-x^{2}}$ para $1\leq x\leq 5$.
%\end{ejemplo}
%
%\emph{Solución.} La derivada de la función es $\frac{dy}{dx}=\frac{1}{2\sqrt{100-x^{2}}}\left( -2x\right) =\frac{-x}{\sqrt{100-x^{2}}}$. Entonces, $\left( \frac{dy}{dx}\right)^{2}=\frac{x^{2}}{100-x^{2}}$ y $1+\left( \frac{dy}{dx}\right)^{2}=1+\frac{x^{2}}{100-x^{2}}=\frac{100}{100-x^{2}}$. Con ello, $\sqrt{1+\left( \frac{dy}{dx}\right)^{2}}=\sqrt{\frac{100}{100-x^{2}}}
%=\frac{10}{\sqrt{100-x^{2}}},$ y el área superficial es
%\begin{align*}
%a(S)&=2\pi \int_{1}^{5}f\left( x\right)\sqrt{1+\left( \frac{dy}{dx}%
%\right) ^{2}}dx\\
%&=2\pi \int_{1}^{5}\sqrt{100-x^{2}}\frac{10}{\sqrt{100-x^{2}}}
%dx=2\pi \int_{1}^{5}10dx=\left( 20\pi x\right) _{1}^{5}=80\pi
%\end{align*}
%es decir, el área de la superficie generada es $80\pi$ \text{unidades de área}.$

\begin{ejemplo}
Halle el área de la superficie de revolución generada al rotar alrededor del eje $x$ la curva dada por $y=x^{3}+\frac{1}{12}x^{-1}$ para $-2\leq x\leq -1$ y \textit{$1\leq x\leq 2$.}
\end{ejemplo}
%%\begin{multicols}{2}

\emph{Solución.} La figura de los segmentos de la curva que se hacen girar se muestra adelante:

\begin{figure}[H]
\centering
\caption{Curva a girar}
\includegraphics[width=0.35\textwidth]{Gráficas/SupRevolX3.pdf}

\propia
\end{figure}

%\begin{center}-----------
%\centerline{ \graf{La superficie}}
%\includegraphics[scale=0.35]{SupRevolX3.pdf}
%\textbf{Fuente:} elaboración propia
%\end{center}

La superficie resultante es simétrica con respecto al eje $y$, es decir, la parte de la izquierda tiene el mismo tamaño que la parte de la derecha.

Por tanto, el área superficial total es igual a dos veces el área de la parte de la derecha.

En las siguientes figuras podemos ver que efectivamente se tiene la simetría y, por tanto, justificamos lo que se manifiesta en las líneas anteriores.

Por ello calcularemos, entonces, la integral solamente para la parte de la derecha:
%%\end{multicols}

\begin{figure}[H]
\centering
\caption{La superficie}
\label{fig:superficie_derecha}
\includegraphics[width=0.55\textwidth]{Gráficas/SupRevolX3Solido.pdf}

\propia
\end{figure}

%\begin{center}----------
%\includegraphics[scale=0.45]{SupRevolX3Solido.pdf}
%\textbf{Fuente:} elaboración propia
%\end{center}

La derivada de $y$ es $y^{\prime}=3x^{2}-\frac{1}{12}x^{-2}$, entonces
{\small
\begin{align*}
  \sqrt{1+\left( y^{\prime}\right) ^{2}}&=\sqrt{1+\left( 3x^{2}-\frac{1}{12}%
  x^{-2}\right) ^{2}}=\sqrt{1+9x^{4}-\frac{1}{2}+\frac{1}{144}x^{-4}} \\
  &=\sqrt{%
  9x^{4}+\frac{1}{2}+\frac{1}{144}x^{-4}}=\sqrt{\left( 3x^{2}+\frac{1}{12}%
  x^{-2}\right) ^{2}}=3x^{2}+\frac{1}{12}x^{-2}
\end{align*}
}

El área de la parte derecha es
{\small
\begin{align*}
  a(S)&=2\pi \int_{1}^{2}f\left( x\right) \sqrt{1+\left( y^{\prime}\right) ^{2}}dx
  =2\pi \int_{1}^{2}\left( x^{3}+\frac{1}{12}x^{-1}\right) \left( 3x^{2}+%
  \frac{x^{-2}}{12}\right) dx \\
  &=2\pi \int_{1}^{2}\left( \frac{1}{3}x+\frac{1}{144}x^{-3}+3x^{5}\right) dx=2\pi\left. \left( \frac{1}{6}x^{2}-\frac{1}{288}x^{-2}+\frac{3}{6}x^{6}\right)\right\vert_{1}^{2}\\
  &=2\pi \left( \frac{4}{6}-\frac{1}{288\times 4}+2^{5}-\frac{1}{6}+\frac{%
  1}{288}-\frac{1}{2}\right)%\\
  %&
  =\frac{12,289}{192}\pi
\end{align*}
}
y, por tanto, el área total es
{\small
\[
  2\times \frac{12,289}{192}\pi=\frac{12,289}{96}\pi
\]
}
En el siguiente enlace puede verse la superficie generada:

\begin{center}
\verb"https://www.geogebra.org/m/wttm4qnz"
\end{center}

%\begin{center}----------
%\centerline{ \graf{Curva por girar}}
%\verb"https://www.geogebra.org/m/wttm4qnz"
%\end{center}

\begin{ejemplo}
Halle el área de la superficie de revolución que se tiene al girar alrededor de la recta $y=3$ la curva dada por $y=\frac{1}{4}x^{-1}+\frac{1}{3}x^{3}$ para \textit{$1\leq x\leq 2$.}
\end{ejemplo}
%%\begin{multicols}{2}

\emph{Solución.} La parte de la figura \ref{fig:superficie_derecha} que gira se muestra en la siguiente manera:

\begin{figure}[H]
\centering
\caption{Curva a girar}
\includegraphics[width=0.30\textwidth]{Gráficas/SupRevolAlredY3.pdf}

\propia
\end{figure}

%\begin{center}------------
%\centerline{ \graf{Superficie de revolución}}
%\includegraphics[scale=0.35]{SupRevolAlredY3.pdf}
%\textbf{Fuente:} elaboración propia
%\end{center}

y la superficie se ve así:

\begin{figure}[H]
\centering
\caption{La superficie}
\includegraphics[width=0.40\textwidth]{Gráficas/SupRevolAlredY3Superf.pdf}

\propia
\end{figure}

%\begin{center}------------
%\includegraphics[scale=0.45]{SupRevolAlredY3Superf.pdf}
%\textbf{Fuente:} elaboración propia
%\end{center}

%%\end{multicols}
%\eqref{AreaSuperfAlrx}

La expresión en \eqref{AreaSuperfAlrx} se dio cuando la curva gira alrededor del eje $x$. Como en este caso está girando alrededor del eje $y=3$, entonces el radio de la cinta es
\[
r=3-f(x)=3-\frac{1}{4}x^{-1}-\frac{1}{3}x^{3}\text{    además  }  f^\prime(x)=x^{2}-\frac{1}{4}x^{-2}
\]
Con ello,
{\small
\begin{align*}
  \sqrt{1+\left( f^\prime(x)\right) ^{2}}&=\sqrt{1+\left( x^{2}-\frac{1}{4}%
  x^{-2}\right) ^{2}}=\sqrt{1+\left( x^{4}-\frac{1}{2}+\frac{1}{16}%
  x^{-4}\right) ^{2}} \\
  &=\sqrt{\left( x^{4}+\frac{1}{2}+\frac{1}{16}x^{-4}\right) ^{2}}=\sqrt{\left(
  x^{2}+\frac{1}{4}x^{-2}\right) ^{2}}=x^{2}+\frac{1}{4}x^{-2}
\end{align*}
}

Entonces el área es
{\small
\begin{align*}
  a(S)&=2\pi \int_{1}^{2}\left( 3-\frac{1}{4x}-\frac{1}{3}x^{3}\right) \left( x^{2}+%
  \frac{1}{4x^{2}}\right) dx \\
  &=2\pi \int_{1}^{2}\left(-\frac{1}{3}x+\frac{3}{4}x^{-2}+3x^{2}-\frac{1}{16}x^{-3}-\frac{1}{3}x^{5}\right)dx \\
  &=2\pi \left.\left( -\frac{1}{6}x^{2}-\frac{3}{4}x^{-1}+x^{3}+\frac{1}{32}x^{-2}-\frac{1}{18}x^{6}\right)\right\vert_{1}^{2} \\
  &=2\pi \left( -\frac{4}{6}-\frac{3}{8}+8+\frac{1}{128}-\frac{64}{18}-\left(
  -\frac{1}{6}-\frac{3}{4}+1+\frac{1}{32}-\frac{1}{18}\right) \right) \\
  &=\frac{429}{64}\pi \approx 21.058 \text{ unidades cuadradas}
\end{align*}
}

La superficie generada puede manipularse en el siguiente enlace:

\begin{center}
\verb"https://www.geogebra.org/m/sfyudysd"
\end{center}

%\begin{center}-----------
%\centerline{ \graf{Curva por girar}}
%\verb"https://www.geogebra.org/m/sfyudysd"
%\end{center}

\begin{ejemplo}
Determine el área de la superficie generada al hacer girar la curva dada por $x=2\ln y$ para $1\leq y\leq 2$ alrededor del eje $x.$
\end{ejemplo}
%%\begin{multicols}{2}
\emph{Solución.} La parte de la curva por girar alrededor del eje $x$ es

\begin{figure}[H]
\centering
\caption{Curva a girar}
\includegraphics[width=0.38\textwidth]{Gráficas/SupRevol2lny.pdf}

\propia
\end{figure}

%\begin{center}-----------
%\centerline{ \graf{Superficie generada}}
%\includegraphics[scale=0.35]{SupRevol2lny.pdf}
%\textbf{Fuente:} elaboración propia
%\end{center}

De la función $x=2\ln y$, despejando $y$, se tiene $e^{\frac{x}{2}}=y$ con $0\leq x\leq 2\ln 2$. Ahora, usando la fórmula para la integral, tenemos
{\small
\[
  1+\left(\left(e^{\frac{x}{2}}\right)^{\prime}\right)^{2}=1+\left( \frac{1}{2}%
  e^{\frac{x}{2}}\right) ^{2}=1+\frac{1}{4}e^{x}
\]
}
Así, el área superficial en la figura \ref{fig:superficie_generada}:

\begin{figure}[H]
\centering
\caption{Superficie generada}
\label{fig:superficie_generada}
\includegraphics[width=0.38\textwidth]{Gráficas/SupRevol2lnySuperf.pdf}

\propia
\end{figure}

%\begin{center}------------
%\includegraphics[scale=0.45]{SupRevol2lnySuperf.pdf}
%\textbf{Fuente:} elaboración propia
%\end{center}

está dada por
%%\end{multicols}
{\small
\begin{align*}
  a\left( S\right)= & \int_{0}^{2\ln \left( 2\right) }2\pi \left( e^{\frac{x}{2}%
  }\right) \sqrt{1+\frac{1}{4}e^{x}}dx =\int_{0}^{2\ln \left( 2\right) }2\pi\left( e^{\frac{x}{2}}\right) \sqrt{\frac{4+e^{x}}{4}}dx\\
  =&2\pi \int_{0}^{2\ln \left( 2\right) }\frac{e^{\frac{x}{2}}}{2}\sqrt{4+e^{x}}dx
\end{align*}
}

Notemos que esta integral no puede resolverse por los métodos de integración tradicionales para conseguir un valor exacto.

Veamos la segunda definición, con la que sí podemos terminar el ejercicio:
{\small
\[
  \begin{aligned}
    a(S) &= \int_{1}^{2} 2\pi y \sqrt{1 + \left( (2\ln y)' \right)^2} \, dy = 2\pi \int_{1}^{2} y \sqrt{1 + \left( \frac{2}{y} \right)^2} \, dy \\
    &= 2\pi \int_{1}^{2} y \sqrt{1 + \frac{4}{y^2}} \, dy = 2\pi \int_{1}^{2} y \sqrt{ \frac{1}{y^2} (y^2 + 4) } \, dy \\
    &= 2\pi \int_{1}^{2} y \sqrt{ \frac{y^2 + 4}{y^2} } \, dy = 2\pi \int_{1}^{2} y \cdot \frac{ \sqrt{y^2 + 4} }{y} \, dy = 2\pi \int_{1}^{2} \sqrt{y^2 + 4} \, dy
  \end{aligned}
\]
}

Usamos en la \'ultima igualdad que, dada la ecuación $1\leq y\leq 2$, se tiene que $y^{2}+4>0$ y $y^{2}>0$.

Ahora tomemos $y=2\tan (z)$, con $\frac{-\pi }{2}<z<\frac{\pi }{2}$, de donde $dy=2\sec
^{2}\left( z\right) dz$ y $z=\arctan \left( \frac{y}{2}\right)$.

Con ello, la integral se convierte en la siguiente, pero recordemos que $y$ está entre $1$ y $2$:

\resizebox{\textwidth}{!}{$
\begin{aligned}
  &= \left( \int 2\pi \sqrt{4\tan^2(z) + 4} \cdot \left( 2\sec^2(z) \right) \, dz \right) = 4\pi \left( \int \sqrt{4(\tan^2(z) + 1)} \cdot \sec^2(z) \, dz \right) \\
  &= 4\pi \left( \int 2\sqrt{\sec^2(z)} \cdot \sec^2(z) \, dz \right) = 8\pi \left( \int \sec^3(z) \, dz \right) \\
  &= 8\pi \left[ \frac{1}{2} \sec(z)\tan(z) + \frac{1}{2} \int \sec(z) \, dz \right]_1^2 = 8\pi \left[ \frac{1}{2} \sec(z)\tan(z) + \frac{1}{2} \ln \left| \sec(z) + \tan(z) \right| \right]_1^2 \\
  &= 8\pi \left[ \frac{1}{2} \cdot \frac{\sqrt{4 + y^2}}{2} \cdot \frac{y}{2} + \frac{1}{2} \ln \left| \frac{\sqrt{4 + y^2}}{2} + \frac{y}{2} \right| \right]_1^2 = 4\pi \left( \sqrt{4 + y^2} \cdot \frac{y}{2} + \ln \left| \frac{\sqrt{4 + y^2}}{2} + \frac{y}{2} \right| \right) \Bigg|_1^2 \\
  % La línea más problemática se divide para garantizar que quepa en el margen
  &= 4\pi \left[ \left( \sqrt{4 + 4} \cdot 1 + \ln \left| \frac{\sqrt{8}}{2} + 1 \right| \right) \right. \\
  & \quad \left. - \left( \sqrt{5} \cdot \frac{1}{2} + \ln \left| \frac{\sqrt{5}}{2} + \frac{1}{2} \right| \right) \right] \\
  &= 4\pi \left( 2\sqrt{2} + \ln \left| \sqrt{2} + 1 \right| - \frac{\sqrt{5}}{2} - \ln \left| \frac{\sqrt{5} + 1}{2} \right| \right) \\
  &= 4\pi \left( \frac{4\sqrt{2} - \sqrt{5}}{2} + \ln \left( \frac{2(\sqrt{2} + 1)}{\sqrt{5} + 1} \right) \right) \\
  &= 4\pi \left( \frac{4\sqrt{2} - \sqrt{5}}{2} + \ln \left( \frac{2\sqrt{2} + 2}{\sqrt{5} + 1} \right) \right) \quad \text{unidades cuadradas}
\end{aligned}
$}

La superficie generada puede manipularse en el siguiente enlace:

\begin{center}
\verb"https://www.geogebra.org/m/cuwqsya6"
\end{center}

%\begin{center}------------
%\centerline{ \graf{Fuerza como función a trozos}}
%\verb"https://www.geogebra.org/m/cuwqsya6"
%\end{center}

\section{Trabajo}

En física, cuando una fuerza constante $F$ mueve un objeto una distancia $d$ en la misma dirección de la fuerza, el trabajo realizado se define como el producto $fuerza \times \ distancia$:

\[
W=F\cdot d
\]

Por ejemplo, si una fuerza de $20$ libras mueve un objeto $8$ pies en la misma
dirección de la fuerza, entonces el trabajo realizado es $160$ pies-lb. En
esta sección veremos cómo encontrar el trabajo realizado por una
fuerza variable.

Antes de examinar el trabajo como integral definida, revisaremos algunas
unidades de medida importantes.

\subsection*{Unidades de medida}
En la tabla 2 se enumeran unidades de uso común de fuerza, distancia y trabajo.
\begin{table}[H]
\caption{Unidades de uso común}
\begin{tabular}{|l|c|c|c|}
  \hline
  \textbf{Cantidad} & \textbf{Sistema ingenieril} & \textbf{SI} & \textbf{cgs} \\ \hline
  Fuerza & Libra (lb) & Newton (N) & Dina \\ \hline
  Distancia & Pie (pie) & Metro (m) & Centímetro (cm) \\ \hline
  Trabajo & Pie-libra & Newton-metro & Dina-centímetro
  \\
  & (pie-lb) & (julio) & (ergio) \\ \hline
\end{tabular}

\propia
\end{table}

Por tanto, si una fuerza de $200$N mueve $12$\,m un objeto, el trabajo realizado
es $W=200\times12=2400$N-m o $2400$ julios.

Para efectos de comparación y conversión de una unidad a otra, tenga en cuenta que
\begin{align*}
1\,\text{N}=& 105\, \text{dinas}=0,2247\, \text{lb} \\
1\,\text{pie-lb} =& 1,356\, \text{julios} =1,356\times 10^{7}\, \text{ergios}
\end{align*}
De modo que, por ejemplo, $60\,$ pies-lb equivalen a $60\times 1,356=81,36$ julios y $5700$ julios equivalen a $\frac{5700}{1,356}=4203,54$ pies-lb.

\subsection*{El trabajo como una integral}

Si $F(x)$ representa una fuerza
variable continua que actúa sobre un intervalo $[a,b]$, entonces el
trabajo no es simplemente el producto $W=F\cdot d$.

Supongamos que $P$ es la partición $a=x_{0}<x_{1}<x_{2}<\cdots <x_{n}=b$ y $%
\Delta x_{k}$ es el ancho del $k-{e}simo$ subintervalo $%
[x_{k-1},x_{k}]$. Sea $x_{k}^{\ast }$\ el punto muestra escogido de manera
arbitraria en cada subintervalo. Si el ancho de cada $[x_{k-1},x_{k}]$\ es
muy pequeño, entonces, puesto que $F$ es continua, los valores
de $F(x)$ no pueden variar mucho en el subintervalo. Por tanto, puede considerarse en forma razonable que la fuerza actúa sobre $[x_{k-1},x_{k}]$ como la constante $F\left( x_{k}^{\ast }\right) $ y
que el trabajo realizado desde $x_{k-1}$ hasta $x_{k}$ está dado por la
aproximación
\[
W_{k}=F\left( x_{k}^{\ast }\right) \triangle x_{k}
\]

Entonces, una aproximación al trabajo total realizado desde $a$ hasta $b$
está dada por la suma de Riemann:
\[
W\approx \sum_{k=1}^{n}W_{k}=\sum_{k=1}^{n}F\left( x_{k}^{\ast }\right)
\triangle x_{k}
\]

Resulta natural suponer que el trabajo realizado por $F$ sobre el intervalo
es
\[
W=\lim_{\left\Vert P\right\Vert \rightarrow 0}\sum_{k=1}^{n}F\left(
x_{k}^{\ast }\right) \triangle x_{k}=\int_{a}^{b}F\left( x\right) dx
\]

El análisis anterior se resume en la siguiente definición.

\begin{defi}
Sea $F$ continua sobre el intervalo $[a,b]$ y sea $F(x)$ la fuerza en un número $x$ en el intervalo. \ Entonces, el trabajo $W$ realizado por la fuerza para mover un objeto de $a$ a $b$ es
\begin{equation}\label{trabajo}
  W=\int_{a}^{b}F\left( x\right) dx
\end{equation}
\end{defi}

\textbf{Nota:} Si $F$ es constante, $F(x)=k$ para toda $x$ en el intervalo, entonces
\eqref{trabajo} se vuelve
\[W=\int_{a}^{b}F\left( x\right) dx=\int_{a}^{b}kdx=k\left( b-a\right)\]
lo cual es consistente con el trabajo para una fuerza constante.

\begin{ejemplo}
Cuando una partícula se localiza a una distancia de $x$ metros desde el origen, una fuerza de $F\left( x\right) =\cos \left( \frac{\pi x}{3}\right) $ newtons actúa sobre ella. \textquestiondown Cuánto trabajo se realiza al mover la partícula desde $x=0$ hasta $x=2?$
\end{ejemplo}

\emph{Solución.} Con la información dada, tenemos que

\resizebox{\textwidth}{!}{$
\begin{aligned}
  W&=\int_{a}^{b}F\left( x\right) dx=\int_{0}^{2}\cos \left( \frac{\pi }{3}%
  x\right) dx=\frac{1}{\frac{\pi }{3}}\left( \sin \left( \frac{\pi }{3}%
  x\right) \right) _{0}^{2}=\frac{3}{\pi }\left( \sin \left( \frac{2\pi }{3}%
  \right) -\sin 0\right) \\
  &=\frac{3\sqrt{3}}{2\pi }\, \text{julios}
\end{aligned}
$}

Es decir, el trabajo realizado por la fuerza dada al mover la partícula es de {\small$\cfrac{3\sqrt{3}}{2\pi }\, \text{julios}$.}

\begin{ejemplo}
Se muestra la figura de una función fuerza (en newtons) que se
incrementa a su máximo valor y luego permanece constante:

\begin{figure}[H]
  \centering
  \caption{Fuerza como función a trozos}
  \includegraphics[width=0.45\textwidth]{Gráficas/FueraFuncTroz.pdf}
  \propia
\end{figure}

%\begin{center}------------
%\includegraphics[scale=0.25]{FueraFuncTroz.pdf}
%\textbf{Fuente:} elaboración propia
%\end{center}

¿Cuánto trabajo realiza la fuerza al mover un objeto una distancia de $8$ m?
\end{ejemplo}

\emph{Solución.} Recordemos que una integral definida de una función positiva es el área bajo la figura de la función. Entonces nuestra integral de trabajo, para este caso particular, la podemos manejar de la siguiente manera:
{\small
\begin{align*}
  W&=\int_{0}^{8}F\left( x\right) dx=\text{área bajo la curva}
  =\text{área del trapecio}\\
  &=\frac{\left( B+b\right) h}{2}=\frac{\left(
  8+4\right) \times 3}{2}=18\text{ julios}
\end{align*}
}

\section{Problemas de resortes}

La ley de Hooke establece que, cuando un resorte se estira (o comprime) más allá de su longitud natural, la fuerza de reconstitución ejercida por el resorte es directamente proporcional a la cantidad de elongación (o compresión) $x$. Así, para estirar un resorte $x$ unidades más allá de su longitud natural, es necesario aplicar una fuerza $F(x)=kx,$ donde $k$ es una constante de proporcionalidad denominada \textit{constante del resorte}:

\begin{figure}[H]
\centering
\caption{Elongación del resorte}
\includegraphics[width=0.45\textwidth]{Gráficas/Resorte.pdf}

\propia
\end{figure}

%\begin{center}-----------
%\includegraphics[scale=0.35]{Resorte.pdf}
%\textbf{Fuente:} elaboración propia
%\end{center}

\begin{ejemplo}
\textbf{Alargamiento de un resorte}: para estirar un resorte de $50$ cm, se requiere una fuerza de $130$ N. Encuentre el trabajo realizado para estirar el resorte $20$ cm más allá de su longitud natural.
\end{ejemplo}

\emph{Solución.} Cuando una fuerza se mide en newtons, las distancias suelen expresarse en
metros. El dato estirar $50$ cm requiere una fuerza de$130$ N se usa para hallar la
constante $k$ del resorte:
\[
F=kx\rightarrow k=\frac{F}{x}
\]
Por otro lado, dado que, cuando $F=130$ N, el estiramiento es de $x=50$ cm$=\cfrac{1}{2}$ m, entonces la constante del resorte es
{\small
\[
  k=\frac{\text{Fuerza}}{\text{Elongación}}=\frac{130N}{\frac{1}{2}m}=260N/m
\]
}
Por tanto, $F\left( x\right) =260x$. Luego, el trabajo realizado para estirar el
resorte 20 CM = 0,2 M %$20 cm=0\cm2 m$
a partir de $x=0~($longitud natural del resorte$)$ es
{\small
\[
  W=\int_{a}^{b}F\left( x\right) \mathrm{d}x=\int_{0}^{0.2}260x\,\mathrm{d}x=\left. 130x^{2}\right|_{0}^{0.2}=130\left( 0.2\right)^{2} =5.2 \text{ julios}
\]
}
\section{Problemas de bombeo}

Cuando un líquido que pesa $\rho $ $lb/pie^{3}$ se bombea desde un
tanque, el trabajo realizado para mover un volumen fijo o una capa de líquido $d$ pies en una dirección vertical es

\resizebox{\textwidth}{!}{$
\begin{aligned}
  W &= \text{fuerza} \times \text{distancia} = (\text{peso por unidad de volumen}) \times (\text{volumen}) \times (\text{distancia}) \\
  W &= \rho \times (\text{volumen}) \times d
\end{aligned}
$}

En física, la cantidad $\rho $ se denomina peso específico del
fluido. Para el agua, {\small $\rho =62,4$} {\small lb/pie$^{3}$} o {\small $9800$ N/m$^{3}.$} Queremos, ahora, hallar el trabajo requerido para bombear el líquido
que se encuentra dentro del tanque hasta una altura $h:$

\begin{figure}[H]
\centering
\caption{Trabajo para bombear un líquido}
\label{fig:trabajo_bombear}
\includegraphics[width=0.45\textwidth]{Gráficas/BombeoDef.pdf}

\propia
\end{figure}

%\begin{center}-----------
%\centerline{ \graf{Trabajo para bombear un líquido}}
%\includegraphics[scale=0.35]{BombeoDef.pdf}
%\textbf{Fuente:} elaboración propia
%\end{center}

Ubicamos el eje $x$ hacia arriba, como la figura \ref{fig:trabajo_bombear}. Partimos de una sección transversal del líquido $S_{x}$ a una altura $x\in \left[ a,b\right],$ con un espesor $dx$; hay que bombear esa sección de agua hasta el punto $c$. El trabajo realizado para bombear la sección $S_{x}$ hasta una
altura de $c$ es

{\small
\begin{align*}
  W_{x}&=\left(
    \begin{tabular}{c}
      Peso específico \\
      del líquido%
    \end{tabular}
  \right)\times \left(
    \begin{tabular}{c}
      Volumen \\
      de S$_{x}$%
    \end{tabular}%
  \right)\times \left(
    \begin{tabular}{l}
      Distancia para \\
      mover $S_{x}$%
    \end{tabular}%
  \right)
  \\
  &=\rho (\text{área}\left( S_{x}\right) \times dx)\left( c-x\right) =\rho\,\text{área}\left( S_{x}\right) \left( c-x\right) dx
\end{align*}
}

El trabajo total es
{\small
\begin{align*}
  W&=\sum_{x\in \left[ a,b\right] }W_{x}=\int_{a}^{b}\rho\, \text{área}\left( S_{x}\right) \left( c-x\right) dx
  =\rho \int_{a}^{b}\text{área}\left( S_{x}\right) \underset{\substack{\text{distancia desde $x$}\\ \text{
  hasta la altura final}}}{\underbrace{\left( c-x\right) }}dx
\end{align*}
}

En general, para bombeo en tanques aplicamos la siguiente fórmula
{\small
\begin{equation}\label{Bombeo}
  W=\int_{a}^{b}\rho \text{área}\left( S_{x}\right)\text{distancia}\, dx
\end{equation}
}

Donde $\rho$ es el peso específico del líquido, $\left[ a,b\right] $ es el intervalo en
el eje $x$ donde está el líquido, $S_{x}$ es la sección transversal del líquido (perpendicular al eje $x$ en el punto $x$) y $distancia$ es la distancia para mover la sección $S_{x}$.

\begin{ejemplo}
Un tanque hemisférico de radio $10$ pies está lleno de agua hasta una profundidad de $8$ pies. Encuentre el trabajo realizado para bombear toda el agua hasta la parte superior del tanque.
\end{ejemplo}

\emph{Solución.} Tomamos un sistema de coordenadas centrado en el fondo del tanque y con el
eje $x$ hacia arriba

\begin{figure}[H]
\centering
\caption{Tanque hemisférico}
\includegraphics[width=0.45\textwidth]{Gráficas/TrabajTanqueHemisf.pdf}

\propia
\end{figure}

%\begin{center}-----------
%\centerline{ \graf{Tanque hemisférico}}
%\includegraphics[scale=0.3]{TrabajTanqueHemisf.pdf}
%\textbf{Fuente:} elaboración propia
%\end{center}

Si vemos el tanque como una región en el plano, entonces la semicircunferencia centrada en $\left( 10,0\right)$ con radio $r=10$ (la proyección del tanque al plano) tiene ecuación $\left( x-10\right) ^{2}+y^{2}=10^{2}$. La sección transversal $S_{x}$ es el disco de radio $y$ donde $y^{2}=100-\left( x-10\right) ^{2}$ y, además,
\[
\text{área}\left( S_{x}\right) =\pi y^{2}=\pi \left( 100-\left( x-10\right)^{2}\right)
\]
y con ello,

\resizebox{\textwidth}{!}{$
\begin{aligned}
  W&=\int_{a}^{b}\rho\times \text{ área}\left( S_{x}\right)\times \text{distancia } dx
  =62,4\int_{0}^{8}\pi \left( 100-\left( x-10\right) ^{2}\right) \left(
  10-x\right) dx\\
  &=62,4\int_{0}^{8}\pi \left( -x^{2}+20x\right) \left( 10-x\right) dx\\
  &=62,4\int_{0}^{8}\pi \left( x^{3}-30x^{2}+200x\right) dx\\
  &=62,4\pi \left( \frac{1}{4}x^{4}-10x^{3}+100x^{2}\right)\Big|_{0}^{8}\\
  &=4,\allowbreak 516\,7\times 10^{5} &\text{pies-lb}
\end{aligned}
$}

Tomemos, para este ejemplo, una solución alterna: un sistema de coordenadas centrado en el centro de la esfera y con el eje $x$\ hacia abajo.

\begin{figure}[H]
\centering
\caption{Tanque hemisférico. Dirección positiva del eje $x$ hacia abajo}
\includegraphics[width=0.45\textwidth]{Gráficas/TrabajTanqueHemisf2.pdf}

\propia
\end{figure}

%\begin{center}-----------
%\centerline{ \graf{Tanque hemisférico. Dirección positiva del eje $x$ hacia abajo}}
%\includegraphics[scale=0.35]{TrabajTanqueHemisf2.pdf}
%\textbf{Fuente:} elaboración propia
%\end{center}

Como antes, \textit{S}$_{x}$ es el disco de radio $y$ donde $x^{2}+y^{2}=10^{2}$. Área $\left(
S_{x}\right) =\pi y^{2}=\pi \left( 100-x^{2}\right)$; por tanto,
{\small
\begin{align*}
  W&=\int_{a}^{b}\rho \times \text{area}\left( S_{x}\right)\times \text{distancia }dx\\
  &= \num{62.4}\int_{2}^{10}\pi \left( 100-x^{2}\right) xdx\\
  &= \num{451.67}\times 10^{3}\,\text{pies-lb}
\end{align*}
}

\begin{ejemplo}
Un tanque parabólico, generado por el giro de la parábola $y=x^{2}$ para $0\leq x\leq 3,$ en pies, alrededor del eje $y$, está lleno de agua. Halle el trabajo realizado para bombear el agua hasta $1$ pie por encima del tanque.
\end{ejemplo}

\emph{Solución.} Dibujamos el tanque:

\begin{figure}[H]
\centering
\caption{Tanque parabólico}
\includegraphics[width=0.38\textwidth]{Gráficas/TrabajTanqueParabolic.pdf}

\propia
\end{figure}

%\begin{center}-----------
%\centerline{ \graf{Tanque parabólico}}
%\includegraphics[scale=0.35]{TrabajTanqueParabolic.pdf}
%\textbf{Fuente:} elaboración propia
%\end{center}

Para un punto $y$ en el intervalo $\left[ 0,4\right] $, sea $S_{y}$ la sección transversal (disco) de agua perpendicular al eje $y$ que pasa por $%
\left( 0,y\right) $. El radio de $S_{y}$ se obtiene despejando $x$ en $%
y=x^{2} $ así, que es $R=\sqrt{y}$, y el área del disco es

\[
\text{área}\left( S_{y}\right) =\pi R^{2}=\pi y
\]

La distancia para mover el disco de agua es $d=5-y$ y el trabajo realizado para
bombear el agua es

\resizebox{\textwidth}{!}{$
\begin{aligned}
  W&=\int_{0}^{4}\rho\, \text{área}\left( S_{y}\right)\times \text{distancia } dy=62,4\int_{0}^{4}\pi y\left( 5-y\right) dy=62,4\int_{0}^{4}\pi \left( 5y-y^{2}\right) dy\\
  &=62,4\pi \left( \frac{5}{2}y^{2}-\frac{1}{3}y^{3}\right) _{0}^{4}=62,4\pi \left( 40-\frac{64}{3}\right) =3659,3\text{ pies-lb}
\end{aligned}
$}

\section{Presión y fuerza de un fluido}

La mayoría de las personas han experimentado que se les "tapan los oídos" e incluso han experimentado dolor en los oídos cuando descienden en avión (o en un elevador), o cuando nadan hacia el fondo de una piscina. Estas sensaciones molestas en los oídos se deben a un incremento en la presión ejercida por el aire o el agua sobre mecanismos en el oído medio. El aire y el agua son ejemplos de fluidos. En esta sección se mostrará la forma en que la integral definida puede
usarse para encontrar la fuerza ejercida por un fluido.

Suponga que una placa horizontal plana se
sumerge en un fluido como el agua. La fuerza ejercida por el fluido exactamente arriba de la placa, denominada \textit{fuerza $F$ del fluido}, se define como
{\small
\begin{equation}\label{FuerzaFluido}
  F=\underbrace{(\text{fuerza por unidad de área})}_{\text{presión del fluido $P$\ }}\times (\text{área de superficie})=PA
  %\underset{}{(\underbrace{\times
\end{equation}
}

Si $\rho $ denota el peso específico del fluido (peso por unidad de
volumen) y $A$ es el área de la placa horizontal sumergida hasta una
profundidad $h$ (\textbf{(figura \ref{fig:palanca_horizontal})}) (figura 135):

\begin{figure}[H]
\centering
\caption{Placa horizontal sumergida}
\label{fig:palanca_horizontal}
\includegraphics[width=0.38\textwidth]{Gráficas/PresPlacaHoriz.pdf}

\propia
\end{figure}

%\begin{center}------------
%\centerline{ \graf{\label{PlacHoriza} Placa horizontal sumergida}}
%\includegraphics[scale=0.35]{PresPlacaHoriz.pdf}
%\textbf{Fuente:} elaboración propia
%\end{center}

entonces la presión $P$
del fluido sobre la placa puede expresarse en términos de $\rho $:
\begin{equation}\label{presion}
P=(\text{peso por unidad de volumen})\times (\text{profundidad})=\rho h
\end{equation}

En consecuencia, la fuerza \eqref{FuerzaFluido} del fluido es la misma que
\begin{align*}
F&=(\text{presión de fluido})\times (\text{área de superficie})=\rho hA\\
&=\text{peso específico$ \,\times $\,profundidad$ \,\times $ área}
\end{align*}

No obstante, cuando se sumerge una placa vertical, la presión del fluido
y la fuerza del fluido sobre un lado de la placa varían con la
profundidad:

\begin{figure}[H]
\centering
\caption{Placa vertical sumergida}
\includegraphics[width=0.30\textwidth]{Gráficas/PresPlacaVertic.pdf}

\propia
\end{figure}

%\begin{center}-----------
%\centerline{ \graf{Placa vertical sumergida} \label{PlacVertica}}
%\includegraphics[scale=0.35]{PresPlacaVertic.pdf}
%\textbf{Fuente:} elaboración propia
%\end{center}

Por ejemplo, la presión del fluido sobre una presa vertical es menor en la parte superior que en su base.

Antes de empezar, considere un ejemplo simple de presión y fuerza de una
placa sumergida horizontalmente.

Sobre una placa vertical, la presión varía de la parte superior al fondo.

\begin{ejemplo} \textbf{Presión y fuerza:} una placa rectangular plana de $5 \times 6$ pies se sumerge horizontalmente en agua a una profundidad de $10$ pies. Determine la presión y la fuerza ejercidas sobre la placa por el agua arriba de esta.
\end{ejemplo}

\begin{figure}[H]
\centering
\caption{Placa horizontal}
\label{fig:palanca_plana}
\includegraphics[width=0.30\textwidth]{Gráficas/LaminaHorizontal5_6.pdf}

\propia
\end{figure}

%\begin{center}------------
%\centerline{ \graf{\label{PlacaPlana1} Placa plana vertical}}
%\includegraphics[scale=0.35]{LaminaHorizontal5_6.pdf}
%\textbf{Fuente:} elaboración propia
%\end{center}
%\centerline{\graf{ Lámina horizontal $5\times 6$}\label{Lamina56}}

\emph{Solución.} Recuerde que el peso específico del agua es $62,4$ lb/pie
$^{3}$. Así, por \eqref{presion}, la presión del fluido es
\[
P=\rho h=(62,4\text{lb/pie}^{3})\cdot (10\text{pies})=624 \text{lb/pie}^{2}
\]

Puesto que el área superficial de la placa es $A=30$ pies $^{2}$, por \eqref{FuerzaFluido},
se concluye que la fuerza del fluido sobre la placa es
\[
F=PA=(\rho h)A=(624\text{lb/pie}^{2})\cdot (30\text{pies}^{2})=18,720\text{lb}
\]

Para determinar la fuerza total $F$ ejercida por un fluido sobre un lado de
una superficie plana sumergida verticalmente, se emplea una forma del
principio de Pascal: \textsl{la presión ejercida por un fluido a una profundidad $h$ es
la misma en todas direcciones}. Entonces, si en un gran contenedor con fondo
plano y paredes verticales se vierte agua hasta una profundidad de $10$ pies,
la presión de $624$ lb/pie$^{2}$ en el fondo se ejerce de la misma forma
sobre las paredes (véase la figura \ref{fig:palanca_plana}).

Considere que el eje $x$ positivo está dirigido hacia abajo con el origen en la superficie del
fluido. Suponga que una placa plana vertical, limitada por las rectas
horizontales $x=a$ y $x=b$, se sumerge en el fluido, como se muestra en la figura \ref{fig:rectangulo_base}.
%%\begin{multicols}{2}

\begin{figure}[H]
\centering
\caption{Placa plana vertical}
\label{fig:rectangulo_base}
\includegraphics[width=0.30\textwidth]{Gráficas/FuerzaFluid1.pdf}

\propia
\end{figure}

%\begin{center}-----------
%\centerline{ \graf{\label{PlacaPlana2} Placa plana vertical. Rectángulo base}}
%\includegraphics[scale=0.3]{FuerzaFluid1.pdf}
%\textbf{Fuente:} elaboración propia
%\end{center}

y también

\begin{figure}[H]
\centering
\caption{Placa plana vertical. Rectángulo base}
\includegraphics[width=0.30\textwidth]{Gráficas/FuerzaFluid2.pdf}

\propia
\end{figure}

%\begin{center}-----------
%\includegraphics[scale=0.3]{FuerzaFluid2.pdf}
%\textbf{Fuente:} elaboración propia
%\end{center}

%%\end{multicols}
Sea $w(x)$ la función que denota el ancho de la placa en cualquier nú%
mero $x$ en $[a,b]$. La superficie $S$ la podemos ver como unión de rectángulos, para cada $x\in \lbrack a,b],$ sea $R_{x}$ el rectángulo de base $w\left( x\right) $ y altura $dx.$ \ La fuerza $F_{x}$ ejercida por el fluido sobre el elemento rectangular $R_{x}$ es
\[
F_{x}=\rho \cdot x\cdot w(x)dx
\]

donde, como antes, $\rho $ denota el peso específico del fluido. As%
í, la fuerza total del fluido sobre la placa es
\[
\sum_{x\in \left[ a,b\right] }F_{x}=\int_{a}^{b}\rho xw\left( x\right) dx
\]

\begin{defi}
Sea $\rho $ el peso específico de un fluido y sea $w(x)$ una función continua sobre $[a,b]$ que describe el ancho de una placa plana sumergida verticalmente a una profundidad $x$. La fuerza $F$ ejercida por el fluido sobre un lado de la placa sumergida es
\[
F=\int_{a}^{b}\rho xw(x)dx
\]
la integral de $\rho \times \text{profundidad}\times \text{ancho}$ horizontal de la placa en un punto \textit{$x\in \left[ a,b\right] $.}
\end{defi}

\begin{ejemplo}
Una placa en forma de triángulo equilátero de $6$ pies de lado y con la base horizontal se sumerge verticalmente en agua, hasta que la base del triá%
ngulo coincida con el nivel del agua. Encuentre la fuerza ejercida por el
agua sobre un lado de la placa.
\end{ejemplo}

\emph{Solución.} Ubicamos un sistema de coordenadas como en la figura \ref{fig:placa_triangular} para usar la simetría del triángulo.

\begin{figure}[H]
\centering
\caption{Placa triangular sumergida}
\label{fig:placa_triangular}
\includegraphics[width=0.38\textwidth]{Gráficas/PlacaTriangularSumerg.pdf}

\propia
\end{figure}

%\begin{center}----------
%\centerline{ \graf{ Placa triangular sumergida}}
%\includegraphics[scale=0.4]{PlacaTriangularSumerg.pdf}
%\textbf{Fuente:} elaboración propia
%\end{center}

La altura de un tríangulo equilátero de lado $L$ es $\frac{\sqrt{3}}{2}%
L,$ así que la altura de este triángulo es $h=3\sqrt{3}$. La placa
queda sumergida en el intervalo $\left[ 0,3\sqrt{3}\right] .$ \ Tome un
segmento horizontal en $x\in \left[ 0,3\sqrt{3}\right]$ sobre la placa
(rectángulo de espesor infinitesimal); le medimos la profundidad $h=x$ y el ancho $w=2y$ con triángulos semejantes: $\frac{w}{6}=\frac{h-x}{h}$. Entonces {\footnotesize $w=\frac{6\left( 3\sqrt{3}-x\right) }{3  \sqrt{3}}=6-\frac{2\sqrt{3}}{3}x$.}

La fuerza ejercida por el agua sobre la placa es

\resizebox{\textwidth}{!}{$
\begin{aligned}
F&=\int_{0}^{3\sqrt{3}}\rho hw\left( x\right) dx=\int_{0}^{3\sqrt{3}}62,4x\left( 6-\frac{2\sqrt{3}}{3}x\right) dx=62,4\int_{0}^{3\sqrt{3}}\left( 6x-\frac{2\sqrt{3}}{3}x^{2}\right) dx\\
&=62,4\left( 3x^{2}-\frac{2\sqrt{3}}{9}x^{3}\right) _{0}^{3\sqrt{3}}=62,4\left( 81-54\right)\approx1684,8\text{\,libras}
\end{aligned}
$}

\begin{ejemplo}
La placa formada por el sector parabólico limitado por $y=x^{2}$ y la recta $y=4$ ($x$ en pies) se sumerge verticalmente en agua, hasta que la recta $y=4$ coincida con el nivel del agua. Encuentre la fuerza ejercida por el agua sobre un lado de la placa.
\end{ejemplo}

\emph{Solución.} Dibujamos la placa parabólica de forma normal: el eje $x$ horizontal y el
eje $y$ vertical:

\begin{figure}[H]
\centering
\caption{Placa parabólica sumergida}
\includegraphics[width=0.38\textwidth]{Gráficas/PlacaParabolSumerg.pdf}

\propia
\end{figure}

%\begin{center}----------
%\centerline{ \graf{ Placa parabólica sumergida}}
%\includegraphics[scale=0.3]{PlacaParabolSumerg.pdf}
%\textbf{Fuente:} elaboración propia
%\end{center}

El intervalo de integración es $\left[ 0,4\right] $ en el eje $y$, para cada $y\in \left[ 0,4\right] $, dibujamos una placa $P_{y},$ infinitesimal de ancho $w$ y altura $dy.$ La profundidad de la placa $P_{y}$ es $4-y.$

Si $y=x^{2},$ entonces $x=\sqrt{y}$ y el $ancho$ de la placa es $w=2x=2\sqrt{%
y}$. La fuerza ejercida por el agua sobre un lado de la placa es
{\small
\begin{align*}
F&=\int_{0}^{4}62,4\ast ancho\ast \text{profundidad}\ast dy=\int_{0}^{4}62,4\left( 2\sqrt{y}\right) \left( 4-y\right) dy\\
&=124,8\int_{0}^{4}\left( 4y^{\frac{1}{2}}-y^{\frac{3}{2}}\right) dy=124,8\left( \frac{8}{3}y^{\frac{3}{2}}-\frac{2}{5}y^{\frac{5}{2}}\right)
_{0}^{4}\\
&=124,8\left( \frac{64}{3}-\frac{64}{5}\right)=1064,96\text { libras}
\end{align*}
}

\section{Centro de masa y momentos de inercia}

\textbf{Densidad y masa:} la densidad es la concentración de masa en un objeto (tridimensional) y se mide frecuentemente en unidades de masa por volumen, por ejemplo, g/cm$^{3}$. Un objeto con densidad uniforme (constante) satisface la relación
\[
\text{Masa} = \text{densidad} \times \text{volumen}
\]
Si el objeto es bidimensional ---por ejemplo, una placa o lámina plana--- y la
densidad es constante, entonces
\[
\text{Masa} = \text{densidad} \times \text{área}
\]
Si el objeto es unidimensional ---como, por ejemplo, una barra, varilla o alambre
delgado--- y la densidad es constante, entonces
\[
\text{Masa} = \text{densidad} \times \text{longitud}
\]
Cuando la densidad de un objeto es variable, estas fórmulas no se cumplen.

Para encontrar, entonces, un valor de la masa, usamos herramientas que nos da el cálculo.

\subsection*{Masa de objetos unidimensionales}

Encontremos la masa de objetos unidimensionales, barras o
varillas delgadas, que se representan como intervalos cerrados sobre el eje $x$.

\begin{figure}[H]
\centering
\caption{Barra o varilla delgada}
\includegraphics[width=0.40\textwidth]{Gráficas/MasaBarra.pdf}

\propia
\end{figure}

%\begin{center}----------
%\centerline{ \graf{ Barra o varilla delgada}}
%\includegraphics[scale=0.35]{MasaBarra.pdf}
%\textbf{Fuente:} elaboración propia
%\end{center}

Supongamos que la barra de la figura 142 de arriba tiene una densidad variable $%
\rho \left( x\right) $ en cada punto $x$ del intervalo $\left[ a,b\right] $.

Para un objeto unidimensional, usamos la densidad lineal con unidades de masa por longitud, por ejemplo, gramos por centímetro. Para saber cuál es la
masa de esta varilla, la dividimos por medio de una partición regular del
intervalo $\left[ a,b\right]$:

\begin{figure}[H]
\centering
\caption{Partición en barra delgada}
\begin{minipage}[t]{0.38\textwidth}
\centering
\includegraphics[width=\linewidth]{Gráficas/MasaBarraParticion.pdf}
\end{minipage}
\hspace{1em}
\begin{minipage}[t]{0.38\textwidth}
\centering

\includegraphics[width=\linewidth]{Gráficas/MasaBarraParticion2.pdf}
\end{minipage}

\propia
\end{figure}

%\begin{center}------------
%\centerline{ \graf{ Partición en barra delgada}}
%\includegraphics[scale=0.35]{MasaBarraParticion.pdf}
%\textbf{Fuente:} elaboración propia
%\end{center}

%\begin{center}------------
%\centerline{ \graf{ Partición en barra delgada (subintervalos)}}
%\includegraphics[scale=0.35]{MasaBarraParticion2.pdf}
%\textbf{Fuente:} elaboración propia
%\end{center}

Comenzamos dividiendo la barra, representada por el intervalo $\left[ a,b%
\right] $, en $n$ subintervalos de igual longitud $\triangle x=\frac{b-a}{n}$, es decir, $\left\{ x_{k}\right\} _{k=0}^{n},$ $x_{k}=a+k\triangle x$ \ $0\leq k\leq n$.

Así, la barra $B$ es la unión de las $n$ barras $B=\cup _{k=1}^{n}$B$_{k}$, la \textit{k}-%
ésima barra B$_k$ está representada por el k-ésimo
subintervalo $\left[ x_{k-1},x_{k}\right] $; por lo tanto, la masa de la \textit{k}-ésima barra la podemos aproximar por $m_{k}\approx \rho \left( x_{k}\right)
\triangle x$, densidad constante por longitud, por lo que la masa de toda la barra la aproximamos por

\[
m=\sum_{k=1}^{n}m_{k}\approx \sum_{k=1}^{n}\rho \left( x_{k}\right)\Delta x
\]

El valor de la masa de la barra es cada vez más exacta cuando $n$ es grande,
así que
\[
m=\lim_{n\rightarrow \infty }\sum_{k=1}^{n}\rho \left( x_{k}\right)
\triangle x
\]
Con ello, y por la definición de integral,

\begin{equation}\label{IntMasa}
m=\int_{a}^{b}\rho \left( x\right) dx
\end{equation}

\begin{ejemplo} Halle la masa de una barra delgada que se encuentra sobre el eje $x$
en el intervalo $\left[ 0,4\right] ,$ si la función de densidad de la
barra es {\small $\rho \left( x\right) =x\sqrt{16-x^{2}},$}
\end{ejemplo}

\emph{Solución.} Por \eqref{IntMasa},
\[
\text{Masa}=m=\int_{0}^{4}\rho \left( x\right) dx=\int_{0}^{4}x\sqrt{16-x^{2}}dx
\]
Si hacemos la sustitución {\small $z=\sqrt{16-x^{2}}$}, tenemos que $z^{2}=16-x^{2}$ y $xdx=-zdz$, entonces
\begin{align*}
m&=\int_{0}^{4}x\sqrt{16-x^{2}}dx=\int_{4}^{0}z\left( -zdz\right)=\int_{0}^{4}z^{2}dz=\frac{1}{3}z^{3}\Big\vert _{0}^{4} =\frac{64}{3}
\end{align*}

\subsection*{Masa de una lámina}

Sea $R$ la región acotada por las curvas $y=f(x)$, $y=g(x)$, $x=a$ \ y \
$x=b$, \ donde $g(x)\leq f(x)$ \ en \ $a\leq x\leq b.$ \

Físicamente, $R$ representa una lámina delgada. Supongamos que esta lámina
tiene densidad constante $\rho $, entonces su masa está dada por
\begin{equation}\label{MasaLamin}
m=\text{densidad}\times \text{área}=\rho \int_{a}^{b}\left( f\left( x\right)
-g\left( x\right) \right) dx
\end{equation}

Supongamos que la densidad de la lámina $R$ es variable y solo depende
de $x$, $\rho =\rho \left( x\right).$

Comenzamos dividiendo la lámina, representada por la región $R$. Sea $%
P=\left\{ x_{k}\right\} _{k=0}^{n}$ una partición regular del intervalo $%
\left[ a,b\right] $, en $n$ subintervalos de igual longitud $\Delta x=\frac{%
b-a}{n},$ $x_{k}=a+k\Delta x$, \ $0\leq k\leq n$:

\begin{figure}[H]
\centering
\caption{Lámina delgada}
\includegraphics[width=0.40\textwidth]{Gráficas/MasaLamina.pdf}

\propia
\end{figure}

%\begin{center}-----------
%\centerline{ \graf{ Lámina delgada}}
%\includegraphics[scale=0.35]{MasaLamina.pdf}
%\textbf{Fuente:} elaboración propia
%\end{center}

\begin{figure}[H]
\centering
\caption{Rectángulos en lámina delgada}
\includegraphics[width=0.48\textwidth]{Gráficas/MasaLaminaPartic.pdf}

\propia
\end{figure}

%\begin{center}------------
%\centerline{ \graf{Rectángulos en lámina delgada}}
%\includegraphics[scale=0.35]{MasaLaminaPartic.pdf}
%\textbf{Fuente:} elaboración propia
%\end{center}

Así, la lámina $R$ se aproxima de las $n$ láminas $R\approx \cup
_{k=1}^{n}R_{k}$; la \textit{k}-ésima lámina R$_{k\text{ }}$ está representada por
el \textit{k}-ésimo rectángulo, cuya base es el subintervalo $\left[
x_{k-1},x_{k}\right] $; así, la masa de la \textit{k}-ésima barra la podemos
aproximar por $m_{k}\approx \rho \left( x_{k}\right) \left( f\left(
x_{k}\right) -g\left( x_{k}\right) \right) \triangle x$ densidad constante\
por el área.

Así que \ la masa de toda la barra la aproximamos por

\[
m=\sum_{k=1}^{n}m_{k}\approx \sum_{k=1}^{n}\rho \left( x_{k}\right) \left(
f\left( x_{k}\right) -g\left( x_{k}\right) \right) \triangle x
\]
El valor de la masa de la lámina es cada vez más exacto cuando $n$ es grande,
así que
\[
m=\lim_{n\rightarrow \infty }\sum_{k=1}^{n}\rho \left( x_{k}\right) \left(
f\left( x_{k}\right) -g\left( x_{k}\right) \right) \triangle x
\]
Por definición de integral,
\[
m=\int_{a}^{b}\rho \left( x\right) \left( f\left( x\right) -g\left(
x\right) \right) dx
\]
\begin{ejemplo}
Halle la masa de la lámina que se encuentra sobre la región
limitada por $y=9-x^{2},y=x+3,$ si la función de densidad es $\rho
\left( x\right) =x^{2}$.
\end{ejemplo}
\emph{Solución.} Dibujamos la lámina:

\begin{figure}[H]
\centering
\caption{Lámina del ejercicio}
\includegraphics[width=0.38\textwidth]{Gráficas/MasaLamParab.pdf}

\propia
\end{figure}

%\begin{center}-----------
%\centerline{ \graf{Lámina del ejercicio}}
%\includegraphics[scale=0.35]{MasaLamParab.pdf}
%\textbf{Fuente:} elaboración propia
%\end{center}

Hallamos los puntos de corte de la parábola y la recta igualando sus ecuaciones:
\begin{align*}
9-x^{2}&=x+3\\
x^{2}+x-6&=0\\
\left( x+3\right) \left(x-2\right) &=0\\
x&=-3\ x=2
\end{align*}
Con ello,

{\small
\begin{align*}
\text{Masa}&=\int_{-3}^{2}\rho \left( x\right) \left( f\left( x\right) -g\left(
x\right) \right) dx\\
&=\int_{-3}^{2}x^{2}\left( 9-x^{2}-\left( x+3\right)
\right) dx=\int_{-3}^{2}x^{2}\left( 6-x-x^{2}\right) dx\\
&=\int_{-3}^{2}\left(6x^{2}-x^{3}-x^{4}\right) dx=\left( 2x^{3}-\frac{1}{4}x^{4}-\frac{1}{5}x^{5}\right)\Big|_{-3}^{2}\\
&=2\left( 2\right) ^{3}-\frac{1}{4}\left( 2\right) ^{4}-\frac{1}{5}\left(
2\right) ^{5}-\left( 2\left( -3\right) ^{3}-\frac{1}{4}\left( -3\right) ^{4}-%
\frac{1}{5}\left( -3\right) ^{5}\right)\\
&=16-4-\frac{32}{5}-\left( -54-\frac{81}{4}+\frac{243}{5}\right) \\
&=16-4-\frac{32}{5}+54+\frac{81}{4}-\frac{243}{5}\\
&=66-\frac{275}{5}+\frac{81}{4}\\
&=\frac{1320-1100+405}{20}=\frac{625}{20}=\frac{125}{4}
\end{align*}
}

\subsection*{Centro de masa para masas puntuales unidimensionales}

Supongamos que se tienen dos partículas unidimensionales con masas $m_{1}$ y $m_{2}$,
ubicadas sobre el eje $x$ y en los puntos con coordenadas $x_1$ y $x_{2}$. El punto $\overline{x}$ sobre el eje $x$ donde se equilibran esas dos masas se llama el centro de masa del sistema.

\begin{figure}[H]
\centering
\caption{Centro de masa}
\includegraphics[width=0.40\textwidth]{Gráficas/CentMasaUnid.pdf}

\propia
\end{figure}

%\begin{center}------------
%\centerline{ \graf{Centro de masa}}
%\includegraphics[scale=0.35]{CentMasaUnid.pdf}
%\textbf{Fuente:} elaboración propia
%\end{center}

El centro de masa $\overline{x}$ cumple $m_{1}d_{1}=m_{2}d_{2}$, \ así
que
{\small
\begin{align*}
m_{1}\left( \overline{x}-x_{1}\right) &=m_{2}\left( x_{2}-\overline{x}%
\right) \\
m_{1}\left( \overline{x}-x_{1}\right) +m_{2}\left(
\overline{x}-x_{2}\right) &=0\\
\overline{x}\left( m_{1}+m_{2}\right)
&=m_{1}x_{1}+m_{2}x_{2}\\
\overline{x}&=\frac{m_{1}x_{1}+m_{2}x_{2}}{  m_{1}+m_{2}}
\end{align*}
}

Generalizando a un sistema de $n$ masas puntuales, $m_1,\ m_2,\ldots,m_n$, ubicadas en puntos del eje
$x$ con coordenadas $x_1,\ x_2,\ldots, x_{n}$, tenemos que el centro de masa se
obtiene de la igualdad
{\small
\begin{align*}
\sum_{k=1}^{n}m_{k}\left( x_{k}-\overline{x}\right) &=0\\
\sum_{k=1}^{n}m_{k}x_{k}-\sum_{k=1}^{n}m_{k}\overline{x}&=0 \\
\overline{x}\sum_{k=1}^{n}m_{k}&=\sum_{k=1}^{n}m_{k}x_{k}
\end{align*}
}
%
o también

\[
\overline{x}=\frac{\sum_{k=1}^{n}m_{k}x_{k}}{\sum_{k=1}^{n}m_{k}}=\frac{%
M_{0}}{m}
\]

donde $M_{0}$ es el momento de inercia en torno al origen del sistema de masas puntuales y $m$ es la masa total (la suma de las masas puntuales).

\begin{ejemplo}
Halle el centro de masa del sistema de masas $m_{1}=7\,\text{kg}$, $m_{2}= 4\,\text{kg}$ y $m_{3}= 5\,\text{kg}$, ubicadas respectivamente en $x_{1}=-2,\ x_{2}=3$ y $x_{3}=5$.
\end{ejemplo}

\emph{Solución.} Usando la fórmula anterior,
\[
\overline{x}=\frac{m_{1}x_{1}+m_{2}x_{2}+m_{3}x_{3}}{%
m_{1}+m_{2}+m_{3}}=\frac{7\left( -2\right) +4\times 3+5\times 5}{7+4+5}=\frac{23}{16}
\]

\subsection*{Centro de masa para varillas, alambres o barras delgadas}

Supongamos que tenemos una barra delgada sobre el intervalo $\left[ a,b\right] $ y
en el eje $x$ con densidad variable $\rho =\rho \left( x\right) $. Entonces,
\[
\text{Masa}=m=\int_{a}^{b}\rho \left( x\right) dx,\text{ y } \overline{x}=\frac{M}{m}
\]
con $M$ el momento respecto al origen,
\[
M=\int_{a}^{b}x\rho \left( x\right) dx \text{ , y entonces, } \overline{x}=\frac{\int_{a}^{b}x\rho \left( x\right) dx}{\int_{a}^{b}\rho  \left( x\right) dx}
\]

\begin{ejemplo}
Halle el centro de masa de una barra delgada que se encuentra
sobre el eje $x$ en el intervalo $\left[ 0,5\right] $ y que tiene funció%
n de densidad variable $\rho \left( x\right) =\sqrt{4+x}$.
\end{ejemplo}

\emph{Solución.} Las integrales que se tienen en este ejemplo pueden calcularse sin mucha dificultad, por eso no se detallan acá:
\[
m=\int_{a}^{b}\rho \left( x\right) dx=\int_{0}^{5}\sqrt{4+x}dx
\]
tomando $u=\sqrt{4+x}\Rightarrow u^{2}=4+x\Rightarrow 2udu=dx$

si $x=0\Rightarrow u=2$

si $x=5\Rightarrow u=3$

$m=\int_{2}^{3}2u^{2}du=\left( \frac{2u^{3}}{3}\right) _{2}^{3}=\frac{2}{3}%
\left( 3^{3}-2^{3}\right) =\frac{2}{3}\left( 27-8\right) =\frac{2}{3}\left(
19\right) =\frac{38}{3}$
\[
M_{0}=\int_{0}^{5}x\rho \left( x\right) dx=\int_{0}^{5}x\sqrt{4+x}dx
\]
tomando $u=\sqrt{4+x}\Rightarrow u^{2}=4+x\Rightarrow 2udu=dx$

$x=u^{2}-4$

si $x=0\Rightarrow u=2$

si $x=5\Rightarrow u=3$

$M_{0}=\int_{2}^{3}\left( u^{2}-4\right) 2u^{2}du=2\int_{2}^{3}\left(
u^{4}-4u^{2}\right) du=2\left( \frac{u^{5}}{5}-\frac{4u^{3}}{3}\right)
_{2}^{3}$

$=2\left( \frac{3^{5}}{5}-\frac{4\left( 3\right) ^{3}}{3}-\left( \frac{2^{5}%
}{5}-\frac{4\left( 2\right) ^{3}}{3}\right) \right) $

$=2\left( \frac{243}{5}-\frac{108}{3}-\frac{32}{5}+\frac{32}{3}\right) $

$=2\left( \frac{211}{5}-\frac{76}{3}\right) $ $=2\left( \frac{633-380}{15}%
\right) =2\left( \frac{253}{15}\right) =\frac{506}{15}$\\
y con ello,

\[
\overline{x}=\frac{\frac{506}{15}}{\frac{38}{3}}=\allowbreak \frac{253}{95}
\]

\begin{ejemplo}
Halle el centro de masa de una barra delgada que se encuentra sobre el eje $x$ en el intervalo $\left[ 1,4\right] $ y que tiene función
de densidad variable $\rho \left( x\right) =\sqrt{1+3x^{2}}$.
\end{ejemplo}

\emph{Solución.} Con la información anterior,
\begin{align*}
m=\int_{1}^{4}\sqrt{1+3x^{2}}dx
\end{align*}
si hacemos $\sqrt{3}x=\tan y$ para $-\frac{\pi }{2}<y<\frac{\pi }{2}$, entonces $\sqrt{3}dx=\sec ^{2}ydy$ y $dx=\cfrac{\sec ^{2}ydy}{\sqrt{3}}$; luego, para $1\leq x\leq 4$,

\resizebox{\textwidth}{!}{$
\begin{aligned}
m&=\int\sqrt{1+3x^{2}}dx
=\left( \int \sqrt{1+\tan ^{2}y}\frac{\sec ^{2}ydy}{\sqrt{3}}\right)\Bigg\vert
_{1}^{4}=\left( \int \sqrt{\sec ^{2}y}\frac{\sec ^{2}ydy}{\sqrt{3}}\right)\Bigg\vert
_{1}^{4}\\
&=\left( \int \frac{\sec ^{3}ydy}{\sqrt{3}}\right)\Bigg\vert _{1}^{4}=\left(
\frac{1}{\sqrt{3}}\int \sec ^{3}ydy\right)\Bigg\vert _{1}^{4}
=\left( \frac{1}{\sqrt{3}}\left( \frac{\sec y\tan y+\ln \left\vert \sec y+\tan y\right\vert }{2}\right) \right)\Bigg\vert_{1}^{4}
\end{aligned}
$}

y usando el triángulo:

\begin{figure}[H]
\centering
\caption{Sustitución para la integral}
\includegraphics[width=0.35\textwidth]{Gráficas/TriangSustTang.pdf}

\propia
\end{figure}

%\begin{center}----------
%\centerline{ \graf{Sustitución para la integral}}
%\includegraphics[scale=0.3]{TriangSustTang.pdf}
%\textbf{Fuente:} elaboración propia
%\end{center}

\resizebox{\textwidth}{!}{$
\begin{aligned}
m&=\left( \frac{1}{\sqrt{3}}\left( \frac{\sec y\tan y+\ln \left\vert \sec
y+\tan y\right\vert }{2}\right) \right)\Bigg\vert_{1}^{4}\\
&=\left( \frac{1}{\sqrt{3}}\left( \frac{\sqrt{1+3x^{2}}\sqrt{3}x+\ln
\left\vert \sqrt{1+3x^{2}}+\sqrt{3}x\right\vert }{2}\right) \right)\Bigg\vert _{1}^{4}\\
&=\left( \frac{1}{2\sqrt{3}}\left( \sqrt{1+3x^{2}}\sqrt{3}x+\ln \left\vert
\sqrt{1+3x^{2}}+\sqrt{3}x\right\vert \right) \right)\Bigg\vert _{1}^{4}\\
&=\frac{1}{2\sqrt{3}}\left( \sqrt{1+3\left( 4\right) ^{2}}\sqrt{3}\left(
  4\right) +\ln \left\vert \sqrt{1+3\left( 4\right) ^{2}}+\sqrt{3}\left(
  4\right) \right\vert\right.\\
  &\left. -\sqrt{1+3\left( 1\right) ^{2}}\sqrt{3}\left( 1\right)
  -\ln \left\vert \sqrt{1+3\left( 1\right) ^{2}}+\sqrt{3}\left( 1\right)
\right\vert \right)\\
&=\frac{1}{2\sqrt{3}}\left( \sqrt{1+3\left( 16\right) ^{{}}}\sqrt{3}\left(
  4\right) +\ln \left\vert \sqrt{1+3\left( 16\right) ^{{}}}+\sqrt{3}\left(
  4\right) \right\vert -\sqrt{4}\sqrt{3}-\ln \left\vert \sqrt{4}+\sqrt{3}%
\right\vert \right) %\\
\end{aligned}
$}

{\small
\begin{align*}
&=\frac{1}{2\sqrt{3}}\left( \sqrt{49^{{}}}\sqrt{3}\left( 4\right) +\ln
  \left\vert \sqrt{49^{{}}}+\sqrt{3}\left( 4\right) \right\vert -2\sqrt{3}-\ln
\left\vert 2+\sqrt{3}\right\vert \right) \\
&=\frac{1}{2\sqrt{3}}\left( 7\sqrt{3}\left( 4\right) +\ln \left\vert 7+\sqrt{%
  3}\left( 4\right) \right\vert -2\sqrt{3}-\ln \left\vert 2+\sqrt{3}%
\right\vert \right) \\
&=\frac{1}{2\sqrt{3}}\left( 28\sqrt{3}+\ln \left\vert 7+4\sqrt{3}\right\vert
-2\sqrt{3}-\ln \left\vert 2+\sqrt{3}\right\vert \right)\\
& =\frac{1}{2\sqrt{3}}%
\left( 26\sqrt{3}+\ln \left\vert \frac{7+4\sqrt{3}}{2+\sqrt{3}}\right\vert
\right) =\frac{1}{2\sqrt{3}}\left( 26\sqrt{3}+\ln \left( \frac{7+4\sqrt{3}}{%
2+\sqrt{3}}\right) \right)
\end{align*}
}
entonces,
\[
m=\int_{1}^{4}\sqrt{1+3x^{2}}dx\approx13,\,38
\]
y
\[
M_{0}=\int_{1}^{4}x\rho \left( x\right) dx=\int_{1}^{4}x\sqrt{1+3x^{2}}dx
\]
y para $z=\sqrt{1+3x^{2}}$, $z^{2}=1+3x^{2}$, $2zdz=6xdx$ y $xdx=\frac{1}{3}zdz$; luego, por el cambio de límites, si $x=1$, entonces $z=\sqrt{4}=2$, y, si $x=4$, entonces $z=\sqrt{49}=7$. Con ello,
\[
M_{0}=\int_{2}^{7}z\frac{1}{3}zdz=\left( \frac{1}{9}z^{3}\right) _{2}^{7}=%
\frac{1}{9}\left( 7^{3}-8\right) =\frac{335}{9}
\]
y
\[
\overline{x}=\frac{\int_{a}^{b}x\rho \left( x\right) dx}{\int_{a}^{b}\rho \left( x\right)
dx}=\frac{\frac{335}{9}}{13\approx 38}\approx 2,7819
\]

\subsection*{Centro de masa para láminas planas formadas por regiones limitadas por dos funciones que dependen de $x$}

Sea \ $\mathcal{R=}\left\{ \left( x,y\right) :a\leq x\leq b,g\left( x\right)
\leq y\leq f\left( x\right) \right\}$ una lámina delgada con densidad
variable $\rho =\rho \left( x\right) $ que depende solo de $x.$

\begin{figure}[H]
\centering
\caption{Lámina delgada}
\includegraphics[width=0.38\textwidth]{Gráficas/MasaLamina.pdf}

\propia
\end{figure}

%\begin{center}----------
%\centerline{ \graf{Lámina delgada}}
%\includegraphics[scale=0.35]{MasaLamina.pdf}
%\textbf{Fuente:} elaboración propia
%\end{center}

La masa de la lámina es
\[
m=\int_{a}^{b}\rho \left( x\right) \left( f\left( x\right) -g\left(
x\right) \right) dx
\]
y las coordenadas del centro de masa son
\[
\overline{x}=\frac{M_{Y}}{m},\ \overline{y}=\frac{M_{X}}{m}
\]
donde
\[
M_{X}=\int_{a}^{b}\frac{\rho \left( x\right) }{2}\left( \left(
f\left( x\right) \right) ^{2}-\left( g\left( x\right) \right) ^{2}\right)
dx
\]
es el momento de la lámina respecto al eje $x$, y
\[
M_{Y}=\int_{a}^{b}x\rho \left( x\right) \left( f\left( x\right) -g\left(
x\right) \right) dx
\]
es momento de la lámina respecto al eje $y$.

El centro de masa de la lámina es $\left( \overline{x},\overline{y}\right).$ Si la densidad
es constante, el centro de masa se denomina \textit{centroide} y solo depende de la geometría de la región.

\begin{ejemplo}
Halle el centro de masa de la placa delgada limitada por $y=9-x^{2}$ y $y=8x^{2}$ con función de densidad $\rho \left( x\right)=x^{2}. $
\end{ejemplo}

\emph{Solución.} Para los puntos de corte, hacemos $9-x^{2}=8x^{2}$, de donde $ x^{2}=1$ y $x=\pm 1.$ La figura de la lámina se muestra a continuación:

\begin{figure}[H]
\centering
\caption{Lámina delgada formada por dos parábolas}
\includegraphics[width=0.38\textwidth]{Gráficas/CentrMasaLamin2Parab.pdf}

\propia
\end{figure}

%\begin{center}------------
%\centerline{ \graf{Lámina delgada formada por dos parábolas}}
%\includegraphics[scale=0.35]{CentrMasaLamin2Parab.pdf}
%\textbf{Fuente:} elaboración propia
%\end{center}

Entonces,

\resizebox{\textwidth}{!}{$
\begin{aligned}
m&=\int_{-1}^{1}x^{2}\left( 9-x^{2}-8x^{2}\right) dx=2\int_{0}^{1}\left(
9x^{2}-9x^{4}\right) dx=18\left( \frac{x^{3}}{3}-\frac{x^{5}}{5}\right)
_{0}^{1}=18\left( \frac{1}{3}-\frac{1}{5}\right)=\frac{12}{5}
\end{aligned}
$}

y también

\begin{align*}
M_{y}&=\int_{-1}^{1}\underset{\text{función impar}}{\underbrace{x\left(
x^{2}\left( 9-x^{2}-8x^{2}\right) \right) }}dx=0
\end{align*}

y

\resizebox{\textwidth}{!}{$
\begin{aligned}
M_{x}&=\int_{-1}^{1}\underset{\text{función par}}{\frac{x^{2}}{2}\left(
\left( 9-x^{2}\right) ^{2}-\left( 8x^{2}\right) ^{2}\right) }dx=2\times \frac{1}{2}\int_{0}^{1}x^{2}\left( 81-18x^{2}+x^{4}-64x^{4}\right)dx\\
&=\int_{0}^{1}\left( 81x^{2}-18x^{4}-63x^{6}\right) dx=\left( 27x^{3}-\frac{18}{5}x^{5}-9x^{7}\right)\Big\vert _{0}^{1}\\
&=27-\frac{18}{5}-9=\frac{72}{5}
\end{aligned}
$}

Por otro lado,
\begin{align*}
\overline{x}&=\frac{M_{y}}{m}=\frac{0}{\frac{12}{5}}=0,\text{ y }
\end{align*}

y

\begin{align*}
\overline{y}&=\frac{M_{x}}{m}=\frac{\frac{72}{5}}{\frac{12}{5}}=6
\end{align*}
De allí, el centro de masa es $\left( \overline{x},\overline{y}\right) =\left( 0,6\right) $.

%\begin{center}------------
%\centerline{ \graf{$Centro de masa de la lámina$}}
%\includegraphics[scale=0.35]{CentrMasaLamin2Parab2.pdf}
%\textbf{Fuente:} elaboración propia
%\end{center}

\begin{ejemplo}
Halle el centro de masa de la placa delgada limitada por $y=x^{2},y=4x$, con función de densidad $\rho \left( x\right) =x^{2}.$
\end{ejemplo}

\emph{Solución.} Calculemos las integrales

\resizebox{\textwidth}{!}{$
\begin{aligned}
m&=\int_{0}^{4}\rho \left( x\right) \left( \text{función parte superior}-\text{función parte inferior}\right) dx\\
&=\int_{0}^{4}x^{2}\left( 4x-x^{2}\right) dx=\int_{0}^{4}\left( 4x^{3}-x^{4}\right) dx=\left( x^{4}-\frac{1}{5}x^{5}\right)\Bigg|_{0}^{4}=4^{4}-\frac{4^{5}}{5}\\
&=\frac{256}{5}
\end{aligned}
$}

Los momentos de la lámina son

\resizebox{\textwidth}{!}{$
\begin{aligned}
M_{y}&=\int_{0}^{4}xx^{2}\left( 4x-x^{2}\right) dx=\int_{0}^{4}\left( 4x^{4}-x^{5}\right) dx=\left( \frac{4x^{5}}{5}-\frac{x^{6}}{6}\right)\biggl| _{0}^{4}=\frac{4^{6}}{5}-\frac{4^{6}}{6}=\frac{2048}{15}
\end{aligned}
$}

y

\resizebox{\textwidth}{!}{$
\begin{aligned}
M_{x}&=\int_{a}^{b}\frac{\rho \left( x\right) }{2}\left( (\text{función parte superior})^2-(\text{función parte inferior})^2\right) dx\\
&=\int_{0}^{4}\frac{x^{2}}{2}\left( 16x^{2}-x^{4}\right)dx=\int_{0}^{4}\left( 8x^{4}-\frac{1}{2}x^{6}\right) dx=\left( \frac{8x^{5}}{5%
}-\frac{x^{7}}{14}\right) _{0}^{4}=\frac{8\ast 4^{5}}{5}-\frac{4^{7}}{14}=\frac{16\,384}{35}
\end{aligned}
$}

y las coordenadas del centro de masa son
\[
\overline{x}=\frac{M_{y}}{m}=\frac{\frac{2048}{15}}{\frac{256}{5}}=\frac{8}{%
3}
\]
y
\[
\overline{y}=\frac{M_{x}}{m}=\frac{\frac{16\,384}{35}}{\frac{256}{5}}=\frac{%
64}{7}
\]
Con ello, el centro de masa es $\left( \overline{x},\overline{y}\right) =\left( \frac{8}{3},%
\frac{64}{7}\right) $

\begin{ejemplo}
Halle el centro de masa de la placa delgada limitada por $y=x^{2}$, $y=x$, con función de densidad $\rho \left( x\right) =12x$.
\end{ejemplo}

\emph{Solución.} Tenemos
{\small
\begin{align*}
m&=\int_{a}^{b}\rho \left( x\right) \left( f\left( x\right) -g\left(
x\right) \right) dx=\int_{0}^{1}12x\left( x-x^{2}\right) dx\\
&=\int_{0}^{1}12\left( x^{2}-x^{3}\right) dx=12\left( \frac{x^{3}}{3}-\frac{%
x^{4}}{4}\right) _{0}^{1}\\
&=12\left( \frac{1}{3}-\frac{1}{4}\right) =12\left( \frac{1}{12}\right)=1\\
\end{align*}
}
{\small
\begin{align*}
M_{x}&=\int_{a}^{b}\frac{\rho \left( x\right) }{2}\left( \left( f\left(
x\right) \right) ^{2}-g\left( x\right) ^{2}\right) dx=\int_{0}^{1}6x\left(
x^{2}-x^{4}\right) dx\\
&=\int_{0}^{1}6\left( x^{3}-x^{5}\right) dx=6\left( \frac{x^{4}}{4}-\frac{%
x^{6}}{6}\right) _{0}^{1}\\
&=6\left( \frac{1}{4}-\frac{1}{6}\right) =6\left( \frac{3-2}{12}\right)
=\frac{1}{2}\\
\end{align*}
}
{\small
\begin{align*}
M_{y}&=\int_{a}^{b}x\rho \left( x\right) \left( f\left( x\right) -g\left(
x\right) \right) dx=\int_{0}^{1}12x^{2}\left( x-x^{2}\right) dx\\
&=\int_{0}^{1}12\left( x^{3}-x^{4}\right) dx=12\left( \frac{x^{4}}{4}-\frac{
x^{5}}{5}\right) _{0}^{1}\\
&=12\left( \frac{1}{4}-\frac{1}{5}\right) =12\left(\frac{1}{20}\right)=\frac{3}{5}\\
\overline{x}&=\frac{M_{y}}{m}=\frac{\frac{3}{5}}{1}=\frac{3}{5}\\
\overline{y}&=\frac{M_{x}}{m}=\frac{\frac{1}{2}}{1}=\frac{1}{2}
\end{align*}
}
Y, entonces, el centro de masa es $\left( \frac{3}{5},\frac{1}{2}\right) $

\begin{ejemplo}
Halle el centro de masa de la placa delgada limitada por $y=x^{2},y=x$, con función de densidad $\rho \left( x\right) =x^{2}$.
\end{ejemplo}

\emph{Solución.} Con las funciones dadas, tenemos
\begin{align*}
m&=\int_{a}^{b}\rho \left( x\right) \left( f\left( x\right) -g\left(
x\right) \right) dx=\int_{0}^{1}x^{2}\left( x-x^{2}\right) dx\\
&=\int_{0}^{1}\left( x^{3}-x^{4}\right) dx=\left( \frac{x^{4}}{4}-\frac{x^{5}%
}{5}\right) _{0}^{1}\\
&=\frac{1}{4}-\frac{1}{5}
=\frac{1}{20}\\
\end{align*}

\begin{align*}
M_{x}&=\int_{a}^{b}\frac{\rho \left( x\right) }{2}\left( \left( f\left(
x\right) \right) ^{2}-g\left( x\right) ^{2}\right) dx=\int_{0}^{1}\frac{1}{2}
x^{2}\left( x^{2}-x^{4}\right) dx\\
&=\int_{0}^{1}\frac{1}{2}\left( x^{4}-x^{6}\right) dx=\frac{1}{2}\left( \frac{%
x^{5}}{5}-\frac{x^{7}}{7}\right) _{0}^{1}\\
&=\frac{1}{2}\left( \frac{1}{5}-%
\frac{1}{7}\right) =\frac{1}{2}\left( \frac{2}{35}\right)= \frac{1}{35}\\
\end{align*}

\begin{align*}
M_{y}&=\int_{a}^{b}x\rho \left( x\right) \left( f\left( x\right) -g\left(
x\right) \right) dx=\int_{0}^{1}x\times x^{2}\left( x-x^{2}\right)
dx\\
&=\int_{0}^{1}\left( x^{4}-x^{5}\right) dx=\left( \frac{1}{5}x^{5}-\frac{1}{
6}x^{6}\right) _{0}^{1}=\frac{1}{5}-\frac{1}{6}=\frac{1}{30}\\
\end{align*}

\[
\overline{x}=\frac{M_{y}}{m}=\frac{\frac{1}{30}}{\frac{1}{20}}=\frac{2}{3}\\
\overline{y}=\frac{M_{x}}{m}=\frac{\frac{1}{35}}{\frac{1}{20}}=\frac{20}{35}%
=\frac{4}{7}
\]

luego, el centro de masa es $\left( \frac{1}{35},\frac{1}{30}\right) .$
\clearpage
\addcontentsline{toc}{section}{Ejercicios del capítulo~\thechapter}

%\group

\begin{ejer}{Para los ejercicios \eqref{areainic} a \eqref{areafinal}, encuentre el área de la región del plano $xy$ acotada por las figuras de las funciones dadas.}
\end{ejer}
{\small
\begin{enumerate}
  %\begin{multicols}{2}

  \label{areainic}
\item {$y=x^{2}+3,\ y=x,\ x=-1$ y $x=1$}{$\frac{20}{3}\ \text{ unidades cuadradas} $}

\item {$y=x^{4},\ y=-x-1,\ x=-2$ y $x=0$}{$\frac{32}{5}\ \text{ unidades cuadradas}$}

\item {$y=x^{2}+3$ y $y=x+5$}{$\frac{9}{2} \ \text{ unidades cuadradas}$}

\item {$y=x$ y $y=x^{2}$}{$\frac{1}{6}\ \text{ unidades cuadradas}$}

\item {$y=x^{2}$ y $ x=y^{2}$}{$\frac{1}{3}\ \text{ unidades cuadradas}$}

\item {$y=4x^{2}$ y $y=x^{2}+3$}{$4 \ \text{ unidades cuadradas}$}

\item {$y=x^{4}-x^{2}$ y $y=1-x^{2}$}{$\frac{8}{5} \ \text{ unidades cuadradas}$}

\item {$x+y^{2}=2$ y $x+y=0$}{$\frac{9}{2} \ \text{ unidades cuadradas}$}

\item {$y=2x-x^{2}$ y $y=x^{3}$}{$\frac{37}{12} \ \text{ unidades cuadradas}$}

\item {$y=\cos x,\ y=\sec ^{2}x,\ x=-\frac{\pi }{4}$ y $x=\frac{\pi}{4}$}{$2-\sqrt{2}\ \text{ unidades cuadradas}$}

\item {$y=\vert x\vert,\ y=(x+1)^{2}-7$ y $x=-4$}{$34\ \text{ unidades cuadradas}$}

\item {$y=\frac{1}{x},\ y=\frac{1}{x^{2}},\ x=1$ y $x=2$}{$\ln 2-\frac{1}{2}\ \text{ unidades cuadradas}$}

\item {$y=\frac{1}{x},\ x=0,\ y=1$ y $y=2$}{$\ln 2\ \text{ unidades cuadradas}$}

\item {$y=1$ y $y=\frac{10}{x^{2}+1}$}{$20\tan^{-1}3-6\ \text{ unidades cuadradas}$}

\item {$y=2^{x},\ y=5^{x},\ x=-1$ y $x=1$}{$\frac{16}{5\ln 5}-\frac{1}{2\ln 2}\ \text{ unidades cuadradas}$}

\item {$y=e^{x},\ y=e^{3x}$ y $x=1$}{$\frac{1}{3}e^{3}-e+\frac{2}{3}\ \text{ unidades cuadradas}$}%\int_{0}^{1}\left( e^{3x}-e^{x}\right) dx$

\item {$y=\sin ^{-1}x,\ y=0$ y $x=0,5$}{$\frac{\pi }{12}+\frac{\sqrt{3}}{2}-1\ \text{ unidades cuadradas}$}

\item {$y=xe^{-x},\ y=0$ y $x=5$}{$1-\frac{6}{e^{5}}\ \text{ unidades cuadradas}$}

\item {$y=5\ln x$ y $y=x\ln x$}{$\frac{25}{2}\ln 5-14\ \text{ unidades cuadradas}$}

\item {$y=\frac{25x^{2}}{x^{2}+9}$ y $y=x^{2}$}{$\frac{472}{3}-150\tan^{-1} \frac{4}{3} \text{ unidades cuadradas}$}

\item {$y=(\sin^{-1}x)^2$ y $y=\sen^{-1} x$}{$2\sin 1-\sqrt{\frac{1}{2}\cos 2+\frac{1}{2}}-1  \text{ unidades cuadradas}$}

\item { $y=x\tan^{-1}x$ y $y=\tan^{-1} x$}{$\frac{1}{2}-\frac{1}{2}\ln 2\ \text{ unidades cuadradas}$}

\item {$y=\frac{2400}{x^{2}-16}$ y $y=\frac{2400}{x^{4}-16}$}{$150\ln \frac{5^4}{3^5}-300\tan^{-1} \frac{1}{2} \text{ unidades cuadradas}$}

\item {$y=2x^{3}-3x^{2}-9x$ y $ y=x^{3}-2x^{2}-3x$}{$\frac{253}{12}\text{ unidades cuadradas}$ }

\item {La región en el primer cuadrante acotada por $y=x^{-2}$, $y=8x$ y $y=\frac{x}{8}$}{$\frac{9}{4}\text{ unidades cuadradas}$}

\item {La región en el primer cuadrante acotada por $y=9-x^{2},\ y=8x$ y $y=\frac{5}{2}x$}{$\frac{17}{3}\text{ unidades cuadradas}$}

\item {Halle el área de la región limitada por la hipérbola dada por la ecuación: $9x^{2}-4y^{2}=36$ y la recta vertical $x=3.$}{$\frac{9\sqrt{5}}{2}-6\ln \frac{3+\sqrt{5}}{2}\text{ unidades cuadradas}$}
  \label{areafinal}
  %\end{multicols}
\end{enumerate}
}

\begin{ejer}{Halle el volumen del sólido cuya base es la región limitada por $%
y=x^{2}$ y $y=4x$ y cuyas secciones transversales perpendiculares al eje $x$
son triángulos rectángulos isósceles con la hipotenusa sobre la
base.
} {$\frac{128}{15} \text{ unidades cúbicas}$}
\end{ejer}

\begin{ejer}{Halle el volumen del sólido cuya base es el disco centrado en $\left(
0,0\right) $ de radio 2 y cuyas secciones transversales perpendiculares al eje
$x$ son triángulos equiláteros con un lado sobre la base.}{$\frac{%
32}{3}\sqrt{3}$ $\text{ unidades cúbicas}$}
\end{ejer}

\begin{ejer}{ Halle el volumen del sólido cuya base es el disco centrado en $\left(
0,0\right) $ de radio 2 y cuyas secciones transversales perpendiculares al eje

$x$ son cuadrados con la diagonal sobre la base.}{$\frac{64}{3}\text{ unidades cúbicas}$}
\end{ejer}

\begin{ejer}{Halle el volumen del sólido cuya base es el disco centrado en $\left(
0,0\right) $ de radio 2 y cuyas secciones transversales perpendiculares al eje
$x$ son cuadrados con un lado sobre la base.} {$\frac{128}{3}\text{ unidades cúbicas}$}
\end{ejer}

\begin{ejer}{Halle el volumen del sólido cuya base es la región limitada por $%
x=y^{2}$ y $x=4$ y cuyas secciones transversales perpendiculares al eje $x$
son discos con el diámetro sobre la base.} {$8\pi \text{ unidades cúbicas}$}
\end{ejer}

\begin{ejer}{Halle el volumen del sólido cuya base es el disco centrado en $\left(
0,0\right) $ de radio 2 y cuyas secciones transversales perpendiculares al
eje $x$ son discos con el diámetro sobre la base.} {$\frac{32}{3}\pi
\text{ unidades cúbicas}$}
\end{ejer}

\begin{ejer}{Halle el volumen del sólido cuya base es la región limitada por $%
y=2\sin x$ y $y=0,0\leq x\leq \pi $, y cuyas secciones transversales perpendiculares al eje $x$
son triángulos equiláteros con un lado sobre la base.} {$\frac{\sqrt{3}%
\pi }{2}\allowbreak \pi \text{ unidades cúbicas}$}
\end{ejer}

\begin{ejer}{Halle el volumen del sólido cuya base es la región limitada por $%
y=\tan x$ y $y=\sec x,x=-\frac{\pi }{4},x=\frac{\pi }{4}$, cuyas secciones
transversales perpendiculares al eje $x$ son rectángulos con el largo
sobre la base y tal que la altura es la mitad del largo.} {$2-\frac{\pi}{4} \text{ unidades cúbicas}$}
\end{ejer}

\begin{ejer}{En los ejercicios \eqref{volinic} a \eqref{volfin}, encuentre el volumen del sólido de revolución que se obtiene al hacer girar la región limitada por las curvas dadas en torno al eje indicado. Dibuje la región y el sólido.}
\end{ejer}

{\small
\begin{enumerate}
  %\begin{multicols}{2}

\item {$y=x^{2},x=1,y=0,$ alrededor del eje $x$}{$V=\frac{\pi}{5}\text{ unidades cúbicas}$}\label{volinic}

\item {$y^{2}=x^{3},x=4,y\geq0,$ alrededor del eje $x$} {$V=64\pi \text{ unidades cúbicas}$}

\item {$y=\sqrt{x-1},x=2,x=5,y=0,$ alrededor del eje $x$} \\ {$V=\frac{15}{2}\pi \text{ unidades cúbicas}$}

\item {$x=y-y^{2},x=0,$ alrededor del eje $y$} {$V=\frac{\pi }{30}\text{ unidades cúbicas}$}

\item {$y=x^{4},y=x^{2},$ alrededor del eje $x$} {$V=\frac{8}{45}\pi \text{ unidades cúbicas}$}

\item {$y=4-x^{2},y=0,$ alrededor del eje $x$} {$V=\frac{512}{15}\pi \text{ unidades cúbicas}$}

\item {$y=x^{2},x=y^{2}$}

\item {Alrededor del eje $x$}{$V=\frac{3}{10}\pi \text{ unidades cúbicas}$}

\item {Alrededor del eje $y$} {$V=\frac{3}{10}\pi \text{ unidades cúbicas}$}

\item {Alrededor de la recta $y=2$}{$V=\frac{31}{30}\pi \text{ unidades cúbicas}$}

\item {$y=\sec x,y=1,x=-1,x=1,$ \ alrededor del eje $x$} {$V=\pi (2\tan (1)-2)\text{ unidades cúbicas}$}

\item {$y=\sin x,y=\cos x,x=0,x=\frac{\pi }{4}$}

\item {En torno al eje $x$}{$V=\frac{\pi }{2}\text{ unidades cúbicas}$}

\item {En torno al eje $y$}{$V=\frac{\pi }{2}\left( \sqrt{2}\pi -4\right)\text{ unidades cúbicas}$}

\item {Eje de giro $x=\frac{\pi }{4}$} {$V=\frac{\pi }{2}\left( 4-\pi \right)
  \text{ unidades cúbicas}$}

\item {Eje \ de rotación $y=1$}{$V=\frac{\pi }{2}\left( 4\sqrt{2}-5\right)\text{ unidades cúbicas}$}

\item {Eje \ de rotación $y=-1$}{$ V=\frac{\pi }{2}\left( 4%
  \sqrt{2}-3\right) \text{ unidades cúbicas}$}

\item {$y=\frac{-1}{x},x=1,y=0,x=3$, \ alrededor del eje $x$} {$V=\frac{2\pi }{3}\text{ unidades cúbicas}$}

\item {$y=\tan x,y=1,x=0,$ alrededor del eje $x$} {$V=\pi (\frac{\pi }{2}-1)\text{ unidades cúbicas}$}

\item {$y=\sqrt{x-1},x=5,y=0,$ alrededor del eje $y$} {$V=\frac{544}{15}\pi \text{ unidades cúbicas}$}

\item {$y=\cos x,x=0,y=0,$ con $0\leq x\leq \frac{\pi }{2}$, alrededor de la recta $y=1$} {$V=2\pi-\frac{1}{4} \pi^2 \text{unidades cúbicas}$} \\ \label{volfin}
  %\end{multicols}
\end{enumerate}
}

%\clearpage

\begin{ejer}{Halle la longitud de la curva dada entre los puntos dados:}
\end{ejer}

{\small
\begin{enumerate}
  %\begin{multicols}{2}

\item {$y=2x+1,\ -1\leq x\leq 3.$} {$L=4\sqrt{5}\text{ unidades}$}

\item {$y^{2}=(x-1)^{3}$ desde $P(1,0)$ a $Q(2,1)$}
  {$L=\frac{1}{27}((13)^{\frac{3}{2}}-8) \text{ unidades}$}

\item {$y=1-x^{\frac{2}{3}}$ desde  $ P(-8,-3)$ a $Q(-1,0)$} {$L=\left( -\frac{13\sqrt{13}}{27}+\frac{80\sqrt{10}}{27}\right)
  \text{ unidades}$}

\item {$9y^{2}=x(x-3)^{2}$ desde $P(1,\frac{2}{3})$ a $ Q(2,\frac{\sqrt{2}}{3})$} {$L=\frac{5\sqrt{2}-4}{3}\text{ unidades}$}

\item {$y=\frac{1}{3}(x^{2}+2)^{\frac{3}{2}}\ 0\leq x\leq 1$} {$L=\frac{4}{3}\text{ unidades}$}

\item {$y=\frac{x^{3}}{6}+\frac{1}{2x}$ \ \ $1\leq x\leq 2$} {$L=\frac{17}{12}\text{ unidades}$}

\item {$y=\frac{x^{4}}{4}+\frac{1}{8x^{2}}$ \ $1\leq x\leq 3$} {$L=\frac{181}{9}\text{ unidades}$}

\item {$y=\ln (\cos x)$ \ $0\leq x\leq \frac{\pi }{4}$} {$L=\ln (\sqrt{2}+1) \text{ unidades}$}

\item {$y=\ln (\sin x)$ \ $\frac{\pi }{6}\leq x\leq \frac{\pi }{3}$} {$L=\ln \left( \frac{\sqrt{3}}{3(2-\sqrt{3})}\right) \text{ unidades}$}

\item {$y=\ln (\frac{e^{x}+1}{e^{x}-1})$ \ $1\leq x\leq 2$
  } {$L=\ln \left( e^{2}+1\right) -1
  \text{ unidades}$}

\item {$\ y=e^{x},0\leq x\leq 1$} {$L=\left( \sqrt{e^{2}+1}+\ln \left( \frac{\sqrt{e^{2}+1}-1}{\sqrt{2}-1}%
  \right) -1-\sqrt{2}\right) \text{ unidades}$}

\item { $y=\ln x,1\leq x\leq \sqrt{3}$} {$L=\left( 2-\ln \left\vert \sqrt{3}(\sqrt{2}-1)\right\vert -\sqrt{2}\right) \text{ unidades}$}

\item {$y=2\sqrt{x},0\leq y\leq 2.$} {$L=\allowbreak \left( \ln \left( \sqrt{2}+1\right) +\sqrt{2}\right) \text{ unidades}$}
  %\end{multicols}
\end{enumerate}
}

\begin{ejer}{Encuentre el área de la superficie que se obtiene al girar la curva dada en torno al eje $x$.}

\exercise{$y=\sqrt{x},4\leq x\leq 9$} {$S=\left( \frac{37\sqrt{37}%
-17\sqrt{17}}{6}\right) \pi \text{ unidades cuadradas}$}

\exercise{$y=x+4,0\leq x\leq 2$} {$ S=20\sqrt{2}\text{ unidades cuadradas}$}

\exercise {$y=\sqrt{4x+4},0\leq x\leq 2$} {$S= \frac{64-16\sqrt{2} }{3}\pi\text{ unidades cuadradas}$}

\exercise {$y=x^{3},0\leq x\leq 2$} {$S=\frac{1}{27}\pi \left( 145\sqrt{145}-1\right) \text{ unidades cuadradas}$}

\exercise {$y=x^{3}+\frac{1}{12x},1\leq x\leq 2$
} {$S= \frac{12\,289}{192}\pi \text{ unidades cuadradas}$}

\exercise {$y=\frac{x^{2}}{4}-\frac{\ln x}{2},1\leq x\leq 4$} {$S=\frac{\pi }{16}\left( 315-128\left( \ln 2\right) -16\ln ^{2}2\right) \\
\text{ unidades cuadradas} $ }

\exercise {$y=\sen x,0\leq x\leq \pi $} {$S=2\pi \left( \sqrt{2}+\ln \left( \sqrt{2}+1\right) \right) \text{ unidades cuadradas}$}

\exercise {$x=\frac{1}{3}(y^{2}+2)^{\frac{3}{2}},1\leq y\leq 2$} { $S= \frac{21}{2}\pi \text{ unidades cuadradas}$}

\exercise {$x=2\ln y,1\leq y\leq 2$} {$S=\pi \left( 4\sqrt{2}+4\ln \left\vert 1+\sqrt{2}\right\vert -\sqrt{5}-4\ln
\frac{1+\sqrt{5}}{2}\right) \\
\text{ unidades cuadradas}$}

\end{ejer}

\begin{ejer}{Encuentre el área de la superficie que se obtiene al girar la curva
$x=\frac{1}{2}y^{4}+\frac{1}{16y^{2}},1\leq y\leq 3$
en torno al eje $y$.} \\ {$S=\frac{2127\,389}{1296}\pi \text{ unidades cuadradas}$}
\end{ejer}

\begin{ejer}{En los siguientes ejercicios, halle la masa de las siguientes barras
delgadas\ representadas en el intervalo dado, con la función de densidad
$\rho \left( x\right) $ indicada.}
\end{ejer}

{\small
\begin{enumerate}
  %\begin{multicols}{2}

\item { $\rho \left( x\right) =1+\cos x;$ $0\leq x\leq \frac{\pi }{2}$}{$\frac{\pi }{2}+1$ unidades de masa}

\item {$\rho \left( x\right) =2+\sen x;$ $0\leq x\leq \pi $}{$2\pi +2$ unidades de masa}

\item {$\rho \left( x\right) =1+x^{4};$ $0\leq x\leq 1$}{$\frac{6}{5}$ unidades de masa}

\item {$\rho \left( x\right) =4+x^{3};$ $0\leq x\leq 2$}{$12$ unidades de masa}

\item {$\rho \left( x\right) =\left\{
      \begin{array}{ccc}
        x^{2} & si & 0\leq x\leq 1 \\
        4-x^{2} & si & 1< x\leq 2%
      \end{array}%
    \right. $}{$2$ unidades de masa}

  \item {$\rho \left( x\right) =\left\{
        \begin{array}{ccc}
          \sen x & si & 0\leq x\leq \pi \\
          2+\cos x & si & \pi < x\leq 2\pi%
        \end{array}%
      \right. $}{$2\pi +2$ unidades de masa}

    \item {$\rho \left( x\right) =e^{x}+e^{2x};$ $0\leq x\leq \ln 2$}{$\frac{5}{2}$ unidades de masa}

    \item {$\rho \left( x\right) =\sen \left( 2x\right) +\tan x;$ $0\leq x\leq \frac{%
      \pi }{4}$}{$\ln \sqrt{2}+\frac{1}{2}$ unidades de masa}

    \item {$\rho \left( x\right) =x\sqrt{16-x^{2}};$ $0\leq x\leq 4$}{$\frac{64}{3}$unidades de masa}

    \item {$\rho \left( x\right) =\frac{x}{\sqrt{x^{2}+16}};$ $0\leq x\leq 3$}{$1$ unidad de masa}

    \item {$\rho \left( x\right) =x^{2}e^{x};$ $0\leq x\leq \ln 3$}{$3\ln^2{3}-6\ln 3+4$ unidad de masa}

    \item {$\rho \left( x\right) =\frac{\sqrt{x^{2}-9}}{x};$ $3\leq x\leq 5$}{$4-3\sec^{-1}\frac{5}{3}$ unidad de masa}

    \item {$\rho \left( x\right) =\frac{x^{2}}{\sqrt{x^{2}+9}};$ $0\leq x\leq 4$}{$10-\frac{9}{2}\ln 3$ unidad de masa}
      %\end{multicols}
  \end{enumerate}
}

%\clearpage

\begin{ejer}{En los ejercicios \eqref{masainic} al \eqref{masafin}, halle la masa y el centro de masa de las siguientes láminas delgadas\ representadas por la región limitada por las curvas o rectas dadas, con la función de densidad $\rho \left(x\right) $ indicada.}
\end{ejer}

{\small
  \begin{enumerate}

    \item {\label{masainic} La región limitada por $y=x^{2}+1,y=2x+1$, con función de densidad $\rho \left( x\right) =x+1$}
      {$m=\frac{8}{3},$ Centro de masa $\left( \frac{11}{10},\frac{14}{5}\right) $}

    \item {La región limitada por $y=x^{2}+1,y=3-x^{2}$, con función de
      densidad $\rho \left( x\right) =2x+1$}{$m=\frac{8}{3},$ Centro de masa $\left( \frac{2}{5},2\right)$}

    \item {La región limitada por $y=2x^{2}+1,y=x^{2}+5$, con función de
      densidad $\rho \left( x\right) =3x+2$} {$m=\frac{64}{3},$ Centro de masa $\left( \frac{6}{5},\frac{21}{5}\right)$}

    \item {La región limitada por $y=x^{4},y=x^{2}$, con función de densidad
      $\rho \left( x\right) =x+3$}{$m=\frac{4}{5},$ Centro de masa $\left( \frac{1}{7},\frac{1}{3}%
      \right)$}

    \item {La región del primer cuadrante acotada por $x^{2}+y^{2}=16,$ \ $\rho
      \left( x\right) =x$}{$m=\frac{64}{3},$ Centro de masa $\left( \frac{3}{4}\pi ,\frac{3}{2}%
      \right)$}

    \item {La región acotada por $y=1-\left\vert x\right\vert $ y el eje $x,$ \
      $\rho \left( x\right) =x^{2}$}{$m=\frac{1}{6},$ Centro de masa $\left( 0,\frac{1}{5}\right)$}

    \item {La región acotada por $y=e^{-x},y=e^{x}$ y $x=\ln 2,$ \ $\rho
      \left( x\right) =e^{x}$}{$m=\frac{3}{2}-\ln 2,$ \ Centro de masa $\left( \frac{2\ln 2-\frac{1}{2}\ln ^{2}2-\frac{3}{4}}{%
      \frac{3}{2}-\ln 2},\frac{\frac{11}{12}}{\frac{3}{2}-\ln 2}\right)$}

    \item { La región acotada por $y=\ln x,y=0$ y $x=e,$ \ $\rho \left(
      x\right) =4$}{$m=4,$ Centro de masa $\left( \frac{e^{2}+1}{4},\frac{2e-4}{4}\right) $}

    \item {La región limitada por $y=\sec ^{2}\left( x\right) ,\tan ^{2}\left(
        x\right) $ y las rectas $x=-\frac{\pi }{4},x=\frac{\pi }{4}$ con función
      de densidad $\rho \left( x\right) =2$}{$m=\pi ,$ $\ $Centro de masa $\left( 0,\frac{8-\pi}{\pi} \right) $}

    \item {La región limitada por $y=x^{2}-x,y=x$, con función de densidad $%
      \rho \left( x\right) =\sqrt{x+1}$}{$\ m=\frac{24}{35}\sqrt{3}+\frac{24}{35},$ Centro de masa $\left(1.051,0.65029\right) $}

    \item {La región encerrada por la semicircunferencia $x^{2}+y^{2}=4,y\geq 0$,
      $\ $con función de densidad $\rho \left( x\right) =x^{2}$}{$m=32\pi ,$ \ Centro de masa $\left( 0,\frac{64}{15 \pi} \right)$}

    \item {La región limitada por $y=e^{x},y=e^{-x}$ y $x=\ln 3$, con funció%
      n de densidad $\rho \left( x\right) =x$}{$m=\frac{10}{3}\ln 3-\frac{8}{3},$ Centro de masa $\left(
      0.83442,1.3977\right) $}

    \item {La región limitada por $y=e^{x},y=x^{2}e^{x}$, con función de
      densidad $\rho \left( x\right) =x^{2}$}{$m=0.32651$, \ Centro de masa $\left( 0.38348,1.2437\right) $}

    \item {La región limitada por $y=2\ln x,y=x\ln x$, con función de
      densidad $\rho \left( x\right) =x^{2}$}{$m=0.30614$, Centro de masa $\left( 1.5484,0.77321\right) $}

    \item {Un tanque en forma de prisma rectangular, con un largo de $6$ pies, ancho de $4$ pies y alto de $10$ pies, está lleno de agua. Encuentre el trabajo
      realizado para bombear toda el agua hasta la parte superior del tanque.}{$74880$ lb-pie}

    \item {Un tanque en forma de pirámide y con base cuadrada, de lado $5$ pies y
        altura de $12$ pies, con el vértice hacia abajo, está lleno de agua. \
        Encuentre el trabajo realizado para bombear toda el agua hasta $2$ pies arriba
      de la parte superior del tanque.}{$18720$ lb-pie}

    \item { Un tanque en forma de cono circular recto, con radio de la base de $6$ pies y altura de $14$ pies y con el vértice hacia abajo está lleno de agua
        hasta una profundidad de $8$ pies. \ Encuentre el trabajo realizado para
        bombear el agua hasta una altura de $3$ pies por encima de la parte superior del
      tanque.}{$\frac{3833856}{245}\pi $ lb-pie}

    \item {Un tanque hemisférico de radio de $30$ pies está lleno de agua
        hasta una profundidad de $20$ pies. \ Encuentre el trabajo realizado para
      bombear toda el agua hasta la parte superior del tanque.}{$9984000\pi $ lb-pie}

    \item { Un tanque hemisférico de radio de 10 pies está lleno de agua
        hasta una profundidad de 5 pies. \ Encuentre el trabajo realizado para
      bombear toda el agua hasta 1 pie arriba de la parte superior del tanque.}{$100750\pi \ \ lb-pie$}

    \item { Un tanque parabólico, generado por el giro de la recta $y=\frac{1}{2}%
        x^{2},$ $0\leq x\leq 2$ ($y$ en pies ) alrededor del eje $y$, está lleno de
        agua. \ Halle el trabajo realizado para bombear el agua hasta 1 pie por
      encima del tanque.}{$1306.9 $ lb-pie}

    \item { Un tanque, generado por el giro de la recta $y=x^{3},$ $0\leq x\leq 2$ ($y$ en pies ) alrededor del eje $y$, está lleno de agua. \ Encuentre el trabajo
      realizado para bombear el agua hasta $2$ pies por encima del tanque.}{$\frac{29952}{5}\pi $ lb-pie}

    \item {Un tanque, generado por el giro de la recta $y=x^{4},$ $0\leq x\leq 2$ ($%
        y$ en pies ) alrededor del eje $y$, está lleno de agua. \ Encuentre el trabajo
      realizado para bombear el agua hasta $1$ pie por encima del tanque.}{$\frac{492544}{25}\pi $ lb-pie}

    \item {Una placa rectangular plana de $20$ pies por $26$ pies se sumerge
        horizontalmente en agua a una profundidad de $15$ pies. Determine la presi%
      ón y la fuerza ejercidas sobre la placa por el agua arriba de esta.}{$486720$ lb}

    \item {Una placa en forma de disco de radio de $8$ pies se sumerge
        horizontalmente en agua a una profundidad de $12$ pies. Determine la presi%
      ón y la fuerza ejercidas sobre la placa por el agua arriba de esta.}{$150555.525$ lb}

    \item {La placa formada por un rectángulo de base $6$ pies y altura $10$ pies se sumerge verticalmente en agua hasta que la base coincida con el nivel del
      agua. Encuentre la fuerza ejercida por el agua sobre un lado de la placa.}{$18720$ lb}

    \item {La placa formada por un triángulo isóceles de base $8$ pies y
        altura $12$ pies se sumerge verticalmente en agua, hasta que la base coincida
        con el nivel del agua. Encuentre la fuerza ejercida por el agua sobre un
      lado de la placa.}{$11980.8$ lb.}

    \item {La placa formada por el sector parabólico limitado por $y=x^{2}$ y la
        recta $y=9$ ($y$ en pies) se sumerge verticalmente en agua hasta que la
        recta $y=11$ coincida con el nivel del agua. Encuentre la fuerza ejercida por
      el agua sobre un lado de la placa.}{$12580$ lb}

    \item {Una placa en forma de mitad de un disco de radio $3$ pies se sumerge
        verticalmente en agua hasta que el diámetro coincida con el nivel del agua.
      Encuentre la fuerza ejercida por el agua sobre un lado de la placa.}{$1523.3$ lb}

    \item {Una placa en forma de triángulo equilátero de lado $5$ pies y con la base horizontal se sumerge verticalmente en agua hasta que la base coincida con
        el nivel del agua. Encuentre la fuerza ejercida por el agua sobre un lado de
      la placa.}{$975$ lb}

    \item {La placa formada por la región limitada por $y=x^{2}-1$ y $y=7-x^{2}$
        ($y$ en pies) se sumerge verticalmente en agua hasta que el vértice de $%
        y=7-x^{2}$ sea tangente al nivel del agua. Encuentre la fuerza ejercida por
      el agua sobre un lado de la placa.}{$5324.8$ lb}

    \item {La placa formada por la región limitada por $y=4-x^{2}$ y $y=4-4x$ ($%
        y$ en pies) se sumerge verticalmente en agua, de tal forma que el nivel del
        agua corresponde con la recta $y=8$. Encuentre la fuerza ejercida por el
      agua sobre un lado de la placa.}{$6922.2$ lb.}

    \item {\label{masafin}La placa formada por la región limitada por $y=\sen x$ y $y=0,0\leq
        x\leq \pi $ ($y$ en pies) se sumerge verticalmente en agua, de tal forma
        que el nivel del agua coincida con la recta $y=2$. Encuentre la fuerza
      ejercida por el agua sobre un lado de la placa.}{$200.59$ lb.}
  \end{enumerate}
}
